%----------------------------------------------------------------------------------------------------%
%----------------------------------------------------------------------------------------------------%
%%%%%%%%%%%                                              HEADER                                                %%%%%%%%%%%
%----------------------------------------------------------------------------------------------------%
%----------------------------------------------------------------------------------------------------%
%
% Universale header.tex – gilt für Arbeitsblätter, Leistungskontrollen, Handreichungen, Nachschlagewerke.
% Neue Pakete und TikZ-Bibliotheken werden HIER ergänzt (nicht im Dokument) und unten protokolliert.
% Quelle: https://github.com/mq90/latex-skills-schule  (CC BY 4.0, Martin Quandt)
%
% Änderungsprotokoll
% 2026-11-09  zusammengeführt aus allen Handreichungs-Headern: TikZ-Bibliotheken fit, shapes.multipart,
%             arrows.meta, shapes.callouts, circuits.logic.US, circuits.logic.IEC; Pakete array, booktabs,
%             longtable; xskak, chessboard, sudoku, tikzpeople, tipa (safe); MusiXTeX als auskommentierte
%             Option (definiert \cup als Achtelnote um – bricht $A \cup B$)
% 2026-11-11  Naturwissenschaften: circuitikz [european] (Schaltpläne), mhchem [version=4] (\ce),
%             chemfig (Strukturformeln), pgf-PeriodicTable (\pgfPT), genealogytree [all] (Stammbäume).
%             NICHT siunitx: kollidiert mit units (\unit[10]{cm}); siunitx nur in eigenständigen Dokumenten
% 2026-11-11  TikZ-Bibliothek decorations.pathmorphing (Federn, Wellen: decoration=coil/snake)
% 2026-11-23  twemojis: farbige Emojis mit pdflatex, z.B. \twemoji[height=1cm]{man standing} (Skill latex-sachskizzen)
%






%----------------------------------------------------------------------------------------------------%
%%%%%%%%%%%%%                            Gestaltung des Dokuments                                %%%%%%%%%%
%----------------------------------------------------------------------------------------------------%



%Definition der Dokumentenklasse
\documentclass[12pt]{article} %book, article, report, standalone



%Pakete für eine saubere Kodierung
\usepackage[utf8]{inputenc} %UTF-8 als Zeichenkodierung  ... ABER nicht immer notwendig, je nach Compiler
\usepackage[T1]{fontenc} %Verwendet westeuropäische Schriftkodierung
\usepackage[ngerman]{babel} %Ermöglicht deutsche Übersetzung und Silbentrennung
\usepackage{lmodern} %Verwendet moderne Schriftarten



%Pakete für die Gestaltung des Dokuments
\usepackage{geometry} %Ermöglicht die Ausrichtung des Blattes mit \geometry{...}
\geometry{a4paper, top=15mm, left=15mm, right=20mm, bottom=15mm, headsep=10mm, footskip=10mm} %gewählte Ausrichtung des Blattes
\pagestyle{empty} %Art der Seitennummerierung, z.B.: empty , plain
\usepackage{fancyhdr} %Ermöglicht die Gestaltung einer Kopfzeile für LK's oder KA's 
\usepackage{lscape} %Ermöglicht eventuelles Querformat des Blattes mit \begin{landscape} ... \end{landscape}
\usepackage{pdfpages} %Ermöglicht einbinden PDF-Dokumente bzw. einzelene Seiten mit \includepdf[pages={1,3,5}]{Dateiname} bzw. pages=-
\usepackage[german]{layout} %Ermöglicht es mit den Befehl \layout sich die Default-Layout beschreiben zu lassen (Bild vom Layout)
%----------------------------------------------------------------------------------------------------%








%----------------------------------------------------------------------------------------------------%
%%%%%%%%%%%%                                        Packages                                                  %%%%%%%%%%%
%----------------------------------------------------------------------------------------------------%



%Pakete für die Strukturierung des Dokuments
\usepackage[shortlabels]{enumitem} %Ermöglicht die Nummerierung selbst festzulegen mit \begin{enumerate}[Art der Nummerierung] ...   \end{enumerate}
\parindent=0px %Setzt die Absatzeinrückung auf 0 Pixel und ernfernt damit unnötige Leerzeichen vor dem ersten Wort eines neuen Abschnitts
\usepackage{setspace} %Ermöglicht verschiedene Funktion zur Abstandsregelung, z.B.: \hspace{2.5cm} ; \onehalfspacing
\onehalfspacing %Setzt den Zeilenabstand auf 1,5
\usepackage[style=german]{csquotes}%Ermöglicht den deutschen Zitierstil darstellen zu lassen mt \enquote{Zitat}
\usepackage[breaklinks=true,%Zeilenumbruch der Verlinkung ist bündig
			colorlinks=true,%Links können farbig sein
			linkcolor=blue,%Linkfarbe wird auf blau gesetzt
			urlcolor=darkgray,%Farb der URL wird auf dunkelgrau gesetzt
			plainpages=false,%
			hyperfootnotes=false,%
			pdfpagelabels=false%
		]{hyperref} %Ermöglicht die Verlinkung innerhalb des Dokuments 
\usepackage{multicol} %Ermöglicht gleichgroße Aufteilung des Dokuments in n Spalten mit \begin{multicols}{n} ... \end{multicols}
\usepackage{xcolor,colortbl} %Ermöglicht Zeilen, Spalten und Zellen einer Tabelle zu färben
\usepackage{caption} %Ermöglicht die Bezeichnung und Strukturierung von Abbildungen oder Quellcode
\usepackage{fancybox} % Ermöglicht die Umrandung einzelner Abschnitte
\usepackage[most]{tcolorbox} %Ermöglicht die Nutzung von Colorboxen für Merksätze, Lösungen, Definitionen, usw mit \newtcolorbox
\tcbuselibrary{breakable} %Aufteilen der Colorboxen, welche über mehrere Seiten gehen
\usepackage{ifthen} % Ermöglicht \ifthenelse Befehle
\allowdisplaybreaks %Erlaubt Seitenumbrüche innerhalb von align-Umgebung



%Nützliche Pakete für Symbole und mathematischen Formelsatz
\usepackage{amsmath, amssymb, mathrsfs, amsfonts, amsthm, dsfont} %sämtliche Pakete für mathematische Zeichen, Symbole und Formelsatz
\usepackage{units} %Ermöglicht die Angabe von Einheiten mit \unit[Wert]{Einheit}
\usepackage{gensymb} %Ermöglicht Nutzung von \degree für das macro des Grad-Symbols °
\usepackage[official]{eurosym} %Ermöglicht die Nutzung von \euro{}
\usepackage{pifont} %Ermöglicht die Nutzung verschiedener Fonts, z.B.: \ding{42} , \ding{52} , \ding{56}
\usepackage{tikzsymbols} %Ermöglicht die Nutzung von u.a. Smileys und anderen Symbolen, z.B.: \Winkey , \Smiley[1.5]



%Pakete zur Darstellung von Grafiken
\usepackage{tikz} %\begin{tikzpicture} .... \end{tikzpicture}
\usepackage{pgf, pgfplots} %Ermöglicht 2D und 3D Plof-Funktion
\pgfplotsset{compat=1.17}
\usetikzlibrary{arrows,%Pfeile
			calc,%Rechnen in Koordinaten, z.B. mit ($(1,2)+5*(0,1)$) auf Punkt (1,7) kommen
			angles, plotmarks, quotes, babel, intersections, shapes.geometric, decorations.pathreplacing, decorations.pathmorphing, automata, positioning, shadows, patterns, patterns.meta,%
			fit,%Rahmen um mehrere Knoten (UML, Automaten)
			shapes.multipart,%geteilte Knoten (UML-Klassen)
			arrows.meta,%moderne Pfeilspitzen, z.B. -{Triangle[open]}
			shapes.callouts,%Sprechblasen
			circuits.logic.US, circuits.logic.IEC%Logikgatter (digitale Schaltungen)
			} %nützliche Bibliotheken für TikZ
\usepackage{animate} %animierte Darstellungen
\usepackage{graphicx} %\includegraphics[width=..px,height=..px]{Dateiname}
\usepackage{multirow} %Zellen über mehrere Zeilen mit \multirow{n}{*}{Text}
\usepackage{array} %erweiterte Spaltentypen, z.B. >{\centering\arraybackslash}p{2cm}
\usepackage{booktabs} %schöne Tabellenlinien \toprule, \midrule, \bottomrule
\usepackage{longtable} %Tabellen über mehrere Seiten



%gelegentlich hilfreiche Zusatzpakete
\usepackage{hologo} %\hologo{LaTeX}
\usepackage{suetterl} %Sütterlin als Schriftart mit \mathfrak{Text}
\usepackage{qrcode} %Einbindung eines QR-Codes 
\usepackage{url} %Einbindung eines URL-Links
\usepackage{ccicons} %Creative Commons Symbole
\usepackage{listings} %\begin{lstlisting} ... \end{lstlisting}
\usepackage{struktex} %Struktogrammdarstellung
\usepackage{ulem} %Unterstreichung mit \uline{} ggf. mit \hfill im Inneren für Unterstrich bis zum Ende der Zeile
\usepackage{xskak} %Schach: Partien und Züge mit \newchessgame, \mainline{1. e4 e5} (lädt skak)
\usepackage{chessboard} %Schachbretter mit \chessboard[setfen=...]
\usepackage{sudoku} %Sudoku-Gitter mit \begin{sudoku} |5|3| | ... |. \end{sudoku}
\usepackage{tikzpeople} %Personen-Figuren als TikZ-Knoten, z.B. \node[alice, minimum size=2cm]{};
\usepackage{twemojis} %farbige Emojis (Twemoji, CC BY 4.0) als Vektorgrafik: \twemoji[height=1cm]{man standing}, im Text \texttwemoji{light bulb}
\usepackage[safe]{tipa} %Lautschrift mit \textipa{...}; [safe], weil tipa sonst \; \: \! \* \s umdefiniert (bricht z.B. \(\;\)) – Langformen siehe latex-lautschrift
\usepackage[european]{circuitikz} %Schaltpläne mit deutschen Symbolen: \begin{circuitikz} \draw (0,0) to[lamp] (2,0); \end{circuitikz}
\usepackage[version=4]{mhchem} %Reaktionsgleichungen mit \ce{2 H2 + O2 -> 2 H2O}
\usepackage{chemfig} %Strukturformeln mit \chemfig{H-C(-[2]H)(-[6]H)-H}, freie Elektronenpaare mit \charge
\usepackage{pgf-PeriodicTable} %Periodensystem mit \pgfPT[...], deutsche Namen mit \pgfPTsetLanguage{de}
\usepackage[all]{genealogytree} %Stammbäume mit \begin{genealogypicture} ... \end{genealogypicture}
%\usepackage{musixtex} %Notensatz – NUR für Musik-Dokumente einkommentieren: definiert u.a. \cup als Note um (bricht $A \cup B$)
%\usepackage[list=true, font=large, labelfont=bf, labelformat=brace, position=top]{subcaption}



%abgeänderte Befehle
\renewcommand{\arraystretch}{1.25} %
%\renewcommand{\familydefault}{\sfdefault} %serifenlose Schriftart
\renewcommand{\lstlistingname}{Quellcode} %Umbenennung
\renewcommand{\lstlistlistingname}{Quellcodeverzeichnis} %Umbenennung
%----------------------------------------------------------------------------------------------------%
%----------------------------------------------------------------------------------------------------%
%----------------------------------------------------------------------------------------------------%
%%%% %%%%%%                                     NEWCOMMANDS                                         %%%%%%%%%%%%
%----------------------------------------------------------------------------------------------------%
%----------------------------------------------------------------------------------------------------%
%
% Universale newcommands.tex – alle eigenen Befehle, Farben, Boxen, Kopfzeilen, TikZ- und Listing-Stile.
% Neue Befehle werden HIER ergänzt (nicht im Dokument) und unten protokolliert.
% Quelle: https://github.com/mq90/latex-skills-schule  (CC BY 4.0, Martin Quandt)
%
% Änderungsprotokoll
% 2026-11-09  zusammengeführt aus allen Handreichungen: \TikZ, Box »ausgabe«, Stile für
%             Programmablaufpläne (input, operation, condition, output) und UML (umlclass, inherit,
%             compose, aggregate, associate, depend); Listing-Stil »latex« wieder wie in den
%             Handreichungen ohne Zeilennummern, nummerierte Variante als »latexnummeriert«;
%             Umlaute und · € ° → ≤ ≥ ≠ – — „ “ … × ² ³ π in allen Listings erlaubt (\lstset literate); \lehrerblase, \schuelerblase; PAP-Stile mit align=center (Zeilenumbruch mit \\);
%             \ipa{...} für Lautschrift; \koordinatensystem[Schritt]{xmin}{xmax}{ymin}{ymax};
%             TikZ-Stile punkt, punktoffen, baumdiagramm, wkeit
% 2026-11-09  TikZ-Stil baumdiagrammvertikal (Baumdiagramm von oben nach unten)
% 2026-11-09  Formularfelder für ausfüllbare PDFs: \feld, \zeichen, \kreuz, \antwort, einheitliches
%             Aussehen über \DefaultOptionsofText/CheckBox/Radio/PopdownBox (nur wirksam in einer Form-Umgebung)
% 2026-11-11  Akzentfarben akzentblau, akzentrot, akzentgruen, akzentorange, akzentlila (Karten,
%             Naturwissenschaften, Schulorganisation); Box »hinweis« mit frei wählbarem Titel
% 2026-11-11  Umgebung »passend«: verkleinert den Inhalt auf \linewidth, aber nur wenn er breiter ist
% 2026-11-11  \oz{+II}{Cu}: Oxidationszahl über ein Elementsymbol (Chemie, auch innerhalb \ce{... $\oz{0}{Zn}$ ...})
% 2026-11-22  Farbige Schachbretter: \schachbrettfarben{braun|blau|gruen|grau|klassisch} (Standard: braun),
%             Farben brett...hell/brett...dunkel und brettmarke (braucht LSBC5 aus der header.tex)
%








%----------------------------------------------------------------------------------------------------%
%%%%%%%%%%%                                   Mathematische Kürzel                                       %%%%%%%%%%
%----------------------------------------------------------------------------------------------------%


%vereinfachte Schreibweise für den Displaystyle
\newcommand{\dis}[1]{\displaystyle{#1}}


%Zahlenbereiche
\newcommand{\N}{\mathbb{N}} %natürliche Zahlen
\newcommand{\Z}{\mathbb{Z}} %ganze Zahlen
\newcommand{\Q}{\mathbb{Q}} %rationale Zahlen
\newcommand{\R}{\mathbb{R}} %reelle Zahlen
\newcommand{\C}{\mathbb{C}} %komplexe Zahlen
\newcommand{\cP}{\mathcal{P}} %Potenzmenge


%Abkürzungen für Analysis
\newcommand{\eps}{\varepsilon} %Kürzel für Epsilon
\newcommand{\limn}{\underset{n\longrightarrow \infty}{\lim}} %Grenzwert für n gegen Unendlich
\newcommand{\limh}{\underset{h\longrightarrow 0}{\lim}} %Grenzwert für h gegen Null
\newcommand{\limxto}[1]{\underset{x\longrightarrow #1}{\lim}} %beliebiger funktionaler Grenzwert x gegen ....
\newcommand{\dx}{\text{\;d} x} %dx Operator für Integrale
\renewcommand{\d}[1]{\text{\;d}#1} %Operator für beliebige Integrationsvariable
\newcommand{\norm}[1]{\left\lVert#1\right\rVert} %\norm{...} für die Normdarstellung ||...||


%Abkürzungen für Analytische Geometrie
\newcommand{\vektor}[3]{\begin{tiny}\begin{pmatrix} \text{\begin{scriptsize}$#1$\end{scriptsize}}\\ \text{\begin{scriptsize}$#2$\end{scriptsize}}\\ \text{\begin{scriptsize}$#3$\end{scriptsize}}\end{pmatrix}\end{tiny}} %dreidimensionaler Vektor
\newcommand{\punktvektor}[2]{\begin{tiny}\begin{pmatrix} \text{\begin{scriptsize}$#1$\end{scriptsize}}\\ \text{\begin{scriptsize}$#2$\end{scriptsize}}\end{pmatrix}\end{tiny}} %zweidimensionaler Vektor


%Abkürzungen für Stochastik
\newcommand{\E}{\mathbb{E}} %Erwartungswert
\newcommand{\V}{\mathbb{V}} %Varianz


%Umgebung für gegeben und gesucht
\newcommand{\gegges}[2]{%
	\begin{multicols}{2}
		\begin{enumerate}
			\item[geg:] #1
			\item[ges:] #2
		\end{enumerate}
	\end{multicols}
}
%----------------------------------------------------------------------------------------------------%









%----------------------------------------------------------------------------------------------------%
%%%%%%%%%%%%                             Verschiedene  hilfreiche Kürzel                                %%%%%%%%%
%----------------------------------------------------------------------------------------------------%


%römische Darstellung einer Zahl \uproman{2021} ergibt MMXXI
\newcommand{\uproman}[1]{\uppercase\expandafter{\romannumeral#1}}


%Automatisch erstellte Lücke wie die Länge des Textes mit \luecke{<Text>}
\newcommand{\luecke}[1]{\hspace{0,1cm}\underline{\hspace{0,25cm}\hphantom{#1}\hspace{0,25cm}}\hspace{0,1cm}}


%fett-kursive Schrift für neu eingeführte Begriffe
\newcommand{\fk}[1]{\textbf{\textit{#1}}}


%Schwierigkeitsstufen für Aufgaben mit \leicht , \mittel , \schwer
\newcommand{\leicht}{\smash{\raisebox{-1pt}{{\color{black}\rule{4pt}{8pt}}\,{\color{lightgray}\rule{4pt}{8pt}}\,{\color{lightgray}\rule{4pt}{8pt}}}}\,}
\newcommand{\mittel}{\smash{\raisebox{-1pt}{{\color{black}\rule{4pt}{8pt}}\,{\color{black}\rule{4pt}{8pt}}\,{\color{lightgray}\rule{4pt}{8pt}}}}\,}
\newcommand{\schwer}{\smash{\raisebox{-1pt}{{\color{black}\rule{4pt}{8pt}}\,{\color{black}\rule{4pt}{8pt}}\,{\color{black}\rule{4pt}{8pt}}}}\,}
%----------------------------------------------------------------------------------------------------%










%----------------------------------------------------------------------------------------------------%
%%%%%%%                                            Verschiedene Farben                                                      %%%%%%%
%----------------------------------------------------------------------------------------------------%


%Formeln in herkömmlichen Farben schreiben
\newcommand{\frot}[1]{\text{\textcolor{red}{$#1$}}} % rote Formel für Lösungen
\newcommand{\fgruen}[1]{\text{\textcolor{green}{$#1$}}} % grüne Formel für Lösungen
\newcommand{\fblau}[1]{\text{\textcolor{blue}{$#1$}}} % blaue Formel für Lösungen
\newcommand{\fmagenta}[1]{\text{\textcolor{magenta}{$#1$}}} % magenta-farbige Formel für Lösungen
\newcommand{\forange}[1]{\text{\textcolor{orange}{$#1$}}} % orangene Formel für Lösungen
\newcommand{\fcyan}[1]{\text{\textcolor{cyan}{$#1$}}} % cyan-farbige Formel für Lösungen
\newcommand{\folive}[1]{\text{\textcolor{olive}{$#1$}}} % olive-farbige Formel für Lösungen


%Custom made colors
\definecolor{plum}{rgb}{0.36078, 0.20784, 0.4}
\definecolor{chameleon}{rgb}{0.30588, 0.60392, 0.023529}
\definecolor{cornflower}{rgb}{0.12549, 0.29020, 0.52941}
\definecolor{scarlet}{rgb}{0.8, 0, 0}
\definecolor{brick}{rgb}{0.64314, 0, 0}
\definecolor{sunrise}{rgb}{0.80784, 0.36078, 0}
\definecolor{lightblue}{rgb}{0.15,0.35,0.75}
\definecolor{carolina}{rgb}{153, 186, 221}

%Akzentfarben für Grafiken, Karten und Diagramme (kräftig, gut unterscheidbar; aufgehellt z.B. akzentblau!18)
\definecolor{akzentblau}{HTML}{2563EB}
\definecolor{akzentrot}{HTML}{DC2626}
\definecolor{akzentgruen}{HTML}{16A34A}
\definecolor{akzentorange}{HTML}{EA580C}
\definecolor{akzentlila}{HTML}{7C3AED}
%----------------------------------------------------------------------------------------------------%









%----------------------------------------------------------------------------------------------------%
%%%%%%%%                            Kopfzeilen für Leistungsüberprüfungen                              %%%%%%%%%
%----------------------------------------------------------------------------------------------------%


%Definition Kopfzeile für Leistungskontrolle mit \lk{Klasse}{Thema}
\newcommand{\lk}[2]{%
	\pagestyle{fancy}
	\fancyhead[L]{%
	\begin{tabular}[b]{l}
		\\Name:\\Punkte:\\Note:
	\end{tabular} 
	}
	\fancyhead[R]{%
	\begin{tabular}[b]{r}
		\\Leistungskontrolle\\Klasse #1\\ \today
	\end{tabular} 
	}
	\fancyhead[C]{%
	\textbf{\large{#2}}\\
	}
	\fancyfoot[C]{}
	\renewcommand{\headrulewidth}{1pt}
	\fancyheadoffset{1cm}
}


%Definition Kopfzeile für Klassenarbeit mit \ka{Klasse}{Thema}
\newcommand{\ka}[2]{%
	\pagestyle{fancy}
	\fancyhead[L]{%
	\begin{tabular}[b]{l}
		\\Name:\\Punkte:\\Note:
	\end{tabular} 
	}
	\fancyhead[R]{%
	\begin{tabular}[b]{r}
		\\Klassenarbeit\\Klasse #1\\ \today
	\end{tabular} 
	}
	\fancyhead[C]{%
	\textbf{\large{#2}}\\
	}
	\fancyfoot[C]{}
	\renewcommand{\headrulewidth}{1pt}
	\fancyheadoffset{1cm}
}


%Definition Kopfzeile für Klausur mit \kl{Klasse}{Thema}
\newcommand{\kl}[2]{%
	\pagestyle{fancy}
	\fancyhead[L]{%
	\begin{tabular}[b]{l}
		\\Name:\\Punkte:\\Note:
	\end{tabular} 
	}
	\fancyhead[R]{%
	\begin{tabular}[b]{r}
		\\Klausur\\Klasse #1\\ \today
	\end{tabular} 
	}
	\fancyhead[C]{%
	\textbf{\large{#2}}\\
	}
	\fancyfoot[C]{}
	\renewcommand{\headrulewidth}{1pt}
	\fancyheadoffset{1cm}
}


%Entfernt den Header für Leistungskontrollen, Klassenarbeiten oder Klausuren
\newcommand{\clearfancyheader}{%
	\fancyhf{}
	\renewcommand{\headrulewidth}{0pt}
	\renewcommand{\headheight}{0pt}
}


%Definition Kopfzeile für Tägliche Übungen mit \tu{Klasse}{Thema}
\newcommand{\tu}[2]{%
	\pagestyle{fancy}
	\fancyhead[L]{%
	\begin{tabular}[b]{l}
		\\Name:\\Punkte:\\Note:
	\end{tabular} 
	}
	\fancyhead[R]{%
	\begin{tabular}[b]{r}
		\\Tägliche Übung\\Klasse #1\\ \today
	\end{tabular} 
	}
	\fancyhead[C]{%
	\textbf{\large{#2}}\\
	}
	\fancyfoot[C]{}
	\renewcommand{\headrulewidth}{1pt}
	\fancyheadoffset{1cm}
}
%----------------------------------------------------------------------------------------------------%








%----------------------------------------------------------------------------------------------------%
%%%%%%%%%                                     Colorboxes für Hervorhebungen                              %%%%%%%%%
%----------------------------------------------------------------------------------------------------%


 %Colorboxes für Lösungen mit \begin{solution}<Text>\end{solution}
\newtcolorbox{solution}{%
        	lower separated=false,%
	colback=gray!5!white,%
	colframe=gray!75!black,%
        	fonttitle=\bfseries,%
	title=Lösung:,%
      	enhanced,%
	breakable,%
        	attach boxed title to top left={xshift=0.5cm,yshift=-2mm}%
}


%Colorboxes für das chronologische Vorgehen 
\newtcolorbox{todo}{%
        	lower separated=false,%
	colback=yellow!5!white,%
	colframe=orange!75!black,%
        	enhanced,%
	breakable%
}


%Colorboxes für Merksätze
\newtcolorbox{merke}{%
        	lower separated=false,%
	colback=red!5!white,%
	colframe=red!75!black,%
        	fonttitle=\bfseries,%
	title=Merke:,%
      	enhanced,%
	breakable,%
        	attach boxed title to top left={xshift=0.5cm,yshift=-2mm}%
}


%Hinweis-Box mit frei wählbarem Titel: \begin{hinweis}[Lizenz] ... \end{hinweis} (Standardtitel: Hinweis)
\newtcolorbox{hinweis}[1][Hinweis]{%
	enhanced,%
	breakable,%
	sharp corners,%
	colback=akzentblau!5,%
	colframe=akzentblau,%
	boxrule=0pt,%
	leftrule=3pt,%
	fonttitle=\bfseries,%
	coltitle=akzentblau!85!black,%
	attach title to upper={\ },%
	title={#1:}%
}


%Colorbox2 für Merksätze 
\newtcolorbox{merke2}{%
        	lower separated=false,%
        	colback=red!40,%
       	title=Merke,%
       	colframe=black,%
        	fonttitle=\bfseries,%
        	colbacktitle=red,%
        	coltitle=black,%
        	enhanced,%
        	breakable,%
        	attach boxed title to top left={xshift=0.5cm,yshift=-2mm},%
        	title=Merke:%
}
%----------------------------------------------------------------------------------------------------%










%----------------------------------------------------------------------------------------------------%
%%%%%%%%%%%%%                                         MC-Fragen                                    %%%%%%%%%%%%%
%----------------------------------------------------------------------------------------------------%

%Definition des Kästchens zum ankreuzen
\newcommand{\casket}{\LARGE{$\square$}}


%Umgebung für horizontale MC-Fragen
\newenvironment{mcquestionh}[1]{%
	#1%
	\begin{center}\begin{tabular}{>{\centering\arraybackslash}p{2.5cm}  >{\centering\arraybackslash}p{2.5cm} >{\centering\arraybackslash}p{2.5cm} >{\centering\arraybackslash}p{2.5cm} >{\centering\arraybackslash}p{2.5cm}}
	\casket & \casket & \casket & \casket & \casket\medskip\\
	}
	%hier wird die fehlende Tabellenzeile angegeben. Aufruf erfolgt durch:
	%\begin{mcquestionh}{<Frage>}
	%   	\answerh{<Antwort>} & \answerh{<Antwort>} & \answerh{<Antwort>} & \answerh{<Antwort>} & \answerh{<Antwort>}
	%\end{mcquestionh}
	{
	\\
	\end{tabular}\end{center}
}


%Antwort für horizontale MC-Fragen mit \answerh{<Antwort>}
\newcommand{\answerh}[1]{%
	\parbox{2.5cm}{\centering{#1}}
}


%Umgebung für vertikale MC-Fragen
\newenvironment{mcquestionv}[1]{%
    	#1%
    	\begin{enumerate}[\raisebox{-3pt}{\casket}]
    	}
    	%hier wird das fehlende \item angegeben werden. Aufruf erfolgt durch:
    	%\begin{mcquestionv}{<Frage>}
    	%   	\answerv{<Antwort>}
    	%  	\answerv{<Antwort>}
    	%   	\answerv{<Antwort>}
    	%   	\answerv{<Antwort>}
    	%   	\answerv{<Antwort>}
    	%\end{mcquestionv}
    	{
    	\end{enumerate}
}


%Antwort für vertikale MC-Frage mit \answerv{<Antwort>}
\newcommand{\answerv}[1]{%
	\item \hspace{0.1cm} #1
}
%----------------------------------------------------------------------------------------------------%









%----------------------------------------------------------------------------------------------------%
%%%%%%%%%%%                                          Tikz Settings                                          %%%%%%%%%%%%
%----------------------------------------------------------------------------------------------------%


%Koordinatensystem im Standardstil der Mathe-Arbeitsblätter (Gitter gray!30, Achsen -latex, Beschriftung \tiny)
%\koordinatensystem[Schritt]{xmin}{xmax}{ymin}{ymax};
%             TikZ-Stile punkt, punktoffen, baumdiagramm, wkeit
% 2026-11-09  TikZ-Stil baumdiagrammvertikal (Baumdiagramm von oben nach unten) – ganze Zahlen, Bereich muss 0 enthalten,
%Striche und Zahlen bei allen Vielfachen von Schritt (Standard 2). Innerhalb von tikzpicture[scale=0.42] verwenden.
\newcommand{\koordinatensystem}[5][2]{%
	\foreach \kosyx in {#2,...,#3} \draw[gray!30, thin] (\kosyx,#4) -- (\kosyx,#5);
	\foreach \kosyy in {#4,...,#5} \draw[gray!30, thin] (#2,\kosyy) -- (#3,\kosyy);
	\draw[-latex, thick] ({#2-0.3},0) -- ({#3+0.5},0) node[right] {$x$};
	\draw[-latex, thick] (0,{#4-0.3}) -- (0,{#5+0.5}) node[above] {$y$};
	\foreach \kosyx in {#2,...,#3} {%
		\pgfmathtruncatemacro{\kosyrest}{mod(abs(\kosyx),#1)}%
		\ifnum\kosyrest=0 \ifnum\kosyx=0 \else
			\draw (\kosyx,0.1) -- (\kosyx,-0.1) node[below, font=\tiny] {$\kosyx$};
		\fi\fi}
	\foreach \kosyy in {#4,...,#5} {%
		\pgfmathtruncatemacro{\kosyrest}{mod(abs(\kosyy),#1)}%
		\ifnum\kosyrest=0 \ifnum\kosyy=0 \else
			\draw (0.1,\kosyy) -- (-0.1,\kosyy) node[left, font=\tiny] {$\kosyy$};
		\fi\fi}
}


%optionale Eigenschaft für ein  Tikzpicture, damit dieses bündig am oberen Rand einer enumerate-Auflistung steht (benötigt tikzlibrary calc!!!)
\tikzset{
  topline/.style={baseline={([yshift=#1]current bounding box.north)}},
  topline/.default=-\ht\strutbox
}


%Punkte in Graphen und Figuren: \node[punkt, label=above:{$P$}] at (2,3) {}; (offen: punktoffen)
%Als Knoten, damit die Größe nicht mit scale=0.42 mitschrumpft (circle (3pt) würde mitskaliert)
\tikzset{
  punkt/.style={circle, fill, inner sep=0pt, minimum size=4pt},
  punktoffen/.style={circle, draw, fill=white, inner sep=0pt, minimum size=4pt}
}


%Baumdiagramme (Stochastik) von links nach rechts: \begin{tikzpicture}[baumdiagramm] \node[punkt]{} child{...}
%Wahrscheinlichkeiten an den Ästen: edge from parent node[wkeit, above] {$0{,}6$} (oberer Ast above, unterer below)
\tikzset{
  baumdiagramm/.style={grow'=right, sloped, level distance=2.6cm,
    level 1/.style={sibling distance=3cm}, level 2/.style={sibling distance=1.5cm},
    level 3/.style={sibling distance=0.75cm}, edge from parent/.style={draw, thin}},
  wkeit/.style={font=\scriptsize, inner sep=2pt}
}


%Baumdiagramme von oben nach unten: \begin{tikzpicture}[baumdiagrammvertikal] \node[punkt]{} child{...}
%erstes child = linker Ast; Wahrscheinlichkeit links mit node[wkeit, left], rechts mit node[wkeit, right]
\tikzset{
  baumdiagrammvertikal/.style={grow=down, level distance=1.6cm,
    level 1/.style={sibling distance=4.8cm}, level 2/.style={sibling distance=2.4cm},
    level 3/.style={sibling distance=1.2cm}, edge from parent/.style={draw, thin}}
}
%----------------------------------------------------------------------------------------------------%









%----------------------------------------------------------------------------------------------------%
%%%%%%%%%%%                                     Handreichungen                                        %%%%%%%%%%%
%----------------------------------------------------------------------------------------------------%


%Schriftzug TikZ im Fließtext (Leerzeichen danach ist eingebaut)
\newcommand{\TikZ}{Ti\textit{k}Z }


%Oxidationszahl über ein Elementsymbol: $\oz{+II}{Cu}$, in \ce{...} zwischen Dollarzeichen: \ce{$\oz{0}{Zn}$ + ...}
\newcommand{\oz}[2]{\overset{\scriptscriptstyle\mathrm{#1}}{\mathrm{\ce{#2}}}}


%Umgebung passend: Inhalt (z.B. eine große tikzpicture) höchstens so breit wie die Zeile setzen.
%Ist er schmaler, bleibt er unverändert. \begin{passend} ... \end{passend} (Inhalt nicht als Argument, # erlaubt)
\newsavebox{\passendbox}
\newenvironment{passend}{\begin{lrbox}{\passendbox}}{\end{lrbox}%
	\ifdim\wd\passendbox>\linewidth\resizebox{\linewidth}{!}{\usebox{\passendbox}}\else\usebox{\passendbox}\fi}


%Box für die Ausgabe unter einem Code-Beispiel: erst lstlisting[style=latex], darunter \begin{ausgabe} ... \end{ausgabe}
\newtcolorbox{ausgabe}{%
	colback=blue!5,%
	colframe=blue!50!black,%
	title=Ausgabe,%
	fonttitle=\bfseries%
}
%----------------------------------------------------------------------------------------------------%









%----------------------------------------------------------------------------------------------------%
%%%%%%%%%%%                                        Lautschrift                                         %%%%%%%%%%%
%----------------------------------------------------------------------------------------------------%


%Lautschrift in eckigen Klammern mit \ipa{"a:b@nt} – schaltet das babel-Kürzel " vorher aus (sonst wird "a zu ä)
\makeatletter
\newcommand{\ipa}{\begingroup\shorthandoff{"}\ipa@aux}
\newcommand{\ipa@aux}[1]{[\textipa{#1}]\endgroup}
\makeatother
%----------------------------------------------------------------------------------------------------%









%----------------------------------------------------------------------------------------------------%
%%%%%%%%%%%                                       Sprechblasen                                         %%%%%%%%%%%
%----------------------------------------------------------------------------------------------------%


%Sprechblase mit Zipfel links unten (Lehrkraft) mit \lehrerblase[TikZ-Optionen]{Text} (benötigt shapes.callouts)
\newcommand{\lehrerblase}[2][]{%
  \tikz[baseline=(B.base)]{%
    \node[rectangle callout, rounded corners=6pt, callout relative pointer={(-13mm,-7mm)},
      draw, thick, fill=blue!6, align=left, inner sep=3.5mm, #1] (B) {#2};}%
}


%Sprechblase mit Zipfel rechts unten (Schülerin/Schüler) mit \schuelerblase[TikZ-Optionen]{Text}
\newcommand{\schuelerblase}[2][]{%
  \tikz[baseline=(B.base)]{%
    \node[rectangle callout, rounded corners=6pt, callout relative pointer={(13mm,-7mm)},
      draw, thick, fill=green!6, align=left, inner sep=3.5mm, #1] (B) {#2};}%
}
%----------------------------------------------------------------------------------------------------%









%----------------------------------------------------------------------------------------------------%
%%%%%%%%%%%                            TikZ-Stile für Programmablaufpläne                            %%%%%%%%%%%
%----------------------------------------------------------------------------------------------------%


%Ein-/Ausgabe als Parallelogramm, Operation als Rechteck, Verzweigung als Raute (benötigt shapes.geometric)
%align=center, damit \\ im Knotentext eine neue Zeile ergibt (ohne align wird \\ stillschweigend verschluckt)
\tikzstyle{input}=[trapezium, trapezium left angle=60, trapezium right angle=120, fill=green!20, draw=green, align=center]
\tikzstyle{operation}=[rectangle, fill=orange!20, draw=orange, align=center]
\tikzstyle{condition}=[diamond, fill=magenta!20, draw=magenta, align=center]
\tikzstyle{output}=[trapezium, trapezium left angle=60, trapezium right angle=120, fill=red!20, draw=red, align=center]
%----------------------------------------------------------------------------------------------------%









%----------------------------------------------------------------------------------------------------%
%%%%%%%%%%%                            TikZ-Stile für UML-Klassendiagramme                           %%%%%%%%%%%
%----------------------------------------------------------------------------------------------------%


%Klasse mit drei Teilen: Name, Attribute, Methoden (benötigt shapes.multipart)
\tikzstyle{umlclass} = [rectangle split, rectangle split parts=3, rectangle split part align={center,left,left}, draw, thick, text width=4cm, minimum height=2cm]


%Beziehungen (benötigt arrows.meta)
\tikzstyle{inherit} = [-{Triangle[open]}, thick]	%Vererbung
\tikzstyle{compose} = [-{Diamond}, thick]	%Komposition
\tikzstyle{aggregate} = [-{Diamond[open]}, thick]	%Aggregation
\tikzstyle{associate} = [-, thick]	%Assoziation
\tikzstyle{depend} = [-, thick, dashed]	%Abhängigkeit
%----------------------------------------------------------------------------------------------------%








%----------------------------------------------------------------------------------------------------%
%%%%%%%%%%%                                         Listing Settings                                       %%%%%%%%%%%%
%----------------------------------------------------------------------------------------------------%


%Umlaute und häufige Sonderzeichen in ALLEN Listings (Python, Java, ...), sonst Fehler »Invalid UTF-8 byte sequence«
\lstset{literate=%
	{ä}{{\"a}}1 {ö}{{\"o}}1 {ü}{{\"u}}1 {Ä}{{\"A}}1 {Ö}{{\"O}}1 {Ü}{{\"U}}1 {ß}{{\ss}}1%
	{·}{{$\cdot$}}1 {€}{{\euro}}1 {°}{{\degree}}1 {→}{{$\rightarrow$}}2 {≤}{{$\leq$}}1 {≥}{{$\geq$}}1 {≠}{{$\neq$}}1%
	{–}{{--}}1 {—}{{---}}1 {„}{{\glqq}}1 {“}{{\grqq}}1 {…}{{\dots}}3 {×}{{$\times$}}1 {²}{{$^2$}}1 {³}{{$^3$}}1 {π}{{$\pi$}}1%
}


%Einstellung der Codeumgebung für LaTeX-Code-Beispiele (Stil der Handreichungen: ohne Zeilennummern)
%Verwendung: \begin{lstlisting}[style=latex] ... \end{lstlisting}, darunter \begin{ausgabe} ... \end{ausgabe}
\lstdefinestyle{latex}{%
  	language=TeX,%
  	frame=single,%
  	backgroundcolor=\color{yellow!20},%
  	rulecolor=\color{lightgray!40},%
  	breaklines=true,%
  	basicstyle=\scriptsize\ttfamily,%
  	keywordstyle=\bfseries\color{purple!40!black},%
  	commentstyle=\itshape\color{green!40!black},%
  	identifierstyle=\color{blue},%
	stringstyle=\color{black},%
	captionpos=b,%
	tabsize=2,%
	showspaces=false,%
	showstringspaces=false,%
  	showtabs=false,%
	deletekeywords={document},%
	morekeywords={begin}%
}


%Variante mit Zeilennummern: \begin{lstlisting}[style=latexnummeriert]
\lstdefinestyle{latexnummeriert}{%
	style=latex,%
	numbers=left,%
	numberstyle=\tiny\color{gray!90!black},%
  	numbersep=-5pt%
}
%----------------------------------------------------------------------------------------------------%






%----------------------------------------------------------------------------------------------------%
%%%%%%%%%%%                          Formularfelder (ausfüllbare PDFs)                          %%%%%%%%%%%
%----------------------------------------------------------------------------------------------------%
% Nur innerhalb von \begin{Form} ... \end{Form} (hyperref) verwendbar. Für Optionsfelder (radio)
% im Dokument zusätzlich \hypersetup{unicode=false}, sonst bekommen sie kaputte Namen.
% Jedes Feld steckt in einer eigenen Box (\raisebox/\hbox): Ohne sie übernimmt pdfTeX die Höhe
% der umgebenden Zeile, in Tabellen werden die Felder dann so hoch wie die ganze Zeile.

%Aussehen aller Felder: blauer Rand, hellblauer Hintergrund, Schrift 11 pt, wird mitgedruckt
\renewcommand{\DefaultOptionsofText}{print, bordercolor={0 0.353 0.612}, backgroundcolor={0.94 0.96 1}, charsize=11pt}
\renewcommand{\DefaultOptionsofCheckBox}{print, bordercolor={0 0.353 0.612}, backgroundcolor={0.94 0.96 1}, charsize=11pt}
\renewcommand{\DefaultOptionsofRadio}{print, notoggletooff, bordercolor={0 0.353 0.612}, backgroundcolor={0.94 0.96 1}, charsize=11pt}
\renewcommand{\DefaultOptionsofPopdownBox}{print, bordercolor={0 0.353 0.612}, backgroundcolor={0.94 0.96 1}, charsize=11pt}
\renewcommand{\DefaultHeightofText}{0.6cm}

%\feld[Breite]{Name} – einzeiliges Textfeld, sitzt auf der Grundlinie (Standardbreite 2,4 cm)
\newcommand{\feld}[2][2.4cm]{\raisebox{-0.17cm}{\TextField[name=#2, width=#1]{}}}

%\zeichen{Name} – Feld für genau ein Zeichen, zentriert (z. B. E/V/A, Buchstabe einer Reihenfolge)
\newcommand{\zeichen}[1]{\raisebox{-0.12cm}{\TextField[name=#1, width=0.8cm, height=0.5cm, maxlen=1, align=1]{}}}

%\kreuz{Name} – Ankreuzkästchen
\newcommand{\kreuz}[1]{\raisebox{-0.08cm}{\CheckBox[name=#1, width=0.42cm, height=0.42cm]{}}}

%\antwort[Höhe]{Name} – mehrzeiliges Antwortfeld über die ganze Zeilenbreite (Standardhöhe 3,2 cm ≈ 4 Zeilen)
\newcommand{\antwort}[2][3.2cm]{\par\vspace{0.15cm}\noindent\hbox{\TextField[name=#2, multiline=true, width=\linewidth, height=#1]{}}}
%----------------------------------------------------------------------------------------------------%




%----------------------------------------------------------------------------------------------------%
%%%%%%%%%%%                                   Schach (farbige Bretter)                                   %%%%%%%%%%
%----------------------------------------------------------------------------------------------------%
% Die Standardschrift von skak schraffiert die dunklen Felder; Figur und Schraffur sind dort ein
% Zeichen und lassen sich nicht getrennt einfärben. Deshalb: Schrift skaknew in der Kodierung LSBC5
% (zeichnet nur die Figuren, weiße Figuren weiß gefüllt), die Felder malt chessboard als Flächen
% darunter. Eigene Markierungen mit colorbackfields liegen über den Feldern und unter den Figuren.

%Feldfarben der Brettschemata (hell/dunkel) und Markierungsfarbe für den letzten Zug
\definecolor{brettblauhell}{HTML}{EEF3FD} %Schema blau, helle Felder
\definecolor{brettblaudunkel}{HTML}{A9B4CC} %Schema blau, dunkle Felder
\definecolor{brettbraunhell}{HTML}{F0D9B5} %Schema braun (klassisches Holzbrett), helle Felder
\definecolor{brettbraundunkel}{HTML}{B58863} %Schema braun, dunkle Felder
\definecolor{brettgruenhell}{HTML}{EEEED2} %Schema gruen (Turnierbrett), helle Felder
\definecolor{brettgruendunkel}{HTML}{769656} %Schema gruen, dunkle Felder
\definecolor{brettgrauhell}{HTML}{F2F2F2} %Schema grau (für Schwarz-Weiß-Druck), helle Felder
\definecolor{brettgraudunkel}{HTML}{B4B4B4} %Schema grau, dunkle Felder
\definecolor{brettmarke}{HTML}{F6D55C} %Start- und Zielfeld eines Zuges hervorheben: pgfstyle=color, color=brettmarke, colorbackfields={e2,e4}

% Setzt die Voreinstellungen von \setchessboard zurück – eigene \setchessboard{...} deshalb DANACH schreiben.
%\schachbrettfarben{Schema} – Farben aller folgenden Schachbretter: braun (Standard), blau, gruen, grau oder klassisch (schraffiert wie in skak)
\makeatletter
\let\brett@grundeinstellung\board@val@savedkeylist
\newcommand{\schachbrettfarben}[1]{%
	\let\board@val@savedkeylist\brett@grundeinstellung
	\ifthenelse{\equal{#1}{klassisch}}{}%
	{\setchessboard{boardfontfamily=skaknew, boardfontencoding=LSBC5,
		piececolor=black, piecemaskcolor=white, setfontcolors,
		pgfstyle=color, opacity=1, color=brett#1hell, colorbackboard, color=brett#1dunkel, colorblackbackfields,
		color=black, labelfont=\sffamily\color{black!65}}}}
\makeatother
\schachbrettfarben{braun}
%----------------------------------------------------------------------------------------------------%


% Physikalisches Wörterbuch Deutsch – Englisch mit Lautschrift
% Skills: latex-grundstil, latex-lautschrift (tipa, \ipa), latex-physik (circuitikz, Federn, Linsen),
%         latex-funktionsgraphen, latex-tikz-grundlagen
% Kompilieren mit pdflatex (dreimal, wegen Inhaltsverzeichnis und Seitenzahlen im Register).

%----------------------------------------------------------------------------------------------------%
% Dokumentspezifische Befehle (nur für dieses Wörterbuch, deshalb lokal und nicht in newcommands.tex)
%----------------------------------------------------------------------------------------------------%

% \eintrag{Marke}{Artikel}{Begriff}{english term}{Lautschrift EN}{Formelzeichen/Skizze}
% Schaltet das babel-Kürzel " aus, BEVOR die Argumente gelesen werden – dadurch darf die Lautschrift
% wie bei \ipa{...} mit " (Hauptbetonung) und "" (Nebenbetonung) geschrieben werden.
\makeatletter
\newcommand\eintrag{\begingroup\shorthandoff{"}\eintrag@aux}
\newcommand\eintrag@aux[6]{\endgroup
	#2 \textbf{#3}\phantomsection\label{e:#1}
	& \textbf{#4}\newline{\small\color{black!70}\ipa{#5}}
	& #6 \\ \midrule}
% \lies{Ausdruck}{english reading}{deutsche Lesart}
\newcommand\lies{\begingroup\shorthandoff{"}\lies@aux}
\newcommand\lies@aux[3]{\endgroup #1 & #2\newline{\color{black!60}#3} \\[4.5pt]}
\makeatother

% Tabellenkopf für jedes Kapitel in Teil I (wird auf jeder Folgeseite wiederholt)
\newenvironment{woerter}{%
	\arrayrulecolor{black!20}%
	\begin{longtable}{>{\raggedright\arraybackslash}m{5.3cm} >{\raggedright\arraybackslash}m{6.1cm} >{\centering\arraybackslash}m{4.6cm}}
		\rowcolor{cornflower!15}\textbf{Deutsch} & \textbf{Englisch} & \textbf{Formelzeichen \,/\, Skizze} \\ \midrule
		\endhead}{%
	\end{longtable}\arrayrulecolor{black}}

% Liste ohne Linien und Rahmen (Teil II): links der Ausdruck, rechts Englisch (schwarz) und Deutsch (grau)
\newenvironment{lesen}{%
	\begin{longtable}{@{}>{\raggedright\arraybackslash}p{4.6cm} >{\raggedright\arraybackslash}p{12.1cm}@{}}}{%
	\end{longtable}\vspace{-1.2em}}

\newcommand{\anmerkung}[1]{\par\smallskip\noindent{\small\sloppy\textbf{Hinweis:} #1\par}\smallskip}
% Verhindert, dass eine Zwischenüberschrift allein am Seitenende steht
\newcommand{\platz}{\par\ifdim\dimexpr\pagegoal-\pagetotal\relax<5.5cm \newpage\fi}
% Einheit aufrecht im Mathemodus
\newcommand{\eh}[1]{\,\mathrm{#1}}

% Kleine Skizzen: Gegenstand blau, Hervorhebung (der Begriff selbst) rot, Hilfslinien gestrichelt
\tikzset{
	fig/.style={blue, very thick},
	mass/.style={red, thick},
	hilf/.style={dashed},
	achse/.style={-{Stealth[length=4pt]}, thin},
	kraft/.style={red, very thick, -{Stealth[length=5pt]}},
}
\newcommand{\skizze}[1]{\tikz[baseline=(current bounding box.center), scale=0.45, line cap=round,
	line join=round, every node/.style={font=\scriptsize}, outer sep=0pt]{#1}}
\newcommand{\achsen}[6]{\draw[achse] (#1,0) -- (#2,0) node[below left] {$#5$}; \draw[achse] (0,#3) -- (0,#4) node[below left] {$#6$};}
% Schaltzeichen (circuitikz, europäische Symbole)
\newcommand{\sz}[1]{\tikz[baseline=-0.6ex]{\ctikzset{bipoles/length=0.9cm}\draw (0,0) to[#1] (1.7,0);}}

% Messen und Auswerten
\newcommand{\skSkala}{\skizze{\draw[fig, thick] (0,0) rectangle (5,0.9);
	\foreach \x in {0,...,10} {\draw[blue] ({0.25+0.45*\x},0.9) -- ++(0,-0.3);}
	\foreach \x in {0,5,10} {\draw[blue, thick] ({0.25+0.45*\x},0.9) -- ++(0,-0.5) node[below, black] {\x};}}}
\newcommand{\messpunkte}{\foreach \x/\y in {0.5/0.45, 1/0.6, 1.5/1.15, 2/1.2, 2.5/1.75, 3/1.9} {\draw[blue, thick] (\x-0.08,\y-0.08) -- (\x+0.08,\y+0.08) (\x-0.08,\y+0.08) -- (\x+0.08,\y-0.08);}}
\newcommand{\skDiagramm}{\skizze{\achsen{-0.3}{3.6}{-0.3}{2.5}{x}{y}\messpunkte}}
\newcommand{\skAusgleich}{\skizze{\achsen{-0.3}{3.6}{-0.3}{2.5}{x}{y}\messpunkte \draw[mass] (0,0.1) -- (3.4,2.15);}}
\newcommand{\skProp}{\skizze{\achsen{-0.3}{3.4}{-0.3}{2.4}{s}{F}\draw[fig] (0,0) -- (3.1,2.1);}}
\newcommand{\skAntiprop}{\skizze{\achsen{-0.3}{3.4}{-0.3}{2.4}{V}{p}\draw[fig, domain=0.3:3.2, samples=30] plot (\x,{0.65/\x});}}

% Bewegungen
\newcommand{\skST}{\skizze{\achsen{-0.3}{3.4}{-0.3}{2.4}{t}{s}\draw[fig] (0,0) -- (3.1,2.1);}}
\newcommand{\skSTParabel}{\skizze{\achsen{-0.3}{3.4}{-0.3}{2.4}{t}{s}\draw[fig, domain=0:3, samples=30] plot (\x,{0.24*\x*\x});}}
\newcommand{\skVT}{\skizze{\achsen{-0.3}{3.4}{-0.3}{2.4}{t}{v}\draw[fig] (0,0) -- (3.1,2.1);}}
\newcommand{\skVTkonst}{\skizze{\achsen{-0.3}{3.4}{-0.3}{2.4}{t}{v}\draw[fig] (0,1.4) -- (3.1,1.4);}}
\newcommand{\skWurf}{\skizze{\draw[thick] (0,0) -- (4,0); \draw[thick] (0,0) -- (0,1.9) -- (-0.8,1.9);
	\draw[fig, domain=0:3.4, samples=30] plot (\x,{1.9-0.164*\x*\x}); \draw[kraft] (0,1.9) -- (1.1,1.9) node[above] {$v_0$};
	\fill[blue] (0,1.9) circle (3pt);}}
\newcommand{\skKreisbew}{\skizze{\draw[fig, hilf, thick] (0,0) circle (1.2); \fill (0,0) circle (1.5pt);
	\draw[thin] (0,0) -- node[below, inner sep=1.5pt] {$r$} (30:1.2); \fill[blue] (30:1.2) circle (3.5pt);
	\draw[kraft] (30:1.2) -- ++(120:1.2) node[above right, inner sep=1pt] {$v$};}}
\newcommand{\skBahn}{\skizze{\draw[fig, hilf, thick] (0,0) .. controls (1,1.6) and (2.2,-0.6) .. (3.6,0.9); \fill[blue] (3.6,0.9) circle (3.5pt);}}

% Kräfte
\newcommand{\klotz}[2]{\filldraw[fill=blue!15, draw=blue, very thick] (#1,#2) rectangle ++(1.2,0.8);}
\newcommand{\skKraft}{\skizze{\klotz{0}{0} \draw[kraft] (1.2,0.4) -- (3.2,0.4) node[above left] {$\vec{F}$};}}
\newcommand{\skAngriff}{\skizze{\klotz{0}{0} \draw[blue, thick, -{Stealth}] (1.2,0.4) -- (3.2,0.4); \fill[red] (1.2,0.4) circle (3.5pt);}}
\newcommand{\skGewicht}{\skizze{\klotz{0}{0} \draw[kraft] (0.6,0.4) -- (0.6,-1.4) node[right] {$\vec{F}_G$}; \fill (0.6,0.4) circle (2pt);}}
\newcommand{\skNormal}{\skizze{\draw[thick] (-0.8,0) -- (2,0); \klotz{0}{0} \draw[kraft] (0.6,0) -- (0.6,1.9) node[right] {$\vec{F}_N$};}}
\newcommand{\skReibung}{\skizze{\draw[thick] (-1.6,0) -- (3,0); \klotz{0}{0} \draw[blue, thick, -{Stealth}] (1.2,0.4) -- (2.8,0.4) node[above] {$\vec{F}$};
	\draw[kraft] (0,0.12) -- (-1.4,0.12) node[above] {$\vec{F}_R$};}}
\newcommand{\feder}[2]{\draw[blue, thick, decorate, decoration={zigzag, segment length=3.2pt, amplitude=3pt}] #1 -- #2;}
\newcommand{\skFeder}{\skizze{\draw[thick] (-0.7,2.4) -- (0.7,2.4); \feder{(0,2.4)}{(0,0.8)} \filldraw[fill=blue!15, draw=blue, very thick] (-0.4,0.1) rectangle (0.4,0.8);}}
\newcommand{\skDehnung}{\skizze{\draw[thick] (-0.7,2.4) -- (2.2,2.4); \feder{(0,2.4)}{(0,1.3)} \feder{(1.5,2.4)}{(1.5,0.5)}
	\filldraw[fill=blue!15, draw=blue, very thick] (1.1,-0.2) rectangle (1.9,0.5);
	\draw[hilf, thin] (-0.3,1.3) -- (2.6,1.3) (1.5,0.5) -- (2.6,0.5); \draw[mass, <->] (2.4,1.3) -- node[right] {$s$} (2.4,0.5);}}
\newcommand{\skResultierende}{\skizze{\draw[hilf, thin] (2.2,0) -- (2.9,1.5) -- (0.7,1.5);
	\draw[blue, very thick, -{Stealth}] (0,0) -- (2.2,0) node[below] {$\vec{F}_1$}; \draw[blue, very thick, -{Stealth}] (0,0) -- (0.7,1.5) node[left] {$\vec{F}_2$};
	\draw[kraft] (0,0) -- (2.9,1.5) node[right] {$\vec{F}_{\mathrm{res}}$};}}
\newcommand{\hebel}{\draw[fig] (-2,0.5) -- (1.4,0.5); \filldraw[fill=black!20, draw=black, thick] (0,0.45) -- (-0.35,-0.2) -- (0.35,-0.2) -- cycle;}
\newcommand{\skHebel}{\skizze{\hebel \draw[kraft] (-1.8,0.5) -- (-1.8,-0.5) node[left] {$F_1$}; \draw[kraft] (1.2,0.5) -- (1.2,-0.9) node[right] {$F_2$};}}
\newcommand{\skHebelarm}{\skizze{\hebel \draw[mass, |<->|] (-1.8,0.95) -- node[above] {$r_1$} (0,0.95); \draw[mass, |<->|] (0,0.95) -- node[above] {$r_2$} (1.2,0.95);}}
\newcommand{\skDrehpunkt}{\skizze{\hebel \fill[red] (0,0.5) circle (4pt);}}
\newcommand{\skEbene}{\skizze{\draw[thick] (0,0) -- (3.6,0) -- (3.6,1.8) -- cycle; \draw[mass] (1,0) arc[start angle=0, end angle=26.57, radius=1]; \node at (13:1.4) {$\alpha$};
	\begin{scope}[shift={(1.8,0.9)}, rotate=26.57] \filldraw[fill=blue!15, draw=blue, very thick] (-0.5,0) rectangle (0.5,0.7);\end{scope}}}
\newcommand{\skRolle}{\skizze{\draw[thick] (-0.8,2.5) -- (0.8,2.5) (0,2.5) -- (0,1.7); \draw[fig] (0,1.7) circle (0.55); \fill (0,1.7) circle (2pt);
	\draw[blue, thick] (-0.55,1.7) -- (-0.55,0.5); \filldraw[fill=blue!15, draw=blue, very thick] (-0.9,-0.1) rectangle (-0.2,0.5);
	\draw[blue, thick] (0.55,1.7) -- (0.55,0.9); \draw[kraft] (0.55,0.9) -- (0.55,-0.1) node[right] {$F$};}}
\newcommand{\skAuftrieb}{\skizze{\fill[blue!12] (-1.4,-1) rectangle (1.4,0.9); \draw[thick] (-1.4,1.3) -- (-1.4,-1) -- (1.4,-1) -- (1.4,1.3); \draw[blue, thick] (-1.4,0.9) -- (1.4,0.9);
	\filldraw[fill=white, draw=blue, very thick] (-0.4,-0.5) rectangle (0.4,0.2); \draw[kraft] (0,0.2) -- (0,1.6) node[right] {$F_A$};}}
\newcommand{\skDruck}{\skizze{\draw[thick] (-1.5,0) -- (1.5,0); \filldraw[fill=blue!15, draw=blue, very thick] (-0.8,0) rectangle (0.8,0.8);
	\draw[kraft] (0,2) -- (0,0.85) node[pos=0.1, right] {$F$}; \draw[red, line width=2pt] (-0.8,0) -- (0.8,0); \node[below] at (0,0) {$A$};}}

% Wärme
\newcommand{\skThermometer}{\skizze{\draw[fig, thick, rounded corners=2pt] (-0.15,0.3) rectangle (0.15,2.8); \filldraw[fill=red, draw=blue, thick] (0,0.1) circle (0.33);
	\fill[red] (-0.07,0.3) rectangle (0.07,1.7); \foreach \y in {0.8,1.2,...,2.5} {\draw[thin] (0.15,\y) -- (0.4,\y);}}}
\newcommand{\teilchen}[1]{\foreach \p in {#1} {\fill[blue] \p circle (2.6pt);}}
\newcommand{\skFest}{\skizze{\draw[thin] (0,0) rectangle (1.9,1.5); \teilchen{(0.3,0.25),(0.75,0.25),(1.2,0.25),(1.65,0.25),(0.3,0.7),(0.75,0.7),(1.2,0.7),(1.65,0.7),(0.3,1.15),(0.75,1.15),(1.2,1.15),(1.65,1.15)}}}
\newcommand{\skFluessig}{\skizze{\draw[thin] (0,0) rectangle (1.9,1.5); \teilchen{(0.25,0.22),(0.7,0.3),(1.15,0.2),(1.6,0.28),(0.4,0.68),(0.9,0.72),(1.4,0.62),(0.22,1.05),(0.75,1.1),(1.25,1.02),(1.68,0.95)}}}
\newcommand{\skGas}{\skizze{\draw[thin] (0,0) rectangle (1.9,1.5); \teilchen{(0.3,0.3),(1.5,0.5),(0.8,1.2),(1.6,1.25),(0.4,0.95)}}}

% Elektrizität und Magnetismus
\newcommand{\skStromkreis}{\tikz[baseline=(current bounding box.center)]{\ctikzset{bipoles/length=0.8cm}
	\draw (0,0) to[battery1] (0,1.3) -- (1.9,1.3) to[lamp] (1.9,0) to[nos] (0,0);}}
\newcommand{\skReihe}{\tikz[baseline=(current bounding box.center)]{\ctikzset{bipoles/length=0.75cm}
	\draw (0,0) to[battery1] (0,1.2) to[R, l=\scriptsize$R_1$] (1.5,1.2) to[R, l=\scriptsize$R_2$] (3,1.2) -- (3,0) -- (0,0);}}
\newcommand{\skParallel}{\tikz[baseline=(current bounding box.center)]{\ctikzset{bipoles/length=0.75cm}
	\draw (0,0) to[battery1] (0,1.3) -- (1.3,1.3) to[R, l_=\scriptsize$R_1$, *-*] (1.3,0) -- (0,0);
	\draw (1.3,1.3) -- (2.6,1.3) to[R, l=\scriptsize$R_2$] (2.6,0) -- (1.3,0);}}
\newcommand{\skMagnet}{\skizze{\filldraw[fill=red!60, draw=black, thick] (-1.3,-0.35) rectangle (0,0.35); \filldraw[fill=green!50!black!60, draw=black, thick] (0,-0.35) rectangle (1.3,0.35);
	\node at (-0.65,0) {N}; \node at (0.65,0) {S};}}
\newcommand{\skMagnetfeld}{\skizze{\foreach \h in {0.9,1.5} {\draw[blue, thick, -{Stealth[length=3pt]}] (-1.1,0.35) .. controls (-1.6,\h+0.3) and (1.6,\h+0.3) .. (1.1,0.35);
		\draw[blue, thick, -{Stealth[length=3pt]}] (-1.1,-0.35) .. controls (-1.6,-\h-0.3) and (1.6,-\h-0.3) .. (1.1,-0.35);}
	\filldraw[fill=red!60, draw=black, thick] (-1.3,-0.35) rectangle (0,0.35); \filldraw[fill=green!50!black!60, draw=black, thick] (0,-0.35) rectangle (1.3,0.35);
	\node at (-0.65,0) {N}; \node at (0.65,0) {S};}}
\newcommand{\skEFeld}{\skizze{\draw[red, line width=2pt] (0,0) -- (0,2) node[above] {$+$}; \draw[blue, line width=2pt] (2.4,0) -- (2.4,2) node[above] {$-$};
	\foreach \y in {0.3,0.77,1.23,1.7} {\draw[thick, -{Stealth[length=3pt]}] (0.1,\y) -- (2.25,\y);}}}

% Schwingungen und Wellen
\newcommand{\skPendel}{\skizze{\draw[thick] (-0.9,2.6) -- (0.9,2.6); \draw[hilf, thin] (0,2.6) -- (0,0.2); \draw[blue, thick] (0,2.6) -- ++(-65:2.3) coordinate (m);
	\fill[blue] (m) circle (5pt); \draw[mass] (0,1.6) arc[start angle=-90, end angle=-65, radius=1]; \draw[hilf, thin, blue] (0,2.6) -- ++(-115:2.3);}}
\newcommand{\sinus}{\achsen{-0.3}{4.6}{-1.3}{1.5}{}{}\draw[fig, domain=0:4.2, samples=60] plot (\x,{sin(180*\x)});}
\newcommand{\skAmplitude}{\skizze{\sinus \draw[mass, <->] (0.5,0) -- node[right, inner sep=1.5pt] {$\hat{y}$} (0.5,1);}}
\newcommand{\skPeriode}{\skizze{\sinus \draw[mass, |<->|] (0.5,1.3) -- node[above] {$T$} (2.5,1.3); \node[below] at (4.4,0) {$t$};}}
\newcommand{\skWellenlaenge}{\skizze{\sinus \draw[mass, |<->|] (0.5,1.3) -- node[above] {$\lambda$} (2.5,1.3); \node[below] at (4.4,0) {$x$};}}
\newcommand{\skBergTal}{\skizze{\sinus \fill[red] (0.5,1) circle (3pt) (2.5,1) circle (3pt); \filldraw[fill=white, draw=red, thick] (1.5,-1) circle (3pt) (3.5,-1) circle (3pt);}}
\newcommand{\skGedaempft}{\skizze{\achsen{-0.3}{4.6}{-1.3}{1.5}{t}{y}\draw[fig, domain=0:4.2, samples=80] plot (\x,{1.2*exp(-0.55*\x)*cos(360*\x)});}}
\newcommand{\skStehend}{\skizze{\draw[thick] (0,-1) -- (0,1) (4,-1) -- (4,1); \draw[fig, domain=0:4, samples=50] plot (\x,{0.8*sin(90*\x)});
	\draw[fig, hilf, thick, domain=0:4, samples=50] plot (\x,{-0.8*sin(90*\x)});}}
\newcommand{\skLaengswelle}{\skizze{\foreach \k in {0,...,24} {\pgfmathsetmacro{\x}{0.17*\k+0.12*sin(30*\k)} \draw[blue, thick] (\x,-0.5) -- (\x,0.5);}}}

% Optik
\newcommand{\skReflexion}{\skizze{\draw[thick] (-2,0) -- (2,0); \foreach \x in {-1.9,-1.6,...,1.9} {\draw[thin] (\x,0) -- ++(-0.15,-0.2);}
	\draw[hilf, thin] (0,0) -- (0,2); \draw[fig, -{Stealth}] (130:2.3) -- (130:1.1); \draw[fig] (130:1.2) -- (0,0); \draw[fig, -{Stealth}] (0,0) -- (50:1.3); \draw[fig] (50:1.2) -- (50:2.3);
	\draw[mass] (90:0.9) arc[start angle=90, end angle=130, radius=0.9]; \draw[mass] (90:0.9) arc[start angle=90, end angle=50, radius=0.9];
	\node at (110:1.25) {$\alpha$}; \node at (70:1.25) {$\alpha'$};}}
\newcommand{\skBrechung}{\skizze{\fill[blue!12] (-2,-1.7) rectangle (2,0); \draw[thick] (-2,0) -- (2,0); \draw[hilf, thin] (0,-1.7) -- (0,1.8);
	\draw[fig, -{Stealth}] (130:2.2) -- (130:1); \draw[fig] (130:1.1) -- (0,0) -- (-65:1.9);
	\draw[mass] (90:0.8) arc[start angle=90, end angle=130, radius=0.8]; \draw[mass] (-90:0.9) arc[start angle=-90, end angle=-65, radius=0.9];
	\node at (110:1.15) {$\alpha$}; \node at (-78:1.3) {$\beta$};}}
\newcommand{\skLot}{\skizze{\draw[thick] (-2,0) -- (2,0); \foreach \x in {-1.9,-1.6,...,1.9} {\draw[thin] (\x,0) -- ++(-0.15,-0.2);}
	\draw[red, very thick, hilf] (0,0) -- (0,2); \draw[fig] (130:2.3) -- (0,0) -- (50:2.3);}}
\newcommand{\linse}{\draw[fig, fill=blue!12] (0,1.3) arc[start angle=140, end angle=220, radius=2.022] arc[start angle=-40, end angle=40, radius=2.022];}
\newcommand{\skSammellinse}{\skizze{\draw[thin, dash dot] (-2.2,0) -- (3,0); \linse
	\foreach \y in {0.7,-0.7} {\draw[red, thick] (-2.2,\y) -- (0,\y) -- (2.8,-\y);} \fill (1.4,0) circle (2pt) node[below] {$F$};}}
\newcommand{\skZerstreuungslinse}{\skizze{\draw[thin, dash dot] (-2.4,0) -- (2.8,0);
	\draw[fig, fill=blue!12] (-0.45,1.3) -- (0.45,1.3) arc[start angle=150, end angle=210, radius=2.6] -- (-0.45,-1.3) arc[start angle=-30, end angle=30, radius=2.6];
	\foreach \s in {1,-1} {\draw[red, thick] (-2.4,0.7*\s) -- (0,0.7*\s) -- (2.6,1.7*\s); \draw[red, thin, hilf] (0,0.7*\s) -- (-1.8,0);} \fill (-1.8,0) circle (2pt) node[below] {$F$};}}
\newcommand{\skBrennweite}{\skizze{\draw[thin, dash dot] (-1.2,0) -- (3,0); \linse \fill (1.9,0) circle (2pt) node[above] {$F$}; \draw[mass, |<->|] (0,-0.6) -- node[below] {$f$} (1.9,-0.6);}}
\newcommand{\skPrisma}{\skizze{\draw[fig, fill=blue!10] (0,0) -- (2.4,0) -- (1.2,2.08) -- cycle; \draw[thick] (-1.6,0.6) -- (0.6,1.04);
	\draw[red, thick] (0.6,1.04) -- (1.75,1.12) -- (3.4,0.7); \draw[green!60!black, thick] (0.6,1.04) -- (1.82,1.0) -- (3.4,0.4); \draw[blue!70!red, thick] (0.6,1.04) -- (1.89,0.88) -- (3.4,0.1);}}
\newcommand{\skSchatten}{\skizze{\fill[black!25] (1.5,0.5) -- (4,1.25) -- (4,-1.25) -- (1.5,-0.5) -- cycle; \fill[yellow!80!orange, draw=black] (0,0) circle (0.22);
	\draw[thin] (0,0) -- (4,1.25) (0,0) -- (4,-1.25); \filldraw[fill=blue!40, draw=blue, thick] (1.5,0) circle (0.5); \draw[thick] (4,-1.6) -- (4,1.6);}}

% Atom- und Kernphysik
\newcommand{\skAtom}{\skizze{\draw[thin] (0,0) circle (0.75) circle (1.3); \fill[red] (0,0) circle (0.25); \fill[blue] (40:0.75) circle (3pt) (200:0.75) circle (3pt) (100:1.3) circle (3pt) (300:1.3) circle (3pt);}}
\newcommand{\skKern}{\skizze{\draw[thin, black!40] (0,0) circle (1.2); \foreach \w/\c in {0/red, 60/black!45, 120/red, 180/black!45, 240/red, 300/black!45} {\fill[\c] (\w:0.3) circle (0.17);} \fill[black!45] (0,0) circle (0.17);}}
\newcommand{\skHuelle}{\skizze{\draw[red, very thick] (0,0) circle (0.75) circle (1.3); \fill[black!45] (0,0) circle (0.25);}}
\newcommand{\skZerfall}{\skizze{\achsen{-0.3}{4.4}{-0.3}{2.6}{t}{N}\draw[fig, domain=0:4.1, samples=40] plot (\x,{1.9*exp(-0.693*\x/1.2)});
	\draw[mass, hilf] (0,0.95) -- (1.2,0.95) -- (1.2,0) node[below] {$T_{1/2}$}; \node[left] at (0,1.9) {$N_0$};}}
\newcommand{\skNiveaus}{\skizze{\foreach \y/\n in {0/1, 1.3/2, 1.9/3, 2.2/4} {\draw[fig, thick] (0,\y) -- (2.4,\y) node[right, black] {$n=\n$};}
	\draw[mass, -{Stealth}, decorate, decoration={snake, amplitude=1.2pt, segment length=4pt, post length=3pt}] (1.2,1.3) -- (1.2,0.05);}}
\newcommand{\skSpaltung}{\skizze{\fill[black!45] (-2,0) circle (0.14); \draw[thin, -{Stealth}] (-1.8,0) -- (-1,0); \filldraw[fill=red!35, draw=red, thick] (0,0) circle (0.7);
	\draw[thin, -{Stealth}] (0.9,0.3) -- (1.6,0.7); \draw[thin, -{Stealth}] (0.9,-0.3) -- (1.6,-0.7); \filldraw[fill=red!35, draw=red, thick] (2.1,0.95) circle (0.45) (2.1,-0.95) circle (0.38);}}


\title{Physikalisches Wörterbuch\\[0.3em] \large Deutsch -- Englisch mit englischer Lautschrift, Formelzeichen und Skizzen}
\author{Martin Quandt}
\date{\today}
\pagestyle{plain}
\setlength{\tabcolsep}{4pt}


% Zweisprachige Überschriften: \teil, \kapitel, \zwischen{Deutsch}{English}{Lautschrift}
% Wie \eintrag wird " ausgeschaltet, BEVOR die Argumente gelesen werden (Betonungszeichen der Lautschrift).
% Ins Inhaltsverzeichnis und in die PDF-Lesezeichen kommt nur „Deutsch -- English“ ohne Lautschrift.
\makeatletter
\newcommand\teil{\begingroup\shorthandoff{"}\teil@aux}
\newcommand\teil@aux[3]{\endgroup\part[#1 -- #2]{\raggedright #1\\[5pt]{\normalfont\Large\color{black!55}#2\hspace{0.35em}\penalty0 \ipa{#3}}}}
\newcommand\kapitel{\begingroup\shorthandoff{"}\kapitel@aux}
\newcommand\kapitel@aux[3]{\endgroup\section[#1 -- #2]{\raggedright #1\\[0pt]{\normalfont\small\color{black!55}#2\hspace{0.35em}\penalty0 \ipa{#3}}}}
\newcommand\zwischen{\begingroup\shorthandoff{"}\zwischen@aux}
\newcommand\zwischen@aux[3]{\endgroup\subsection*{\raggedright #1 {\normalfont\small\color{black!55}\,--\, #2\hspace{0.35em}\penalty0 \ipa{#3}}}}
% Die zweite Überschriftenzeile kostet Platz: Zeilen in den Listen und Tabellen etwas enger setzen
\renewcommand\lies@aux[3]{\endgroup #1 & #2\newline{\color{black!60}#3} \\[1.5pt]}
\makeatother
\setlength{\aboverulesep}{0.15ex}\setlength{\belowrulesep}{0.3ex}\renewcommand{\arraystretch}{0.93}

\begin{document}

\maketitle

Dieses Wörterbuch enthält 295 physikalische Fachbegriffe aus der Sekundarstufe I und II,
geordnet nach Sachgebieten. Zu jedem Begriff stehen der englische Fachausdruck mit seiner Aussprache und
-- wo es hilft -- das Formelzeichen mit Einheit, eine Formel oder eine kleine Skizze. Teil~II zeigt, wie
man Einheiten, Größenangaben und Formeln auf Englisch vorliest und wie man Experimente und Diagramme
beschreibt. Am Ende finden Sie je ein alphabetisches Register Deutsch -- Englisch und Englisch -- Deutsch.

\zwischen{Hinweise zur Lautschrift}{Notes on the phonetic transcription}{"n@Uts 6n D@ f@"netIk tr\ae n"skrIpS\textsyllabic{n}}

Die Lautschrift folgt dem Internationalen Phonetischen Alphabet (IPA) und gibt die britische
Aussprache wieder, wie sie in Schulwörterbüchern steht (dort wird \textipa{r} statt
\textipa{\textturnr} geschrieben). Der Strich \textipa{\textprimstress} steht \textit{vor} der
betonten Silbe, \textipa{\textsecstress} vor einer Silbe mit Nebenbetonung, \textipa{:} bedeutet lang.

\medskip
{\small
\begin{tabular}{>{\centering\arraybackslash}p{0.9cm} l @{\hspace{1.2cm}} >{\centering\arraybackslash}p{0.9cm} l}
	\toprule
	\textbf{Laut} & \textbf{Beispiel} & \textbf{Laut} & \textbf{Beispiel} \\ \midrule
	\textipa{@} & pow\textbf{er} \ipa{"paU@} & \textipa{T} & \textbf{th}ermal \ipa{"T3:m\textsyllabic{l}} \\
	\textipa{\ae} & m\textbf{a}ss \ipa{m\ae s} & \textipa{S} & pre\textbf{ss}ure \ipa{"preS@} \\
	\textipa{2} & c\textbf{u}rrent \ipa{"k2r@nt} & \textipa{Z} & mea\textbf{s}ure \ipa{"meZ@} \\
	\textipa{6} & v\textbf{o}lume \ipa{"v6lju:m} & \textipa{dZ} & \textbf{j}oule \ipa{dZu:l} \\
	\textipa{3:} & w\textbf{or}k \ipa{w3:k} & \textipa{N} & spri\textbf{ng} \ipa{sprIN} \\
	\textipa{@U} & \textbf{oh}m \ipa{@Um} & \textipa{\textsyllabic{l}} & partic\textbf{le} \ipa{"pA:tIk\textsyllabic{l}} \\
	\bottomrule
\end{tabular}}

\medskip
\textbf{Dritte Spalte:} Das Formelzeichen ist das im deutschen Unterricht übliche; weicht das
englische Zeichen ab, steht es klein darunter. \textbf{Farben in den Skizzen:} Gegenstände sind blau,
hervorgehobene Teile (der Begriff selbst, Kräfte, Maße) rot, Hilfslinien gestrichelt.

\newpage
{\footnotesize\setlength{\parskip}{0pt}\tableofcontents}
\newpage

\teil{Wörterbuch}{Dictionary}{"dIkS@n@ri}

%%%%  GRÖSSEN, EINHEITEN UND MESSEN  %%%%
\kapitel{Größen, Einheiten und Messen}{Quantities, units and measurement}{"kw6nt@tiz, "ju:nIts @n "meZ@m@nt}

\begin{woerter}
\eintrag{1}{die}{physikalische Größe}{physical quantity}{"fIzIk\textsyllabic{l} "kw6nt@ti}{$l = 2{,}5\eh{m}$}
\eintrag{2}{der}{Zahlenwert}{numerical value}{nju:"merIk\textsyllabic{l} "v\ae lju:}{$l = \frot{2{,}5}\eh{m}$}
\eintrag{3}{die}{Einheit}{unit}{"ju:nIt}{$l = 2{,}5\,\frot{\mathrm{m}}$}
\eintrag{4}{die}{Basiseinheit}{base unit}{"beIs ""ju:nIt}{m, kg, s, A, K, mol, cd}
\eintrag{5}{die}{abgeleitete Einheit}{derived unit}{dI"raIvd "ju:nIt}{$1\eh{N} = 1\eh{kg \cdot m/s^2}$}
\eintrag{6}{das}{Formelzeichen}{symbol}{"sImb\textsyllabic{l}}{$F;\ m;\ a$}
\eintrag{7}{der}{Vorsatz}{prefix}{"pri:fIks}{k, M, m, $\mu$}
\eintrag{8}{}{messen}{measure}{"meZ@}{}
\eintrag{9}{die}{Messung}{measurement}{"meZ@m@nt}{}
\eintrag{10}{der}{Messwert}{reading, measured value}{"ri:dIN, "meZ@d "v\ae lju:}{}
\eintrag{11}{das}{Messgerät}{measuring instrument}{"meZ@rIN "Instr@m@nt}{}
\eintrag{12}{die}{Skala}{scale}{skeIl}{\skSkala}
\eintrag{13}{die}{Messgenauigkeit}{accuracy}{"\ae kj@r@si}{}
\eintrag{14}{die}{Messunsicherheit}{uncertainty (of measurement)}{2n"s3:t\textsyllabic{n}ti}{$l = (2{,}50 \pm 0{,}01)\eh{m}$}
\eintrag{15}{der}{Mittelwert}{mean (value)}{mi:n}{$\bar{x}$}
\eintrag{16}{der}{Versuch, das Experiment}{experiment}{Ik"sperIm@nt}{}
\eintrag{17}{der}{Versuchsaufbau}{experimental set-up}{Ik""sperI"ment\textsyllabic{l} "set2p}{}
\eintrag{18}{die}{Beobachtung}{observation}{""6bz@"veIS\textsyllabic{n}}{}
\eintrag{19}{die}{Auswertung}{analysis, evaluation}{@"n\ae l@sIs, I""v\ae lju"eIS\textsyllabic{n}}{}
\eintrag{20}{die}{Vermutung, die Hypothese}{hypothesis}{haI"p6T@sIs}{}
\eintrag{21}{das}{Gesetz}{law}{lO:}{}
\eintrag{22}{das}{Modell}{model}{"m6d\textsyllabic{l}}{}
\eintrag{23}{das}{Diagramm}{graph}{grA:f}{\skDiagramm}
\eintrag{24}{die}{Ausgleichsgerade}{line of best fit}{"laIn @v ""best "fIt}{\skAusgleich}
\eintrag{25}{}{proportional}{directly proportional}{d@"rektli pr@"pO:S@n\textsyllabic{l}}{$F \sim s$ \newline \skProp}
\eintrag{26}{}{umgekehrt proportional}{inversely proportional}{In"v3:sli pr@"pO:S@n\textsyllabic{l}}{$p \sim \dfrac{1}{V}$ \newline \skAntiprop}
\eintrag{27}{die}{Konstante}{constant}{"k6nst@nt}{}
\eintrag{28}{der}{Körper}{body, object}{"b6di, "6bdZIkt}{}
\eintrag{29}{der}{Stoff}{substance, material}{"s2bst@ns, m@"tI@ri@l}{}
\end{woerter}

\platz\zwischen{Beispiele}{Examples}{Ig"zA:mp\textsyllabic{l}z}

\begin{lesen}
\lies{$l = 2{,}5\eh{m}$}{The length is two point five metres. Two point five is the numerical value, the metre is the unit.}{Die Länge beträgt zwei Komma fünf Meter. Zwei Komma fünf ist der Zahlenwert, Meter die Einheit.}
\lies{\skSkala}{Read the scale of the measuring instrument. Take the reading at eye level.}{Lies die Skala des Messgeräts ab. Lies den Messwert auf Augenhöhe ab.}
\lies{$t_1;\ t_2;\ \ldots;\ t_5$ \newline $\bar{t} = 1{,}42\eh{s}$}{Repeat the measurement five times and calculate the mean.}{Wiederhole die Messung fünfmal und berechne den Mittelwert.}
\lies{$F \sim s$}{The force is directly proportional to the extension: if the extension doubles, the force doubles.}{Die Kraft ist proportional zur Verlängerung: Verdoppelt sich die Verlängerung, so verdoppelt sich die Kraft.}
\lies{\skAusgleich}{Plot the readings and draw the line of best fit.}{Trage die Messwerte in ein Diagramm ein und zeichne die Ausgleichsgerade.}
\end{lesen}


%%%%  BEWEGUNGEN  %%%%
\newpage
\kapitel{Bewegungen}{Motion}{"m@US\textsyllabic{n}}

\begin{woerter}
\eintrag{30}{die}{Bewegung}{motion}{"m@US\textsyllabic{n}}{}
\eintrag{31}{die}{Ruhe}{rest}{rest}{$v = 0$}
\eintrag{32}{das}{Bezugssystem}{frame of reference}{"freIm @v "refr@ns}{}
\eintrag{33}{die}{Bahn}{path, trajectory}{pA:T, tr@"dZekt@ri}{\skBahn}
\eintrag{34}{der}{Ort}{position}{p@"zIS\textsyllabic{n}}{$x$ in $\mathrm{m}$}
\eintrag{35}{der}{Weg}{distance}{"dIst@ns}{$s$ in $\mathrm{m}$}
\eintrag{36}{die}{Zeit}{time}{taIm}{$t$ in $\mathrm{s}$}
\eintrag{37}{die}{Geschwindigkeit}{velocity, speed}{v@"l6s@ti, spi:d}{$v$ in $\mathrm{m/s}$ \newline $v = \dfrac{s}{t}$}
\eintrag{38}{die}{Durchschnittsgeschwindigkeit}{average speed}{"\ae v@rIdZ spi:d}{$\bar{v} = \dfrac{\Delta s}{\Delta t}$}
\eintrag{39}{die}{Momentangeschwindigkeit}{instantaneous velocity}{""Inst@n"teIni@s v@"l6s@ti}{}
\eintrag{40}{die}{Beschleunigung}{acceleration}{@k""sel@"reIS\textsyllabic{n}}{$a$ in $\mathrm{m/s^2}$ \newline $a = \dfrac{\Delta v}{\Delta t}$}
\eintrag{41}{die}{Verzögerung}{deceleration}{""di:sel@"reIS\textsyllabic{n}}{$a < 0$}
\eintrag{42}{die}{gleichförmige Bewegung}{uniform motion}{"ju:nIfO:m "m@US\textsyllabic{n}}{$s = v \cdot t$ \newline \skVTkonst}
\eintrag{43}{die}{gleichmäßig beschleunigte Bewegung}{uniformly accelerated motion}{"ju:nIfO:mli @k"sel@reItId "m@US\textsyllabic{n}}{$s = \dfrac{a}{2} \cdot t^2$ \newline \skSTParabel}
\eintrag{44}{das}{Weg-Zeit-Diagramm}{distance–time graph}{"dIst@ns "taIm grA:f}{\skST}
\eintrag{45}{das}{Geschwindigkeit-Zeit-Diagramm}{velocity–time graph}{v@"l6s@ti "taIm grA:f}{\skVT}
\eintrag{46}{der}{freie Fall}{free fall}{""fri: "fO:l}{$s = \dfrac{g}{2} \cdot t^2$}
\eintrag{47}{die}{Fallbeschleunigung}{acceleration due to gravity}{@k""sel@"reIS\textsyllabic{n} dju: t@ "gr\ae v@ti}{$g = 9{,}81\eh{m/s^2}$}
\eintrag{48}{der}{waagerechte Wurf}{horizontal projection}{""h6rI"z6nt\textsyllabic{l} pr@"dZekS\textsyllabic{n}}{\skWurf}
\eintrag{49}{die}{Wurfweite}{range}{reIndZ}{}
\eintrag{50}{die}{Kreisbewegung}{circular motion}{"s3:kj@l@ "m@US\textsyllabic{n}}{\skKreisbew}
\eintrag{51}{die}{Umlaufzeit}{period}{"pI@ri@d}{$T$ in $\mathrm{s}$}
\eintrag{52}{die}{Frequenz}{frequency}{"fri:kw@nsi}{$f$ in $\mathrm{Hz}$ \newline $f = \dfrac{1}{T}$}
\eintrag{53}{die}{Drehzahl}{rotational speed}{r@U"teIS@n\textsyllabic{l} spi:d}{$n$ in $\mathrm{1/min}$}
\eintrag{54}{die}{Winkelgeschwindigkeit}{angular velocity}{"\ae Ngj@l@ v@"l6s@ti}{$\omega$ in $\mathrm{1/s}$ \newline $\omega = \dfrac{2\pi}{T}$}
\eintrag{55}{die}{Bahngeschwindigkeit}{tangential speed}{t\ae n"dZenS\textsyllabic{l} spi:d}{$v = \omega \cdot r$}
\end{woerter}

\platz\zwischen{Beispiele}{Examples}{Ig"zA:mp\textsyllabic{l}z}

\begin{lesen}
\lies{$s = 100\eh{m},\ t = 12{,}5\eh{s}$ \newline $v = \dfrac{s}{t} = 8\eh{m/s}$}{A sprinter covers one hundred metres in twelve point five seconds. Her average speed is eight metres per second.}{Eine Sprinterin legt hundert Meter in zwölf Komma fünf Sekunden zurück. Ihre Durchschnittsgeschwindigkeit beträgt acht Meter pro Sekunde.}
\lies{$a = \dfrac{20\eh{m/s}}{5\eh{s}} = 4\eh{m/s^2}$}{The car accelerates from rest to twenty metres per second in five seconds. Its acceleration is four metres per second squared.}{Das Auto beschleunigt aus der Ruhe in fünf Sekunden auf zwanzig Meter pro Sekunde. Seine Beschleunigung beträgt vier Meter pro Sekunde zum Quadrat.}
\lies{\skVTkonst}{The body moves at constant velocity: the velocity–time graph is a horizontal line.}{Der Körper bewegt sich gleichförmig: Das Geschwindigkeit-Zeit-Diagramm ist eine waagerechte Gerade.}
\lies{$s = \dfrac{g}{2} \cdot t^2 = 4{,}9\eh{m/s^2} \cdot (2\eh{s})^2$ \newline $s \approx 19{,}6\eh{m}$}{A stone is dropped from a bridge. After two seconds it has fallen about nineteen point six metres.}{Ein Stein wird von einer Brücke fallen gelassen. Nach zwei Sekunden ist er etwa neunzehn Komma sechs Meter gefallen.}
\lies{$T = 0{,}5\eh{s}$ \newline $f = \dfrac{1}{T} = 2\eh{Hz}$}{The wheel completes one revolution in half a second, so its frequency is two hertz.}{Das Rad braucht für eine Umdrehung eine halbe Sekunde, seine Frequenz beträgt also zwei Hertz.}
\end{lesen}


%%%%  KRÄFTE UND IHRE WIRKUNGEN  %%%%
\newpage
\kapitel{Kräfte und ihre Wirkungen}{Forces and their effects}{"fO:sIz @n De@r I"fekts}

\begin{woerter}
\eintrag{56}{die}{Masse}{mass}{m\ae s}{$m$ in $\mathrm{kg}$}
\eintrag{57}{das}{Volumen}{volume}{"v6lju:m}{$V$ in $\mathrm{m^3}$}
\eintrag{58}{die}{Dichte}{density}{"dens@ti}{$\varrho$ in $\mathrm{kg/m^3}$ \newline $\varrho = \dfrac{m}{V}$}
\eintrag{59}{die}{Kraft}{force}{fO:s}{$F$ in $\mathrm{N}$ \newline \skKraft}
\eintrag{60}{der}{Angriffspunkt}{point of application}{"pOInt @v ""\ae plI"keIS\textsyllabic{n}}{\skAngriff}
\eintrag{61}{der}{Betrag (einer Kraft)}{magnitude}{"m\ae gnItju:d}{$|\vec{F}| = F$}
\eintrag{62}{die}{Richtung}{direction}{d@"rekS\textsyllabic{n}}{}
\eintrag{63}{die}{Gewichtskraft}{weight}{weIt}{$F_G = m \cdot g$ \newline \skGewicht}
\eintrag{64}{die}{Normalkraft}{normal force}{"nO:m\textsyllabic{l} fO:s}{\skNormal}
\eintrag{65}{die}{Reibungskraft}{friction, frictional force}{"frIkS\textsyllabic{n}}{$F_R = \mu \cdot F_N$ \newline \skReibung}
\eintrag{66}{die}{Reibungszahl}{coefficient of friction}{""k@UI"fIS\textsyllabic{n}t @v "frIkS\textsyllabic{n}}{$\mu$}
\eintrag{67}{die}{Haftreibung}{static friction}{"st\ae tIk "frIkS\textsyllabic{n}}{}
\eintrag{68}{die}{Gleitreibung}{sliding friction}{"slaIdIN "frIkS\textsyllabic{n}}{}
\eintrag{69}{der}{Luftwiderstand}{air resistance, drag}{"e@ rI""zIst@ns, dr\ae g}{}
\eintrag{70}{die}{Feder}{spring}{sprIN}{\skFeder}
\eintrag{71}{die}{Dehnung, die Verlängerung}{extension}{Ik"stenS\textsyllabic{n}}{\skDehnung}
\eintrag{72}{die}{Federkonstante}{spring constant}{"sprIN ""k6nst@nt}{$D$ in $\mathrm{N/m}$ \newline \footnotesize engl. $k$}
\eintrag{73}{das}{hookesche Gesetz}{Hooke's law}{"hUks lO:}{$F = D \cdot s$}
\eintrag{74}{die}{resultierende Kraft}{resultant force, net force}{rI"z2lt@nt fO:s}{\skResultierende}
\eintrag{75}{das}{Kräftegleichgewicht}{equilibrium}{""i:kwI"lIbri@m}{$\vec{F}_{\mathrm{res}} = \vec{0}$}
\eintrag{76}{die}{Trägheit}{inertia}{I"n3:S@}{}
\eintrag{77}{das}{newtonsche Grundgesetz}{Newton's second law}{"nju:t\textsyllabic{n}z "sek@nd lO:}{$F = m \cdot a$}
\eintrag{78}{das}{Wechselwirkungsgesetz}{Newton's third law (action and reaction)}{"\ae kS\textsyllabic{n} @nd ri"\ae kS\textsyllabic{n}}{$\vec{F}_{AB} = -\vec{F}_{BA}$}
\eintrag{79}{der}{Impuls}{momentum}{m@"ment@m}{$p$ in $\mathrm{kg \cdot m/s}$ \newline $p = m \cdot v$}
\eintrag{80}{der}{Kraftstoß}{impulse}{"Imp2ls}{$F \cdot \Delta t = \Delta p$}
\eintrag{81}{das}{Drehmoment}{torque, moment (of a force)}{tO:k, "m@Um@nt}{$M$ in $\mathrm{N \cdot m}$ \newline $M = F \cdot r$}
\eintrag{82}{der}{Hebel}{lever}{"li:v@}{$F_1 \cdot r_1 = F_2 \cdot r_2$ \newline \skHebel}
\eintrag{83}{der}{Hebelarm}{lever arm}{"li:v@r A:m}{\skHebelarm}
\eintrag{84}{der}{Drehpunkt}{pivot, fulcrum}{"pIv@t, "fUlkr@m}{\skDrehpunkt}
\eintrag{85}{die}{geneigte Ebene}{inclined plane}{In"klaInd pleIn}{\skEbene}
\eintrag{86}{die}{Rolle}{pulley}{"pUli}{\skRolle}
\eintrag{87}{der}{Flaschenzug}{block and tackle}{"bl6k @n "t\ae k\textsyllabic{l}}{}
\eintrag{88}{der}{Schwerpunkt}{centre of gravity}{"sent@r @v "gr\ae v@ti}{}
\eintrag{89}{der}{Druck}{pressure}{"preS@}{$p$ in $\mathrm{Pa}$ \newline $p = \dfrac{F}{A}$ \newline \skDruck}
\eintrag{90}{der}{Schweredruck}{hydrostatic pressure}{""haIdr@"st\ae tIk "preS@}{$p = \varrho \cdot g \cdot h$}
\eintrag{91}{der}{Luftdruck}{atmospheric pressure}{""\ae tm@s"ferIk "preS@}{$p_0 \approx 1013\eh{hPa}$}
\eintrag{92}{der}{Auftrieb}{buoyancy, upthrust}{"bOI@nsi, "2pTr2st}{$F_A = \varrho \cdot g \cdot V$ \newline \skAuftrieb}
\eintrag{93}{die}{Gravitation}{gravitation, gravity}{""gr\ae vI"teIS\textsyllabic{n}, "gr\ae v@ti}{$F = G \cdot \dfrac{m_1 \cdot m_2}{r^2}$}
\eintrag{94}{die}{Radialkraft}{centripetal force}{sen"trIpIt\textsyllabic{l} fO:s}{$F = \dfrac{m \cdot v^2}{r}$}
\end{woerter}

\platz\zwischen{Beispiele}{Examples}{Ig"zA:mp\textsyllabic{l}z}

\begin{lesen}
\lies{$m = 5\eh{kg}$ \newline $F_G = m \cdot g \approx 49\eh{N}$}{The weight of a five-kilogram mass is about forty-nine newtons.}{Die Gewichtskraft eines Körpers der Masse fünf Kilogramm beträgt etwa neunundvierzig Newton.}
\lies{$F = m \cdot a$ \newline $F = 1200\eh{kg} \cdot 2\eh{m/s^2} = 2400\eh{N}$}{A force of two thousand four hundred newtons is needed to accelerate the car at two metres per second squared.}{Um das Auto mit zwei Meter pro Sekunde zum Quadrat zu beschleunigen, ist eine Kraft von zweitausendvierhundert Newton nötig.}
\lies{$D = \dfrac{F}{s} = \dfrac{2\eh{N}}{0{,}05\eh{m}} = 40\eh{N/m}$}{A force of two newtons stretches the spring by five centimetres. The spring constant is forty newtons per metre.}{Eine Kraft von zwei Newton dehnt die Feder um fünf Zentimeter. Die Federkonstante beträgt vierzig Newton pro Meter.}
\lies{\skHebel}{The lever is balanced when the clockwise moment equals the anticlockwise moment.}{Der Hebel ist im Gleichgewicht, wenn rechtsdrehendes und linksdrehendes Drehmoment gleich groß sind.}
\lies{$p = \dfrac{F}{A} = \dfrac{600\eh{N}}{0{,}03\eh{m^2}}$ \newline $p = 20\,000\eh{Pa}$}{The pressure is the force divided by the area: twenty thousand pascals.}{Der Druck ist der Quotient aus Kraft und Fläche: zwanzigtausend Pascal.}
\end{lesen}


%%%%  ENERGIE, ARBEIT UND LEISTUNG  %%%%
\newpage
\kapitel{Energie, Arbeit und Leistung}{Energy, work and power}{"en@dZi, "w3:k @n "paU@}

\begin{woerter}
\eintrag{95}{die}{Energie}{energy}{"en@dZi}{$E$ in $\mathrm{J}$}
\eintrag{96}{die}{Arbeit}{work}{w3:k}{$W$ in $\mathrm{J}$ \newline $W = F \cdot s$}
\eintrag{97}{die}{potentielle Energie}{(gravitational) potential energy}{p@"tenS\textsyllabic{l} "en@dZi}{$E_{\mathrm{pot}} = m \cdot g \cdot h$}
\eintrag{98}{die}{kinetische Energie}{kinetic energy}{kI"netIk "en@dZi}{$E_{\mathrm{kin}} = \dfrac{1}{2} \cdot m \cdot v^2$}
\eintrag{99}{die}{Spannenergie}{elastic potential energy}{I"l\ae stIk p@"tenS\textsyllabic{l} "en@dZi}{$E_{\mathrm{sp}} = \dfrac{1}{2} \cdot D \cdot s^2$}
\eintrag{100}{die}{thermische Energie}{thermal energy}{"T3:m\textsyllabic{l} "en@dZi}{}
\eintrag{101}{die}{elektrische Energie}{electrical energy}{I"lektrIk\textsyllabic{l} "en@dZi}{$E = U \cdot I \cdot t$}
\eintrag{102}{die}{chemische Energie}{chemical energy}{"kemIk\textsyllabic{l} "en@dZi}{}
\eintrag{103}{die}{Energieumwandlung}{energy conversion, energy transfer}{"en@dZi k@n"v3:S\textsyllabic{n}}{}
\eintrag{104}{der}{Energieerhaltungssatz}{law of conservation of energy}{""k6ns@"veIS\textsyllabic{n} @v "en@dZi}{$E_{\mathrm{ges}} = \text{konstant}$}
\eintrag{105}{die}{Leistung}{power}{"paU@}{$P$ in $\mathrm{W}$ \newline $P = \dfrac{W}{t}$}
\eintrag{106}{der}{Wirkungsgrad}{efficiency}{I"fIS\textsyllabic{n}si}{$\eta = \dfrac{E_{\mathrm{nutz}}}{E_{\mathrm{zu}}}$}
\eintrag{107}{die}{Nutzenergie}{useful energy}{"ju:sf\textsyllabic{l} "en@dZi}{}
\eintrag{108}{die}{Abwärme}{waste heat}{"weIst hi:t}{}
\eintrag{109}{der}{Energieträger}{energy source}{"en@dZi sO:s}{}
\eintrag{110}{die}{erneuerbare Energie}{renewable energy}{rI"nju:@b\textsyllabic{l} "en@dZi}{}
\eintrag{111}{das}{Kraftwerk}{power station}{"paU@ ""steIS\textsyllabic{n}}{}
\eintrag{112}{die}{Kilowattstunde}{kilowatt-hour}{"kIl@w6t ""aU@}{$1\eh{kWh} = 3{,}6\eh{MJ}$}
\end{woerter}

\platz\zwischen{Beispiele}{Examples}{Ig"zA:mp\textsyllabic{l}z}

\begin{lesen}
\lies{$W = F \cdot s$ \newline $W = 50\eh{N} \cdot 4\eh{m} = 200\eh{J}$}{The work done is the force times the distance moved in the direction of the force: two hundred joules.}{Die verrichtete Arbeit ist Kraft mal Weg in Kraftrichtung: zweihundert Joule.}
\lies{$E_{\mathrm{pot}} = 2\eh{kg} \cdot 9{,}81\eh{m/s^2} \cdot 10\eh{m}$ \newline $E_{\mathrm{pot}} \approx 196\eh{J}$}{A two-kilogram mass is lifted ten metres. Its potential energy increases by about one hundred and ninety-six joules.}{Ein Körper der Masse zwei Kilogramm wird zehn Meter angehoben. Seine potentielle Energie nimmt um etwa hundertsechsundneunzig Joule zu.}
\lies{$E_{\mathrm{pot}} \rightarrow E_{\mathrm{kin}}$}{As the ball falls, potential energy is converted into kinetic energy.}{Beim Fallen des Balls wird potentielle Energie in kinetische Energie umgewandelt.}
\lies{$P = \dfrac{W}{t} = \dfrac{3000\eh{J}}{60\eh{s}} = 50\eh{W}$}{The power is the work done per unit time: fifty watts.}{Die Leistung ist die pro Zeit verrichtete Arbeit: fünfzig Watt.}
\lies{$\eta = \dfrac{30\eh{J}}{100\eh{J}} = 0{,}3 = 30\,\%$}{The efficiency of the engine is thirty per cent; seventy per cent of the energy is transferred as waste heat.}{Der Wirkungsgrad des Motors beträgt dreißig Prozent; siebzig Prozent der Energie werden als Abwärme abgegeben.}
\end{lesen}


%%%%  WÄRMELEHRE  %%%%
\newpage
\kapitel{Wärmelehre}{Thermal physics}{"T3:m\textsyllabic{l} "fIzIks}

\begin{woerter}
\eintrag{113}{die}{Temperatur}{temperature}{"tempr@tS@}{$\vartheta$ in $\mathrm{\degree C}$}
\eintrag{114}{die}{absolute Temperatur}{absolute temperature, thermodynamic temperature}{"\ae bs@lu:t "tempr@tS@}{$T$ in $\mathrm{K}$}
\eintrag{115}{der}{absolute Nullpunkt}{absolute zero}{"\ae bs@lu:t "zI@r@U}{$0\eh{K} = -273{,}15\eh{\degree C}$}
\eintrag{116}{das}{Thermometer}{thermometer}{T@"m6mIt@}{\skThermometer}
\eintrag{117}{die}{Wärme}{heat}{hi:t}{$Q$ in $\mathrm{J}$}
\eintrag{118}{die}{innere Energie}{internal energy}{In"t3:n\textsyllabic{l} "en@dZi}{$U$ in $\mathrm{J}$}
\eintrag{119}{die}{spezifische Wärmekapazität}{specific heat capacity}{sp@"sIfIk "hi:t k@""p\ae s@ti}{$c$ in $\mathrm{J/(kg \cdot K)}$ \newline $Q = c \cdot m \cdot \Delta T$}
\eintrag{120}{die}{Wärmeleitung}{conduction}{k@n"d2kS\textsyllabic{n}}{}
\eintrag{121}{die}{Wärmeströmung, die Konvektion}{convection}{k@n"vekS\textsyllabic{n}}{}
\eintrag{122}{die}{Wärmestrahlung}{thermal radiation}{"T3:m\textsyllabic{l} ""reIdi"eIS\textsyllabic{n}}{}
\eintrag{123}{die}{Wärmedämmung}{thermal insulation}{"T3:m\textsyllabic{l} ""Insju"leIS\textsyllabic{n}}{}
\eintrag{124}{das}{Teilchenmodell}{particle model}{"pA:tIk\textsyllabic{l} ""m6d\textsyllabic{l}}{}
\eintrag{125}{der}{Aggregatzustand}{state of matter}{"steIt @v "m\ae t@}{}
\eintrag{126}{}{fest}{solid}{"s6lId}{\skFest}
\eintrag{127}{}{flüssig}{liquid}{"lIkwId}{\skFluessig}
\eintrag{128}{}{gasförmig}{gaseous}{"g\ae si@s}{\skGas}
\eintrag{129}{}{schmelzen}{melt}{melt}{fest $\rightarrow$ flüssig}
\eintrag{130}{}{erstarren}{freeze, solidify}{fri:z, s@"lIdIfaI}{flüssig $\rightarrow$ fest}
\eintrag{131}{}{verdampfen}{evaporate, vaporise}{I"v\ae p@reIt, "veIp@raIz}{flüssig $\rightarrow$ gasförmig}
\eintrag{132}{}{sieden}{boil}{bOIl}{}
\eintrag{133}{}{kondensieren}{condense}{k@n"dens}{gasförmig $\rightarrow$ flüssig}
\eintrag{134}{}{sublimieren}{sublime}{s@"blaIm}{fest $\rightarrow$ gasförmig}
\eintrag{135}{die}{Schmelztemperatur}{melting point}{"meltIN pOInt}{Eis: $0\eh{\degree C}$}
\eintrag{136}{die}{Siedetemperatur}{boiling point}{"bOIlIN pOInt}{Wasser: $100\eh{\degree C}$}
\eintrag{137}{die}{Schmelzwärme}{latent heat of fusion}{"leIt\textsyllabic{n}t "hi:t @v "fju:Z\textsyllabic{n}}{$Q = q_s \cdot m$}
\eintrag{138}{die}{Verdampfungswärme}{latent heat of vaporisation}{"leIt\textsyllabic{n}t "hi:t @v ""veIp@raI"zeIS\textsyllabic{n}}{$Q = q_v \cdot m$}
\eintrag{139}{die}{Längenausdehnung}{linear expansion}{"lIni@r Ik"sp\ae nS\textsyllabic{n}}{$\Delta l = \alpha \cdot l_0 \cdot \Delta T$}
\eintrag{140}{das}{ideale Gas}{ideal gas}{aI"dI@l g\ae s}{$\dfrac{p \cdot V}{T} = \text{konstant}$}
\eintrag{141}{die}{Wärmekraftmaschine}{heat engine}{"hi:t ""endZIn}{}
\eintrag{142}{die}{Entropie}{entropy}{"entr@pi}{$S$ in $\mathrm{J/K}$}
\end{woerter}

\platz\zwischen{Beispiele}{Examples}{Ig"zA:mp\textsyllabic{l}z}

\begin{lesen}
\lies{$\vartheta = 20\eh{\degree C}$ \newline $T = 293{,}15\eh{K}$}{Twenty degrees Celsius corresponds to two hundred and ninety-three point one five kelvin.}{Zwanzig Grad Celsius entsprechen zweihundertdreiundneunzig Komma eins fünf Kelvin.}
\lies{$Q = c \cdot m \cdot \Delta T$ \newline $Q = 4{,}19\eh{\tfrac{kJ}{kg \cdot K}} \cdot 2\eh{kg} \cdot 30\eh{K}$ \newline $Q \approx 251\eh{kJ}$}{About two hundred and fifty-one kilojoules of heat are needed to raise the temperature of two kilograms of water by thirty kelvin.}{Um zwei Kilogramm Wasser um dreißig Kelvin zu erwärmen, sind etwa zweihunderteinundfünfzig Kilojoule Wärme nötig.}
\lies{\skFest \; \skFluessig \; \skGas}{In a solid the particles vibrate about fixed positions; in a liquid they can move past each other; in a gas they move freely.}{In einem festen Körper schwingen die Teilchen um feste Plätze; in einer Flüssigkeit sind sie gegeneinander verschiebbar; in einem Gas bewegen sie sich frei.}
\lies{$100\eh{\degree C}$}{Water boils at one hundred degrees Celsius. While it is boiling, its temperature stays constant.}{Wasser siedet bei hundert Grad Celsius. Während des Siedens bleibt die Temperatur konstant.}
\lies{Wärmeleitung}{Metals are good conductors of heat; air is a good insulator.}{Metalle sind gute Wärmeleiter; Luft ist ein guter Isolator.}
\end{lesen}


%%%%  ELEKTRIZITÄT UND MAGNETISMUS  %%%%
\newpage
\kapitel{Elektrizität und Magnetismus}{Electricity and magnetism}{I""lek"trIs@ti @n "m\ae gn@tIz@m}

\begin{woerter}
\eintrag{143}{die}{elektrische Ladung}{electric charge}{I"lektrIk tSA:dZ}{$Q$ in $\mathrm{C}$}
\eintrag{144}{die}{elektrische Stromstärke}{electric current}{I"lektrIk "k2r@nt}{$I$ in $\mathrm{A}$ \newline $I = \dfrac{Q}{t}$}
\eintrag{145}{die}{elektrische Spannung}{voltage, potential difference}{"v@UltIdZ, p@"tenS\textsyllabic{l} "dIfr@ns}{$U$ in $\mathrm{V}$ \newline \footnotesize engl. $V$}
\eintrag{146}{der}{elektrische Widerstand}{resistance}{rI"zIst@ns}{$R$ in $\mathrm{\Omega}$ \newline $R = \dfrac{U}{I}$}
\eintrag{147}{das}{ohmsche Gesetz}{Ohm's law}{"@Umz lO:}{$U = R \cdot I$}
\eintrag{148}{der}{spezifische Widerstand}{resistivity}{""ri:zIs"tIv@ti}{$R = \varrho \cdot \dfrac{l}{A}$}
\eintrag{149}{der}{Leitwert}{conductance}{k@n"d2kt@ns}{$G$ in $\mathrm{S}$}
\eintrag{150}{der}{Leiter}{conductor}{k@n"d2kt@}{}
\eintrag{151}{der}{Isolator, der Nichtleiter}{insulator}{"InsjuleIt@}{}
\eintrag{152}{der}{Halbleiter}{semiconductor}{""semik@n"d2kt@}{}
\eintrag{153}{der}{Stromkreis}{electric circuit}{I"lektrIk "s3:kIt}{\skStromkreis}
\eintrag{154}{der}{Schaltplan}{circuit diagram}{"s3:kIt ""daI@gr\ae m}{}
\eintrag{155}{die}{Spannungsquelle}{voltage source, power supply}{"v@UltIdZ sO:s, "paU@ s@""plaI}{\sz{battery1}}
\eintrag{156}{die}{Batterie}{battery, cell}{"b\ae tri, sel}{\sz{battery1}}
\eintrag{157}{der}{Schalter}{switch}{swItS}{\sz{nos}}
\eintrag{158}{die}{Glühlampe}{lamp, bulb}{l\ae mp, b2lb}{\sz{lamp}}
\eintrag{159}{der}{Widerstand (Bauteil)}{resistor}{rI"zIst@}{\sz{R}}
\eintrag{160}{der}{Strommesser}{ammeter}{"\ae mi:t@}{\sz{rmeter, t=A}}
\eintrag{161}{der}{Spannungsmesser}{voltmeter}{"v@Ult""mi:t@}{\sz{rmeter, t=V}}
\eintrag{162}{die}{Reihenschaltung}{series connection}{"sI@ri:z k@"nekS\textsyllabic{n}}{$R = R_1 + R_2$ \newline \skReihe}
\eintrag{163}{die}{Parallelschaltung}{parallel connection}{"p\ae r@lel k@"nekS\textsyllabic{n}}{$\dfrac{1}{R} = \dfrac{1}{R_1} + \dfrac{1}{R_2}$ \newline \skParallel}
\eintrag{164}{der}{Kurzschluss}{short circuit}{""SO:t "s3:kIt}{}
\eintrag{165}{die}{Sicherung}{fuse}{fju:z}{}
\eintrag{166}{der}{Gleichstrom}{direct current (DC)}{d@"rekt "k2r@nt}{}
\eintrag{167}{der}{Wechselstrom}{alternating current (AC)}{"O:lt@neItIN "k2r@nt}{}
\eintrag{168}{die}{elektrische Leistung}{electrical power}{I"lektrIk\textsyllabic{l} "paU@}{$P = U \cdot I$}
\eintrag{169}{der}{Kondensator}{capacitor}{k@"p\ae sIt@}{\sz{C}}
\eintrag{170}{die}{Kapazität}{capacitance}{k@"p\ae sIt@ns}{$C$ in $\mathrm{F}$ \newline $C = \dfrac{Q}{U}$}
\eintrag{171}{das}{elektrische Feld}{electric field}{I"lektrIk fi:ld}{\skEFeld}
\eintrag{172}{die}{elektrische Feldstärke}{electric field strength}{I"lektrIk "fi:ld streNT}{$E$ in $\mathrm{V/m}$ \newline $E = \dfrac{F}{Q}$}
\eintrag{173}{die}{Feldlinie}{field line}{"fi:ld laIn}{}
\eintrag{174}{der}{Magnet}{magnet}{"m\ae gn@t}{\skMagnet}
\eintrag{175}{der}{Dauermagnet}{permanent magnet}{"p3:m@n@nt "m\ae gn@t}{}
\eintrag{176}{der}{Elektromagnet}{electromagnet}{I"lektr@U""m\ae gn@t}{}
\eintrag{177}{der}{Nordpol, der Südpol}{north pole, south pole}{"nO:T p@Ul, "saUT p@Ul}{N, S}
\eintrag{178}{das}{Magnetfeld}{magnetic field}{m\ae g"netIk fi:ld}{\skMagnetfeld}
\eintrag{179}{die}{magnetische Flussdichte}{magnetic flux density}{m\ae g"netIk "fl2ks ""dens@ti}{$B$ in $\mathrm{T}$}
\eintrag{180}{der}{magnetische Fluss}{magnetic flux}{m\ae g"netIk fl2ks}{$\Phi$ in $\mathrm{Wb}$ \newline $\Phi = B \cdot A$}
\eintrag{181}{die}{Spule}{coil}{kOIl}{\sz{cute inductor}}
\eintrag{182}{die}{Windungszahl}{number of turns}{"n2mb@r @v "t3:nz}{$N$}
\eintrag{183}{die}{Induktivität}{inductance}{In"d2kt@ns}{$L$ in $\mathrm{H}$}
\eintrag{184}{die}{elektromagnetische Induktion}{electromagnetic induction}{I""lektr@Um\ae g"netIk In"d2kS\textsyllabic{n}}{$U_{\mathrm{ind}} = -N \cdot \dfrac{\Delta\Phi}{\Delta t}$}
\eintrag{185}{die}{Lorentzkraft}{Lorentz force}{"l6r@nts fO:s}{$F = Q \cdot v \cdot B$}
\eintrag{186}{der}{Elektromotor}{electric motor}{I"lektrIk "m@Ut@}{}
\eintrag{187}{der}{Generator}{generator}{"dZen@reIt@}{}
\eintrag{188}{der}{Transformator}{transformer}{tr\ae ns"fO:m@}{$\dfrac{U_1}{U_2} = \dfrac{N_1}{N_2}$}
\eintrag{189}{die}{Erdung}{earth (BE), ground (AE)}{3:T, graUnd}{}
\end{woerter}

\platz\zwischen{Beispiele}{Examples}{Ig"zA:mp\textsyllabic{l}z}

\begin{lesen}
\lies{$R = \dfrac{U}{I} = \dfrac{12\eh{V}}{0{,}5\eh{A}} = 24\eh{\Omega}$}{The voltage across the resistor is twelve volts and the current through it is half an ampere. Its resistance is twenty-four ohms.}{Am Widerstand liegt eine Spannung von zwölf Volt an, die Stromstärke beträgt ein halbes Ampere. Der Widerstand beträgt vierundzwanzig Ohm.}
\lies{\skReihe}{The two resistors are connected in series: the current through them is the same.}{Die beiden Widerstände sind in Reihe geschaltet: Die Stromstärke ist überall gleich.}
\lies{\skParallel}{The two resistors are connected in parallel: the voltage across them is the same.}{Die beiden Widerstände sind parallel geschaltet: An beiden liegt die gleiche Spannung an.}
\lies{\sz{rmeter, t=A} \; \sz{rmeter, t=V}}{An ammeter is connected in series, a voltmeter in parallel with the component.}{Ein Strommesser wird in Reihe, ein Spannungsmesser parallel zum Bauteil geschaltet.}
\lies{\skMagnet}{Like poles repel each other, unlike poles attract each other.}{Gleichnamige Pole stoßen einander ab, ungleichnamige Pole ziehen einander an.}
\end{lesen}


%%%%  SCHWINGUNGEN UND WELLEN  %%%%
\newpage
\kapitel{Schwingungen und Wellen}{Oscillations and waves}{""6sI"leIS\textsyllabic{n}z @n "weIvz}

\begin{woerter}
\eintrag{190}{die}{Schwingung}{oscillation, vibration}{""6sI"leIS\textsyllabic{n}, vaI"breIS\textsyllabic{n}}{}
\eintrag{191}{das}{Fadenpendel}{simple pendulum}{"sImp\textsyllabic{l} "pendj@l@m}{$T = 2\pi \cdot \sqrt{\dfrac{l}{g}}$ \newline \skPendel}
\eintrag{192}{der}{Federschwinger}{mass on a spring}{"m\ae s 6n @ "sprIN}{$T = 2\pi \cdot \sqrt{\dfrac{m}{D}}$}
\eintrag{193}{die}{Ruhelage}{equilibrium position}{""i:kwI"lIbri@m p@"zIS\textsyllabic{n}}{}
\eintrag{194}{die}{Auslenkung}{displacement}{dIs"pleIsm@nt}{$y$ in $\mathrm{m}$}
\eintrag{195}{die}{Amplitude}{amplitude}{"\ae mplItju:d}{$\hat{y}$ \newline \skAmplitude}
\eintrag{196}{die}{Periodendauer}{period}{"pI@ri@d}{$T$ in $\mathrm{s}$ \newline \skPeriode}
\eintrag{197}{die}{Frequenz}{frequency}{"fri:kw@nsi}{$f$ in $\mathrm{Hz}$ \newline $f = \dfrac{1}{T}$}
\eintrag{198}{die}{gedämpfte Schwingung}{damped oscillation}{"d\ae mpt ""6sI"leIS\textsyllabic{n}}{\skGedaempft}
\eintrag{199}{die}{Eigenfrequenz}{natural frequency}{"n\ae tSr@l "fri:kw@nsi}{$f_0$}
\eintrag{200}{die}{Resonanz}{resonance}{"rez@n@ns}{}
\eintrag{201}{die}{Welle}{wave}{weIv}{}
\eintrag{202}{die}{Wellenlänge}{wavelength}{"weIvleNT}{$\lambda$ in $\mathrm{m}$ \newline \skWellenlaenge}
\eintrag{203}{die}{Ausbreitungsgeschwindigkeit}{wave speed}{"weIv spi:d}{$c = \lambda \cdot f$}
\eintrag{204}{der}{Wellenberg, das Wellental}{crest, trough}{krest, tr6f}{\skBergTal}
\eintrag{205}{die}{Querwelle (Transversalwelle)}{transverse wave}{"tr\ae nzv3:s weIv}{}
\eintrag{206}{die}{Längswelle (Longitudinalwelle)}{longitudinal wave}{""l6ndZI"tju:dIn\textsyllabic{l} weIv}{\skLaengswelle}
\eintrag{207}{die}{Reflexion}{reflection}{rI"flekS\textsyllabic{n}}{}
\eintrag{208}{die}{Beugung}{diffraction}{dI"fr\ae kS\textsyllabic{n}}{}
\eintrag{209}{die}{Interferenz}{interference}{""Int@"fI@r@ns}{}
\eintrag{210}{die}{Überlagerung}{superposition}{""su:p@p@"zIS\textsyllabic{n}}{}
\eintrag{211}{die}{stehende Welle}{standing wave, stationary wave}{"st\ae ndIN weIv}{\skStehend}
\eintrag{212}{der}{Schall}{sound}{saUnd}{}
\eintrag{213}{die}{Schallgeschwindigkeit}{speed of sound}{"spi:d @v "saUnd}{Luft: $c \approx 340\eh{m/s}$}
\eintrag{214}{die}{Tonhöhe}{pitch}{pItS}{}
\eintrag{215}{die}{Lautstärke}{loudness, volume}{"laUdn@s, "v6lju:m}{}
\eintrag{216}{der}{Schallpegel}{sound level}{"saUnd ""lev\textsyllabic{l}}{$L$ in $\mathrm{dB}$}
\eintrag{217}{das}{Echo}{echo}{"ek@U}{}
\eintrag{218}{der}{Ultraschall}{ultrasound}{"2ltr@saUnd}{$f > 20\eh{kHz}$}
\eintrag{219}{der}{Doppler-Effekt}{Doppler effect}{"d6pl@r I""fekt}{}
\end{woerter}

\platz\zwischen{Beispiele}{Examples}{Ig"zA:mp\textsyllabic{l}z}

\begin{lesen}
\lies{$l = 1\eh{m}$ \newline $T = 2\pi \cdot \sqrt{\dfrac{l}{g}} \approx 2{,}0\eh{s}$}{A simple pendulum one metre long has a period of about two seconds.}{Ein Fadenpendel von einem Meter Länge hat eine Periodendauer von etwa zwei Sekunden.}
\lies{$f = 440\eh{Hz}$ \newline $\lambda = \dfrac{c}{f} = \dfrac{340\eh{m/s}}{440\eh{Hz}} \approx 0{,}77\eh{m}$}{The note A has a frequency of four hundred and forty hertz. In air its wavelength is about seventy-seven centimetres.}{Der Ton a hat eine Frequenz von vierhundertvierzig Hertz. In Luft beträgt seine Wellenlänge etwa siebenundsiebzig Zentimeter.}
\lies{\skAmplitude}{The larger the amplitude, the louder the sound; the higher the frequency, the higher the pitch.}{Je größer die Amplitude, desto lauter der Ton; je höher die Frequenz, desto höher der Ton.}
\lies{\skLaengswelle}{Sound is a longitudinal wave: the particles oscillate parallel to the direction of travel.}{Schall ist eine Längswelle: Die Teilchen schwingen parallel zur Ausbreitungsrichtung.}
\lies{$s = \dfrac{340\eh{m/s} \cdot 2\eh{s}}{2} = 340\eh{m}$}{The echo returns after two seconds, so the wall is three hundred and forty metres away.}{Das Echo kommt nach zwei Sekunden zurück, die Wand ist also dreihundertvierzig Meter entfernt.}
\end{lesen}


%%%%  OPTIK  %%%%
\newpage
\kapitel{Optik}{Optics}{"6ptIks}

\begin{woerter}
\eintrag{220}{das}{Licht}{light}{laIt}{}
\eintrag{221}{die}{Lichtquelle}{light source}{"laIt sO:s}{}
\eintrag{222}{der}{Lichtstrahl}{ray of light, light ray}{"reI @v "laIt}{}
\eintrag{223}{das}{Lichtbündel}{beam of light}{"bi:m @v "laIt}{}
\eintrag{224}{die}{Lichtgeschwindigkeit}{speed of light}{"spi:d @v "laIt}{$c \approx 3 \cdot 10^8\eh{m/s}$}
\eintrag{225}{der}{Schatten}{shadow}{"S\ae d@U}{\skSchatten}
\eintrag{226}{der}{Kernschatten, der Halbschatten}{umbra, penumbra}{"2mbr@, pI"n2mbr@}{}
\eintrag{227}{die}{Reflexion}{reflection}{rI"flekS\textsyllabic{n}}{\skReflexion}
\eintrag{228}{das}{Reflexionsgesetz}{law of reflection}{"lO: @v rI"flekS\textsyllabic{n}}{$\alpha = \alpha'$}
\eintrag{229}{das}{Einfallslot}{normal}{"nO:m\textsyllabic{l}}{\skLot}
\eintrag{230}{der}{Einfallswinkel}{angle of incidence}{"\ae Ng\textsyllabic{l} @v "InsId@ns}{$\alpha$}
\eintrag{231}{der}{Reflexionswinkel}{angle of reflection}{"\ae Ng\textsyllabic{l} @v rI"flekS\textsyllabic{n}}{$\alpha'$}
\eintrag{232}{der}{Spiegel}{mirror}{"mIr@}{}
\eintrag{233}{der}{Hohlspiegel, der Wölbspiegel}{concave mirror, convex mirror}{"k6nkeIv "mIr@, "k6nveks "mIr@}{}
\eintrag{234}{die}{Brechung}{refraction}{rI"fr\ae kS\textsyllabic{n}}{\skBrechung}
\eintrag{235}{der}{Brechungswinkel}{angle of refraction}{"\ae Ng\textsyllabic{l} @v rI"fr\ae kS\textsyllabic{n}}{$\beta$}
\eintrag{236}{die}{Brechzahl}{refractive index}{rI"fr\ae ktIv "Indeks}{$n = \dfrac{c_0}{c}$}
\eintrag{237}{das}{Brechungsgesetz}{Snell's law, law of refraction}{"snelz lO:}{$\dfrac{\sin\alpha}{\sin\beta} = \dfrac{n_2}{n_1}$}
\eintrag{238}{die}{Totalreflexion}{total internal reflection}{"t@Ut\textsyllabic{l} In"t3:n\textsyllabic{l} rI"flekS\textsyllabic{n}}{}
\eintrag{239}{der}{Grenzwinkel}{critical angle}{"krItIk\textsyllabic{l} "\ae Ng\textsyllabic{l}}{$\alpha_G$}
\eintrag{240}{die}{Linse}{lens}{lenz}{}
\eintrag{241}{die}{Sammellinse}{converging lens, convex lens}{k@n"v3:dZIN lenz}{\skSammellinse}
\eintrag{242}{die}{Zerstreuungslinse}{diverging lens, concave lens}{daI"v3:dZIN lenz}{\skZerstreuungslinse}
\eintrag{243}{die}{optische Achse}{principal axis}{"prIns@p\textsyllabic{l} "\ae ksIs}{}
\eintrag{244}{der}{Brennpunkt}{focal point, principal focus}{"f@Uk\textsyllabic{l} pOInt}{$F$}
\eintrag{245}{die}{Brennweite}{focal length}{"f@Uk\textsyllabic{l} leNT}{$f$ in $\mathrm{m}$ \newline \skBrennweite}
\eintrag{246}{die}{Gegenstandsweite}{object distance}{"6bdZIkt ""dIst@ns}{$g$ in $\mathrm{m}$ \newline \footnotesize engl. $u$}
\eintrag{247}{die}{Bildweite}{image distance}{"ImIdZ ""dIst@ns}{$b$ in $\mathrm{m}$ \newline \footnotesize engl. $v$}
\eintrag{248}{die}{Linsengleichung}{lens equation}{"lenz I""kweIZ\textsyllabic{n}}{$\dfrac{1}{f} = \dfrac{1}{g} + \dfrac{1}{b}$}
\eintrag{249}{das}{Bild}{image}{"ImIdZ}{}
\eintrag{250}{das}{reelle Bild}{real image}{"ri:@l "ImIdZ}{}
\eintrag{251}{das}{virtuelle Bild}{virtual image}{"v3:tSu@l "ImIdZ}{}
\eintrag{252}{der}{Abbildungsmaßstab}{magnification}{""m\ae gnIfI"keIS\textsyllabic{n}}{$A = \dfrac{B}{G} = \dfrac{b}{g}$}
\eintrag{253}{die}{Lupe}{magnifying glass}{"m\ae gnIfaIIN glA:s}{}
\eintrag{254}{das}{Prisma}{prism}{"prIz@m}{\skPrisma}
\eintrag{255}{das}{Spektrum}{spectrum}{"spektr@m}{}
\eintrag{256}{die}{Farbzerlegung, die Dispersion}{dispersion}{dI"sp3:S\textsyllabic{n}}{}
\eintrag{257}{die}{Lichtstärke}{luminous intensity}{"lu:mIn@s In"tens@ti}{$I_v$ in $\mathrm{cd}$}
\eintrag{258}{die}{Beleuchtungsstärke}{illuminance}{I"lu:mIn@ns}{$E_v$ in $\mathrm{lx}$}
\end{woerter}

\platz\zwischen{Beispiele}{Examples}{Ig"zA:mp\textsyllabic{l}z}

\begin{lesen}
\lies{\skReflexion}{The angle of incidence equals the angle of reflection. Both angles are measured from the normal.}{Der Einfallswinkel ist gleich dem Reflexionswinkel. Beide Winkel werden zum Lot gemessen.}
\lies{\skBrechung}{When light passes from air into water, it is refracted towards the normal.}{Beim Übergang von Luft in Wasser wird das Licht zum Lot hin gebrochen.}
\lies{\skSammellinse}{Rays parallel to the principal axis are refracted through the focal point.}{Parallelstrahlen werden so gebrochen, dass sie durch den Brennpunkt verlaufen.}
\lies{$f = 10\eh{cm},\ g = 30\eh{cm}$ \newline $\dfrac{1}{b} = \dfrac{1}{f} - \dfrac{1}{g}$ \newline $b = 15\eh{cm}$}{The object is thirty centimetres from the lens. The image is formed fifteen centimetres behind the lens; it is real, inverted and diminished.}{Der Gegenstand steht dreißig Zentimeter vor der Linse. Das Bild entsteht fünfzehn Zentimeter hinter der Linse; es ist reell, umgekehrt und verkleinert.}
\lies{\skPrisma}{A prism splits white light into the colours of the spectrum.}{Ein Prisma zerlegt weißes Licht in die Spektralfarben.}
\end{lesen}


%%%%  ATOM- UND KERNPHYSIK  %%%%
\newpage
\kapitel{Atom- und Kernphysik}{Atomic and nuclear physics}{@"t6mIk @n "nju:kli@ "fIzIks}

\begin{woerter}
\eintrag{259}{das}{Atom}{atom}{"\ae t@m}{\skAtom}
\eintrag{260}{der}{Atomkern}{nucleus}{"nju:kli@s}{\skKern}
\eintrag{261}{die}{Atomhülle}{electron shell}{I"lektr6n Sel}{\skHuelle}
\eintrag{262}{das}{Elektron}{electron}{I"lektr6n}{$\mathrm{e^-}$}
\eintrag{263}{das}{Proton}{proton}{"pr@Ut6n}{$\mathrm{p}$}
\eintrag{264}{das}{Neutron}{neutron}{"nju:tr6n}{$\mathrm{n}$}
\eintrag{265}{das}{Nukleon}{nucleon}{"nju:kli6n}{}
\eintrag{266}{die}{Elementarladung}{elementary charge}{""elI"mentri tSA:dZ}{\small$e = 1{,}602 \cdot 10^{-19}\eh{C}$}
\eintrag{267}{die}{Kernladungszahl, die Ordnungszahl}{atomic number, proton number}{@"t6mIk "n2mb@}{$Z$}
\eintrag{268}{die}{Massenzahl}{mass number, nucleon number}{"m\ae s ""n2mb@}{$A$}
\eintrag{269}{die}{Neutronenzahl}{neutron number}{"nju:tr6n ""n2mb@}{$N = A - Z$}
\eintrag{270}{das}{Nuklid}{nuclide}{"nju:klaId}{$^{A}_{Z}\mathrm{X}$, z.\,B.\ $^{14}_{\ 6}\mathrm{C}$}
\eintrag{271}{das}{Isotop}{isotope}{"aIs@t@Up}{$^{12}_{\ 6}\mathrm{C}$, $^{14}_{\ 6}\mathrm{C}$}
\eintrag{272}{das}{Ion}{ion}{"aI@n}{$\mathrm{Na^+}$}
\eintrag{273}{die}{Radioaktivität}{radioactivity}{""reIdi@U\ae k"tIv@ti}{}
\eintrag{274}{der}{radioaktive Zerfall}{radioactive decay}{""reIdi@U"\ae ktIv dI"keI}{}
\eintrag{275}{die}{Alphastrahlung}{alpha radiation}{"\ae lf@ ""reIdi"eIS\textsyllabic{n}}{$\alpha$: $^{4}_{2}\mathrm{He}$}
\eintrag{276}{die}{Betastrahlung}{beta radiation}{"bi:t@ ""reIdi"eIS\textsyllabic{n}}{$\beta^-$: $\mathrm{e^-}$}
\eintrag{277}{die}{Gammastrahlung}{gamma radiation}{"g\ae m@ ""reIdi"eIS\textsyllabic{n}}{$\gamma$}
\eintrag{278}{die}{ionisierende Strahlung}{ionising radiation}{"aI@naIzIN ""reIdi"eIS\textsyllabic{n}}{}
\eintrag{279}{das}{Zählrohr}{Geiger counter}{"gaIg@ ""kaUnt@}{}
\eintrag{280}{die}{Halbwertszeit}{half-life}{"hA:f laIf}{$T_{1/2}$ \newline \skZerfall}
\eintrag{281}{das}{Zerfallsgesetz}{decay law}{dI"keI lO:}{$N = N_0 \cdot \mathrm{e}^{-\lambda \cdot t}$}
\eintrag{282}{die}{Zerfallskonstante}{decay constant}{dI"keI ""k6nst@nt}{$\lambda = \dfrac{\ln 2}{T_{1/2}}$}
\eintrag{283}{die}{Aktivität}{activity}{\ae k"tIv@ti}{$A$ in $\mathrm{Bq}$}
\eintrag{284}{die}{Äquivalentdosis}{equivalent dose}{I"kwIv@l@nt d@Us}{$H$ in $\mathrm{Sv}$}
\eintrag{285}{die}{Kernspaltung}{nuclear fission}{"nju:kli@ "fIS\textsyllabic{n}}{\skSpaltung}
\eintrag{286}{die}{Kernfusion}{nuclear fusion}{"nju:kli@ "fju:Z\textsyllabic{n}}{}
\eintrag{287}{die}{Kettenreaktion}{chain reaction}{"tSeIn ri""\ae kS\textsyllabic{n}}{}
\eintrag{288}{der}{Kernreaktor}{nuclear reactor}{"nju:kli@ ri"\ae kt@}{}
\eintrag{289}{der}{Massendefekt}{mass defect}{"m\ae s ""di:fekt}{$E = \Delta m \cdot c^2$}
\eintrag{290}{die}{Bindungsenergie}{binding energy}{"baIndIN ""en@dZi}{}
\eintrag{291}{das}{Photon}{photon}{"f@Ut6n}{$E = h \cdot f$}
\eintrag{292}{das}{plancksche Wirkungsquantum}{Planck constant}{"pl\ae Nk ""k6nst@nt}{\small$h = 6{,}626 \cdot 10^{-34}\eh{J \cdot s}$}
\eintrag{293}{der}{Fotoeffekt}{photoelectric effect}{""f@Ut@UI"lektrIk I"fekt}{}
\eintrag{294}{das}{Energieniveau}{energy level}{"en@dZi ""lev\textsyllabic{l}}{\skNiveaus}
\eintrag{295}{die}{Hauptquantenzahl}{principal quantum number}{"prIns@p\textsyllabic{l} "kw6nt@m ""n2mb@}{$n = 1;\ 2;\ 3;\ \ldots$}
\end{woerter}

\platz\zwischen{Beispiele}{Examples}{Ig"zA:mp\textsyllabic{l}z}

\begin{lesen}
\lies{$^{14}_{\ 6}\mathrm{C}$}{carbon-fourteen: the nucleus contains six protons and eight neutrons}{Kohlenstoff-14: Der Kern enthält sechs Protonen und acht Neutronen}
\lies{$^{226}_{\ 88}\mathrm{Ra} \rightarrow {}^{222}_{\ 86}\mathrm{Rn} + {}^{4}_{2}\mathrm{He}$}{Radium-two two six decays into radon-two two two by emitting an alpha particle.}{Radium-226 zerfällt unter Aussendung eines Alphateilchens in Radon-222.}
\lies{$T_{1/2} = 8\eh{d}$}{The half-life of iodine-one three one is eight days: after eight days half of the nuclei have decayed, after sixteen days three quarters.}{Die Halbwertszeit von Iod-131 beträgt acht Tage: Nach acht Tagen ist die Hälfte der Kerne zerfallen, nach sechzehn Tagen drei Viertel.}
\lies{$\alpha,\ \beta,\ \gamma$}{Alpha radiation is stopped by a sheet of paper, beta radiation by a few millimetres of aluminium; gamma radiation is only reduced by thick lead.}{Alphastrahlung wird durch ein Blatt Papier abgeschirmt, Betastrahlung durch einige Millimeter Aluminium; Gammastrahlung wird erst durch dickes Blei geschwächt.}
\lies{$E = h \cdot f$}{The energy of a photon equals the Planck constant times the frequency.}{Die Energie eines Photons ist gleich dem planckschen Wirkungsquantum mal der Frequenz.}
\end{lesen}



%%%%  TEIL II: PHYSIK LESEN UND SPRECHEN  %%%%
\newpage
\teil{Physik lesen und sprechen}{Reading and speaking physics}{"ri:dIN @n "spi:kIN "fIzIks}

In diesem Teil steht, wie man Einheiten, Größenangaben und Formeln auf Englisch (und zum Vergleich
auf Deutsch) vorliest und wie man Experimente und Diagramme beschreibt. Links steht der Ausdruck,
rechts die englische Lesart und darunter grau die deutsche.

%--------------------------------------------------------------------------------------------------
\kapitel{Die Basiseinheiten des SI}{The SI base units}{Di ""es "aI "beIs ""ju:nIts}

Das Internationale Einheitensystem (\textit{International System of Units}, SI) hat sieben
Basiseinheiten. Alle anderen Einheiten werden daraus abgeleitet.

\begin{lesen}
\lies{$\mathrm{m}$}{metre \ipa{"mi:t@} (AE: meter) -- the unit of length}{Meter -- die Einheit der Länge}
\lies{$\mathrm{kg}$}{kilogram \ipa{"kIl@gr\ae m} -- the unit of mass}{Kilogramm -- die Einheit der Masse}
\lies{$\mathrm{s}$}{second \ipa{"sek@nd} -- the unit of time}{Sekunde -- die Einheit der Zeit}
\lies{$\mathrm{A}$}{ampere \ipa{"\ae mpe@} (short: amp) -- the unit of electric current}{Ampere -- die Einheit der elektrischen Stromstärke}
\lies{$\mathrm{K}$}{kelvin \ipa{"kelvIn} -- the unit of thermodynamic temperature}{Kelvin -- die Einheit der absoluten Temperatur}
\lies{$\mathrm{mol}$}{mole \ipa{m@Ul} -- the unit of amount of substance}{Mol -- die Einheit der Stoffmenge}
\lies{$\mathrm{cd}$}{candela \ipa{k\ae n"del@} -- the unit of luminous intensity}{Candela -- die Einheit der Lichtstärke}
\end{lesen}

\anmerkung{Einheitennamen schreibt man im Englischen klein, auch wenn sie nach Personen benannt sind
(\textit{newton, joule, watt, kelvin}); nur das Einheitenzeichen ist groß (N, J, W, K). Es heißt
\textit{kelvin}, nicht \textit{degree Kelvin} -- das „Grad“ wurde 1968 abgeschafft.}

\platz\zwischen{Beispiele}{Examples}{Ig"zA:mp\textsyllabic{l}z}

\begin{lesen}
\lies{$1{,}80\eh{m}$}{one point eight metres \;/\; one metre eighty}{ein Meter achtzig}
\lies{$75\eh{kg}$}{seventy-five kilograms (short: kilos)}{fünfundsiebzig Kilogramm}
\lies{$0{,}5\eh{A}$}{nought point five amperes \;/\; half an amp}{null Komma fünf Ampere}
\end{lesen}

%--------------------------------------------------------------------------------------------------
\newpage
\kapitel{Abgeleitete Einheiten mit eigenem Namen}{Derived units with special names}{dI"raIvd "ju:nIts wID "speS\textsyllabic{l} "neImz}

\begin{lesen}
\lies{$1\eh{Hz} = 1\eh{s^{-1}}$}{hertz \ipa{h3:ts} -- the unit of frequency: one hertz is one per second}{Hertz -- die Einheit der Frequenz: ein Hertz ist eins durch Sekunde}
\lies{$1\eh{N} = 1\eh{kg \cdot m/s^2}$}{newton \ipa{"nju:t\textsyllabic{n}} -- the unit of force: one newton is one kilogram metre per second squared}{Newton -- die Einheit der Kraft: ein Newton ist ein Kilogramm mal Meter durch Sekunde zum Quadrat}
\lies{$1\eh{Pa} = 1\eh{N/m^2}$}{pascal \ipa{p\ae s"k\ae l} -- the unit of pressure: one pascal is one newton per square metre}{Pascal -- die Einheit des Drucks: ein Pascal ist ein Newton pro Quadratmeter}
\lies{$1\eh{J} = 1\eh{N \cdot m}$}{joule \ipa{dZu:l} -- the unit of energy and work: one joule is one newton metre}{Joule -- die Einheit von Energie und Arbeit: ein Joule ist ein Newtonmeter}
\lies{$1\eh{W} = 1\eh{J/s}$}{watt \ipa{w6t} -- the unit of power: one watt is one joule per second}{Watt -- die Einheit der Leistung: ein Watt ist ein Joule pro Sekunde}
\lies{$1\eh{C} = 1\eh{A \cdot s}$}{coulomb \ipa{"ku:l6m} -- the unit of electric charge: one coulomb is one ampere second}{Coulomb -- die Einheit der elektrischen Ladung: ein Coulomb ist eine Amperesekunde}
\lies{$1\eh{V} = 1\eh{W/A}$}{volt \ipa{v@Ult} -- the unit of voltage: one volt is one watt per ampere}{Volt -- die Einheit der elektrischen Spannung: ein Volt ist ein Watt pro Ampere}
\lies{$1\eh{\Omega} = 1\eh{V/A}$}{ohm \ipa{@Um} -- the unit of resistance: one ohm is one volt per ampere}{Ohm -- die Einheit des elektrischen Widerstands: ein Ohm ist ein Volt pro Ampere}
\lies{$1\eh{F} = 1\eh{C/V}$}{farad \ipa{"f\ae r\ae d} -- the unit of capacitance: one farad is one coulomb per volt}{Farad -- die Einheit der Kapazität: ein Farad ist ein Coulomb pro Volt}
\lies{$1\eh{T} = 1\eh{V \cdot s/m^2}$}{tesla \ipa{"tesl@} -- the unit of magnetic flux density}{Tesla -- die Einheit der magnetischen Flussdichte}
\lies{$1\eh{Wb} = 1\eh{V \cdot s}$}{weber \ipa{"veIb@} -- the unit of magnetic flux: one weber is one volt second}{Weber -- die Einheit des magnetischen Flusses: ein Weber ist eine Voltsekunde}
\lies{$1\eh{H} = 1\eh{V \cdot s/A}$}{henry \ipa{"henri} -- the unit of inductance}{Henry -- die Einheit der Induktivität}
\lies{$1\eh{Bq} = 1\eh{s^{-1}}$}{becquerel \ipa{"bek@rel} -- the unit of activity: one decay per second}{Becquerel -- die Einheit der Aktivität: ein Zerfall pro Sekunde}
\lies{$\mathrm{lm},\ \mathrm{lx}$}{lumen \ipa{"lu:mIn} and lux \ipa{l2ks} -- the units of luminous flux and illuminance}{Lumen und Lux -- die Einheiten von Lichtstrom und Beleuchtungsstärke}
\lies{$\mathrm{\degree C}$}{degree Celsius \ipa{"selsi@s} -- the everyday unit of temperature}{Grad Celsius -- die Einheit der Temperatur im Alltag}
\end{lesen}

\platz\zwischen{Beispiele}{Examples}{Ig"zA:mp\textsyllabic{l}z}

\begin{lesen}
\lies{$230\eh{V}$}{The mains voltage is two hundred and thirty volts.}{Die Netzspannung beträgt zweihundertdreißig Volt.}
\lies{$50\eh{Hz}$}{The mains frequency is fifty hertz.}{Die Netzfrequenz beträgt fünfzig Hertz.}
\lies{$9{,}81\eh{N}$}{A mass of one kilogram has a weight of nine point eight one newtons.}{Ein Körper der Masse ein Kilogramm hat eine Gewichtskraft von neun Komma acht eins Newton.}
\lies{$1013\eh{hPa}$}{Standard atmospheric pressure is one thousand and thirteen hectopascals.}{Der Normalluftdruck beträgt eintausenddreizehn Hektopascal.}
\lies{$60\eh{W}$}{a sixty-watt lamp}{eine Sechzig-Watt-Lampe}
\end{lesen}

%--------------------------------------------------------------------------------------------------
\newpage
\kapitel{Vorsätze für Einheiten}{Prefixes for units}{"pri:fIksIz f@ "ju:nIts}

Vorsätze (\textit{prefixes}) bezeichnen dezimale Vielfache und Teile einer Einheit.

\begin{center}
{\small\renewcommand{\arraystretch}{1.25}
\begin{tabular}{@{}c l l @{\hspace{1.4cm}} c l l@{}}
	\toprule
	\textbf{Zeichen} & \textbf{Vorsatz (engl.)} & \textbf{Faktor} & \textbf{Zeichen} & \textbf{Vorsatz (engl.)} & \textbf{Faktor} \\ \midrule
	T & tera- \ipa{"ter@} & $10^{12}$ & d & deci- \ipa{"desI} & $10^{-1}$ \\
	G & giga- \ipa{"gIg@} & $10^{9}$ & c & centi- \ipa{"sentI} & $10^{-2}$ \\
	M & mega- \ipa{"meg@} & $10^{6}$ & m & milli- \ipa{"mIlI} & $10^{-3}$ \\
	k & kilo- \ipa{"kIl@} & $10^{3}$ & $\mu$ & micro- \ipa{"maIkr@U} & $10^{-6}$ \\
	h & hecto- \ipa{"hekt@U} & $10^{2}$ & n & nano- \ipa{"n\ae n@U} & $10^{-9}$ \\
	da & deca- \ipa{"dek@} & $10^{1}$ & p & pico- \ipa{"pi:k@U} & $10^{-12}$ \\
	\bottomrule
\end{tabular}}
\end{center}

\anmerkung{Die Betonung liegt im Englischen fast immer auf dem Vorsatz: \textit{kilometre}
\ipa{"kIl@mi:t@} (häufig auch \ipa{kI"l6mIt@}), \textit{millimetre} \ipa{"mIlimi:t@},
\textit{megawatt} \ipa{"meg@w6t}.}

\platz\zwischen{Beispiele}{Examples}{Ig"zA:mp\textsyllabic{l}z}

\begin{lesen}
\lies{$5\eh{km} = 5000\eh{m}$}{five kilometres equals five thousand metres}{fünf Kilometer sind fünftausend Meter}
\lies{$3\eh{ms} = 0{,}003\eh{s}$}{three milliseconds equals nought point nought nought three seconds}{drei Millisekunden sind null Komma null null drei Sekunden}
\lies{$2{,}2\eh{k\Omega}$}{two point two kilohms \ipa{"kIl@Umz}}{zwei Komma zwei Kiloohm}
\lies{$470\eh{\mu F}$}{four hundred and seventy microfarads}{vierhundertsiebzig Mikrofarad}
\lies{$650\eh{nm}$}{six hundred and fifty nanometres (the wavelength of red light)}{sechshundertfünfzig Nanometer (die Wellenlänge von rotem Licht)}
\lies{$1{,}2\eh{GW}$}{one point two gigawatts}{eins Komma zwei Gigawatt}
\end{lesen}

%--------------------------------------------------------------------------------------------------
\newpage
\kapitel{Einheiten und Größenangaben lesen}{Reading units and quantities}{"ri:dIN "ju:nIts @n "kw6nt@tiz}

\begin{lesen}
\lies{$\mathrm{m/s}$ \;oder\; $\mathrm{m \cdot s^{-1}}$}{metres per second}{Meter pro Sekunde \;/\; Meter durch Sekunde}
\lies{$\mathrm{m/s^2}$}{metres per second squared \;/\; metres per second per second}{Meter pro Sekunde zum Quadrat \;/\; Meter pro Quadratsekunde}
\lies{$\mathrm{km/h}$}{kilometres per hour}{Kilometer pro Stunde}
\lies{$\mathrm{m^2},\ \mathrm{m^3}$}{square metres, cubic metres}{Quadratmeter, Kubikmeter}
\lies{$\mathrm{kg/m^3}$}{kilograms per cubic metre}{Kilogramm pro Kubikmeter}
\lies{$\mathrm{N \cdot m}$}{newton metres}{Newtonmeter}
\lies{$\mathrm{N/m}$}{newtons per metre}{Newton pro Meter}
\lies{$\mathrm{J/(kg \cdot K)}$}{joules per kilogram kelvin}{Joule pro Kilogramm und Kelvin}
\lies{$\mathrm{kWh}$}{kilowatt-hours}{Kilowattstunden}
\lies{$20\eh{\degree C}$}{twenty degrees Celsius}{zwanzig Grad Celsius}
\lies{$3 \cdot 10^8\eh{m/s}$ \newline \footnotesize engl. $3 \times 10^8\eh{m/s}$}{three times ten to the eight metres per second}{drei mal zehn hoch acht Meter pro Sekunde}
\lies{$1{,}6 \cdot 10^{-19}\eh{C}$}{one point six times ten to the minus nineteen coulombs}{eins Komma sechs mal zehn hoch minus neunzehn Coulomb}
\lies{$(2{,}50 \pm 0{,}01)\eh{m}$}{two point five nought plus or minus nought point nought one metres}{zwei Komma fünf null plus minus null Komma null eins Meter}
\end{lesen}

\anmerkung{Gesprochen stehen Einheiten im Englischen im Plural, sobald der Zahlenwert nicht genau 1
ist: \textit{five newtons, nought point five volts}. Das Einheitenzeichen bekommt nie ein Plural-s
(5\,N, nicht 5\,Ns). Unverändert bleiben \textit{hertz, lux, siemens} und meist \textit{kelvin}.}

\platz\zwischen{Beispiele}{Examples}{Ig"zA:mp\textsyllabic{l}z}

\begin{lesen}
\lies{$v = 72\eh{km/h} = 20\eh{m/s}$}{seventy-two kilometres per hour equals twenty metres per second}{zweiundsiebzig Kilometer pro Stunde sind zwanzig Meter pro Sekunde}
\lies{$\varrho = 1000\eh{kg/m^3}$}{The density of water is one thousand kilograms per cubic metre.}{Die Dichte von Wasser beträgt eintausend Kilogramm pro Kubikmeter.}
\end{lesen}

%--------------------------------------------------------------------------------------------------
\newpage
\kapitel{Angloamerikanische Einheiten}{Imperial and US units}{Im"pI@ri@l @n ""ju: "es "ju:nIts}

In Großbritannien und vor allem in den USA sind neben den SI-Einheiten ältere Einheiten im
Gebrauch (\textit{imperial units} bzw.\ \textit{US customary units}).

\begin{center}
\begin{tikzpicture}[font=\footnotesize, x=1cm]
	\draw[blue, thick, fill=blue!6] (0,0) rectangle (13.2,2.1);
	\foreach \k in {0,...,5} {\draw[blue, thick] (0.25+2.54*\k,2.1) -- ++(0,-0.45) node[below, black] {\k};}
	\foreach \k in {0,...,19} {\draw[blue] (0.25+0.635*\k,2.1) -- ++(0,-0.22);}
	\foreach \k in {0,...,12} {\draw[red, thick] (0.25+\k,0) -- ++(0,0.35) node[above, black] {\k};}
	\foreach \k in {0,...,24} {\draw[red] (0.25+0.5*\k,0) -- ++(0,0.2);}
	\node[anchor=west] at (13.4,1.8) {inches -- \color{black!60}Zoll};
	\node[anchor=west] at (13.4,0.25) {centimetres -- \color{black!60}Zentimeter};
\end{tikzpicture}
\end{center}

\platz\zwischen{Länge}{Length}{leNT}
\begin{lesen}
\lies{$1\eh{in} = 2{,}54\eh{cm}$ \newline \footnotesize auch $1''$}{one inch \ipa{IntS} equals two point five four centimetres}{ein Zoll sind zwei Komma fünf vier Zentimeter}
\lies{$1\eh{ft} = 12\eh{in} = 30{,}48\eh{cm}$ \newline \footnotesize auch $1'$}{one foot (plural: feet) equals twelve inches}{ein Fuß sind zwölf Zoll}
\lies{$1\eh{yd} = 3\eh{ft} = 91{,}44\eh{cm}$}{one yard \ipa{jA:d} equals three feet}{ein Yard sind drei Fuß}
\lies{$1\eh{mi} = 1760\eh{yd} \approx 1{,}609\eh{km}$}{one mile equals one thousand seven hundred and sixty yards}{eine Meile sind eintausendsiebenhundertsechzig Yard}
\end{lesen}

\platz\zwischen{Volumen und Masse}{Volume and mass}{"v6lju:m @n "m\ae s}
\begin{lesen}
\lies{$1\eh{pt} \approx 568\eh{ml}$ (GB) \newline $1\eh{pt} \approx 473\eh{ml}$ (US)}{one pint \ipa{paInt}}{ein Pint}
\lies{$1\eh{gal} \approx 4{,}546\eh{l}$ (GB) \newline $1\eh{gal} \approx 3{,}785\eh{l}$ (US)}{one gallon \ipa{"g\ae l@n} equals eight pints}{eine Gallone sind acht Pints}
\lies{$1\eh{oz} \approx 28{,}35\eh{g}$}{one ounce \ipa{aUns}}{eine Unze}
\lies{$1\eh{lb} = 16\eh{oz} \approx 453{,}6\eh{g}$}{one pound \ipa{paUnd} equals sixteen ounces}{ein Pfund sind sechzehn Unzen}
\lies{$1\eh{st} = 14\eh{lb} \approx 6{,}35\eh{kg}$}{one stone equals fourteen pounds (used in Britain for body weight)}{ein Stone sind vierzehn Pfund (in Großbritannien für das Körpergewicht)}
\end{lesen}

\platz\zwischen{Geschwindigkeit, Druck, Leistung, Energie}{Speed, pressure, power, energy}{"spi:d, "preS@, "paU@, "en@dZi}
\begin{lesen}
\lies{$1\eh{mph} \approx 1{,}609\eh{km/h}$}{one mile per hour}{eine Meile pro Stunde}
\lies{$1\eh{psi} \approx 6895\eh{Pa}$}{one pound per square inch}{ein Pfund pro Quadratzoll}
\lies{$1\eh{hp} \approx 746\eh{W}$}{one horsepower \ipa{"hO:spaU@}}{eine (britische) Pferdestärke; $1\eh{PS} \approx 735{,}5\eh{W}$}
\lies{$1\eh{Btu} \approx 1055\eh{J}$}{one British thermal unit}{eine britische Wärmeeinheit}
\lies{$1\eh{cal} \approx 4{,}19\eh{J}$}{one calorie \ipa{"k\ae l@ri}}{eine Kalorie}
\end{lesen}

\anmerkung{Vorsicht bei \textit{gallon} und \textit{pint}: Die britischen Einheiten sind etwa 20\,\%
größer als die amerikanischen. Körpergröße wird in \textit{feet} und \textit{inches} angegeben:
$5'\,11''$ -- \textit{five foot eleven}.}

\platz\zwischen{Beispiele}{Examples}{Ig"zA:mp\textsyllabic{l}z}

\begin{lesen}
\lies{$6\eh{ft} \approx 1{,}83\eh{m}$}{He is six feet tall, that is about one metre eighty-three.}{Er ist sechs Fuß groß, das sind etwa ein Meter dreiundachtzig.}
\lies{$70\eh{mph} \approx 113\eh{km/h}$}{The speed limit on British motorways is seventy miles per hour.}{Auf britischen Autobahnen gilt ein Tempolimit von siebzig Meilen pro Stunde.}
\lies{$27\eh{in} \approx 68{,}6\eh{cm}$}{a twenty-seven-inch monitor}{ein Siebenundzwanzig-Zoll-Monitor}
\lies{$32\eh{psi} \approx 2{,}2\eh{bar}$}{The tyre pressure should be thirty-two pounds per square inch.}{Der Reifendruck soll etwa zwei Komma zwei Bar betragen.}
\lies{$150\eh{lb} \approx 68\eh{kg}$}{one hundred and fifty pounds is about sixty-eight kilograms}{hundertfünfzig Pfund sind etwa achtundsechzig Kilogramm}
\end{lesen}

%--------------------------------------------------------------------------------------------------
\newpage
\kapitel{Temperaturskalen}{Temperature scales}{"tempr@tS@ skeIlz}

\begin{center}
\begin{tikzpicture}[font=\small]
	\foreach \y in {0.3,2.4,3.6,5.0} {\draw[black!30, thin] (-0.3,\y) -- (9.2,\y);}
	\foreach \x/\name in {0/Kelvin, 3.4/Celsius, 6.8/Fahrenheit} {
		\draw[blue, very thick] (\x,0) -- (\x,1.0); \draw[blue, very thick, dashed] (\x,1.0) -- (\x,1.7); \draw[blue, very thick] (\x,1.7) -- (\x,5.5);
		\node[above, font=\small\bfseries] at (\x,5.6) {\name};
		\foreach \y in {0.3,2.4,3.6,5.0} {\draw[blue, thick] (\x-0.12,\y) -- (\x+0.12,\y);}}
	\foreach \y/\k/\c/\f in {0.3/0\eh{K}/-273{,}15\eh{\degree C}/-459{,}67\eh{\degree F}, 2.4/273{,}15\eh{K}/0\eh{\degree C}/32\eh{\degree F},
			3.6/310{,}15\eh{K}/37\eh{\degree C}/98{,}6\eh{\degree F}, 5.0/373{,}15\eh{K}/100\eh{\degree C}/212\eh{\degree F}} {
		\node[right, fill=white, inner sep=2pt, xshift=5pt] at (0,\y) {$\k$}; \node[right, fill=white, inner sep=2pt, xshift=5pt] at (3.4,\y) {$\c$};
		\node[right, fill=white, inner sep=2pt, xshift=5pt] at (6.8,\y) {$\f$};}
	\node[anchor=west] at (9.4,0.3) {absolute zero -- \color{black!60}absoluter Nullpunkt};
	\node[anchor=west] at (9.4,2.4) {freezing point of water -- \color{black!60}Gefrierpunkt};
	\node[anchor=west] at (9.4,3.6) {body temperature -- \color{black!60}Körpertemperatur};
	\node[anchor=west] at (9.4,5.0) {boiling point of water -- \color{black!60}Siedepunkt};
\end{tikzpicture}
\end{center}

\begin{lesen}
\lies{$\vartheta_C,\ \vartheta_F$}{the temperature in degrees Celsius \ipa{"selsi@s} and in degrees Fahrenheit \ipa{"f\ae r@nhaIt}}{die Temperatur in Grad Celsius und in Grad Fahrenheit}
\lies{$T$}{the absolute (thermodynamic) temperature in kelvin}{die absolute (thermodynamische) Temperatur in Kelvin}
\lies{$T = \vartheta_C + 273{,}15$}{To convert degrees Celsius into kelvin, add two hundred and seventy-three point one five.}{Zur Umrechnung von Grad Celsius in Kelvin addiert man zweihundertdreiundsiebzig Komma eins fünf.}
\lies{$\vartheta_F = \dfrac{9}{5} \cdot \vartheta_C + 32$}{To convert degrees Celsius into degrees Fahrenheit, multiply by nine fifths and add thirty-two.}{Zur Umrechnung von Grad Celsius in Grad Fahrenheit multipliziert man mit neun Fünfteln und addiert zweiunddreißig.}
\lies{$\vartheta_C = \dfrac{5}{9} \cdot (\vartheta_F - 32)$}{To convert degrees Fahrenheit into degrees Celsius, subtract thirty-two and multiply by five ninths.}{Zur Umrechnung von Grad Fahrenheit in Grad Celsius subtrahiert man zweiunddreißig und multipliziert mit fünf Neunteln.}
\lies{$\Delta T = 15\eh{K}$}{a temperature difference (a rise in temperature) of fifteen kelvin}{eine Temperaturdifferenz (ein Temperaturanstieg) von fünfzehn Kelvin}
\end{lesen}

\anmerkung{Die Umrechnungsformeln sind Zahlenwertgleichungen: Man setzt nur die Zahlenwerte ein.
Statt \textit{degrees Celsius} hört man im Alltag auch noch \textit{degrees centigrade}
\ipa{"sentIgreId}. Temperaturdifferenzen gibt man in Kelvin an.}

\platz\zwischen{Beispiele}{Examples}{Ig"zA:mp\textsyllabic{l}z}

\begin{lesen}
\lies{$68\eh{\degree F} \mathrel{\widehat{=}} 20\eh{\degree C}$}{Sixty-eight degrees Fahrenheit corresponds to twenty degrees Celsius: sixty-eight minus thirty-two is thirty-six; thirty-six times five ninths is twenty.}{Achtundsechzig Grad Fahrenheit entsprechen zwanzig Grad Celsius: achtundsechzig minus zweiunddreißig ist sechsunddreißig; sechsunddreißig mal fünf Neuntel ist zwanzig.}
\lies{$-40\eh{\degree C} \mathrel{\widehat{=}} -40\eh{\degree F}$}{At minus forty degrees the Celsius and the Fahrenheit scales show the same value.}{Bei minus vierzig Grad zeigen die Celsius- und die Fahrenheit-Skala denselben Wert.}
\lies{$-5\eh{\degree C}$}{minus five degrees Celsius \;/\; five degrees below zero}{minus fünf Grad Celsius \;/\; fünf Grad unter null}
\lies{$0\eh{K}$}{Absolute zero is the lowest possible temperature.}{Der absolute Nullpunkt ist die tiefste mögliche Temperatur.}
\end{lesen}

%--------------------------------------------------------------------------------------------------
\newpage
\kapitel{Formeln lesen: Mechanik}{Reading formulae: mechanics}{"ri:dIN "fO:mjuli:, mI"k\ae nIks}

\begin{lesen}
\lies{$v = \dfrac{s}{t}$}{$v$ equals $s$ over $t$: speed equals distance divided by time}{$v$ gleich $s$ durch $t$: Geschwindigkeit gleich Weg durch Zeit}
\lies{$a = \dfrac{\Delta v}{\Delta t}$}{$a$ equals delta $v$ over delta $t$: acceleration is the change in velocity divided by the time taken}{$a$ gleich Delta $v$ durch Delta $t$: Die Beschleunigung ist die Geschwindigkeitsänderung durch die dafür benötigte Zeit}
\lies{$s = \dfrac{1}{2} \cdot a \cdot t^2$}{$s$ equals a half $a$ $t$ squared}{$s$ gleich ein Halb $a$ $t$ Quadrat}
\lies{$F = m \cdot a$}{$F$ equals $m$ $a$: force equals mass times acceleration}{$F$ gleich $m$ mal $a$: Kraft gleich Masse mal Beschleunigung}
\lies{$F_G = m \cdot g$}{weight equals mass times the acceleration due to gravity}{Gewichtskraft gleich Masse mal Fallbeschleunigung}
\lies{$\varrho = \dfrac{m}{V}$}{rho equals $m$ over $V$: density equals mass divided by volume}{Rho gleich $m$ durch $V$: Dichte gleich Masse durch Volumen}
\lies{$p = \dfrac{F}{A}$}{$p$ equals $F$ over $A$: pressure equals force divided by area}{$p$ gleich $F$ durch $A$: Druck gleich Kraft durch Fläche}
\lies{$W = F \cdot s$}{work done equals force times distance (moved in the direction of the force)}{Arbeit gleich Kraft mal Weg (in Kraftrichtung)}
\lies{$E_{\mathrm{kin}} = \dfrac{1}{2} \cdot m \cdot v^2$}{kinetic energy equals a half $m$ $v$ squared}{kinetische Energie gleich ein Halb $m$ $v$ Quadrat}
\lies{$E_{\mathrm{pot}} = m \cdot g \cdot h$}{potential energy equals $m$ $g$ $h$}{potentielle Energie gleich $m$ $g$ $h$}
\lies{$P = \dfrac{W}{t}$}{power equals work done divided by time taken}{Leistung gleich Arbeit durch Zeit}
\lies{$F = G \cdot \dfrac{m_1 \cdot m_2}{r^2}$}{$F$ equals capital $G$ times $m$ one $m$ two over $r$ squared}{$F$ gleich groß $G$ mal $m$ eins mal $m$ zwei durch $r$ Quadrat}
\end{lesen}

\anmerkung{Einige Formelzeichen unterscheiden sich: Im Englischen steht meist $d$ oder $s$ für den
Weg (\textit{distance, displacement}), $k$ für die Federkonstante, $W$ für die Gewichtskraft
(\textit{weight}) und $u$ für die Anfangsgeschwindigkeit: $v = u + a\,t$.}

\platz\zwischen{Beispiele}{Examples}{Ig"zA:mp\textsyllabic{l}z}

\begin{lesen}
\lies{\textit{Gegeben -- gesucht}}{\textit{Given:} $m$ equals two kilograms, $a$ equals three metres per second squared. \textit{Find:} the force $F$.}{\textit{Gegeben:} $m$ gleich zwei Kilogramm, $a$ gleich drei Meter pro Sekunde zum Quadrat. \textit{Gesucht:} die Kraft $F$.}
\lies{$F = m \cdot a$ \newline $F = 2\eh{kg} \cdot 3\eh{m/s^2}$}{Write down the equation and substitute the values.}{Schreibe die Gleichung auf und setze die Werte ein.}
\lies{$F = 6\eh{N}$}{The force is six newtons.}{Die Kraft beträgt sechs Newton.}
\lies{$v = \dfrac{s}{t} \ \Rightarrow\ t = \dfrac{s}{v}$}{Rearrange the equation to make $t$ the subject \;/\; solve the equation for $t$.}{Stelle die Gleichung nach $t$ um \;/\; löse die Gleichung nach $t$ auf.}
\lies{$E_{\mathrm{kin}} = \frac12 \cdot 1000\eh{kg} \cdot (20\eh{m/s})^2$ \newline $E_{\mathrm{kin}} = 200\eh{kJ}$}{A car with a mass of one thousand kilograms travelling at twenty metres per second has a kinetic energy of two hundred kilojoules.}{Ein Auto der Masse eintausend Kilogramm hat bei zwanzig Meter pro Sekunde eine kinetische Energie von zweihundert Kilojoule.}
\end{lesen}

%--------------------------------------------------------------------------------------------------
\newpage
\kapitel{Formeln lesen: Elektrizitätslehre}{Reading formulae: electricity}{"ri:dIN "fO:mjuli:, I""lek"trIs@ti}

\begin{center}
\begin{circuitikz}[font=\footnotesize]
	\draw (0,0) to[battery1, l=$U$] (0,2.6) to[rmeter, t=A] (2.6,2.6) to[lamp, l=$L_1$] (5.2,2.6) to[R, l=$R$] (5.2,0) to[nos] (2.6,0) -- (0,0);
	\draw (5.2,2.6) -- (6.7,2.6) to[rmeter, t=V] (6.7,0) -- (5.2,0);
	\fill (5.2,2.6) circle (1.6pt) (5.2,0) circle (1.6pt);
	\node[anchor=west, align=left] at (8,2.5) {\sz{battery1}\; cell, battery -- \color{black!60}Batterie};
	\node[anchor=west, align=left] at (8,1.95) {\sz{nos}\; switch -- \color{black!60}Schalter};
	\node[anchor=west, align=left] at (8,1.4) {\sz{lamp}\; lamp -- \color{black!60}Lampe};
	\node[anchor=west, align=left] at (8,0.85) {\sz{R}\; resistor -- \color{black!60}Widerstand};
	\node[anchor=west, align=left] at (8,0.3) {\sz{rmeter, t=A}\; ammeter -- \color{black!60}Strommesser};
	\node[anchor=west, align=left] at (8,-0.25) {\sz{rmeter, t=V}\; voltmeter -- \color{black!60}Spannungsmesser};
\end{circuitikz}
\end{center}

\begin{lesen}
\lies{$I = \dfrac{Q}{t}$}{$I$ equals $Q$ over $t$: current is the charge flowing per unit time}{$I$ gleich $Q$ durch $t$: Die Stromstärke ist die pro Zeit fließende Ladung}
\lies{$U = R \cdot I$ \newline \footnotesize engl. $V = I\,R$}{$V$ equals $I$ $R$: voltage equals current times resistance (Ohm's law)}{$U$ gleich $R$ mal $I$: Spannung gleich Widerstand mal Stromstärke (ohmsches Gesetz)}
\lies{$P = U \cdot I$ \newline \footnotesize engl. $P = I\,V$}{power equals current times voltage}{Leistung gleich Spannung mal Stromstärke}
\lies{$E = U \cdot I \cdot t$}{energy transferred equals voltage times current times time}{elektrische Energie gleich Spannung mal Stromstärke mal Zeit}
\lies{$R_{\mathrm{ges}} = R_1 + R_2$}{in series: the total resistance equals $R$ one plus $R$ two}{Reihenschaltung: Der Gesamtwiderstand ist gleich $R$ eins plus $R$ zwei}
\lies{$\dfrac{1}{R_{\mathrm{ges}}} = \dfrac{1}{R_1} + \dfrac{1}{R_2}$}{in parallel: one over the total resistance equals one over $R$ one plus one over $R$ two}{Parallelschaltung: eins durch $R$ gesamt gleich eins durch $R$ eins plus eins durch $R$ zwei}
\lies{$C = \dfrac{Q}{U}$}{capacitance equals charge divided by voltage}{Kapazität gleich Ladung durch Spannung}
\lies{$\dfrac{U_1}{U_2} = \dfrac{N_1}{N_2}$}{the ratio of the voltages equals the ratio of the numbers of turns (transformer equation)}{die Spannungen verhalten sich wie die Windungszahlen (Transformatorgleichung)}
\end{lesen}

\anmerkung{Im Englischen ist das Formelzeichen für die Spannung $V$ (\textit{voltage}), nicht $U$ --
Zeichen und Einheit sehen dann gleich aus: $V = 12\eh{V}$. Man sagt \textit{the voltage across
the lamp} (an der Lampe) und \textit{the current through the lamp} (durch die Lampe).}

\platz\zwischen{Beispiele}{Examples}{Ig"zA:mp\textsyllabic{l}z}

\begin{lesen}
\lies{$I = \dfrac{U}{R} = \dfrac{6\eh{V}}{12\eh{\Omega}} = 0{,}5\eh{A}$}{A six-volt battery is connected to a twelve-ohm resistor. The current is nought point five amperes.}{An eine Sechs-Volt-Batterie wird ein Zwölf-Ohm-Widerstand angeschlossen. Die Stromstärke beträgt null Komma fünf Ampere.}
\lies{$R_{\mathrm{ges}} = 100\eh{\Omega} + 220\eh{\Omega} = 320\eh{\Omega}$}{Two resistors of one hundred ohms and two hundred and twenty ohms in series have a total resistance of three hundred and twenty ohms.}{Zwei Widerstände von hundert Ohm und zweihundertzwanzig Ohm in Reihe haben einen Gesamtwiderstand von dreihundertzwanzig Ohm.}
\lies{$R_{\mathrm{ges}} = \dfrac{60\eh{\Omega} \cdot 30\eh{\Omega}}{60\eh{\Omega} + 30\eh{\Omega}} = 20\eh{\Omega}$}{Sixty ohms and thirty ohms in parallel give a total resistance of twenty ohms.}{Sechzig Ohm und dreißig Ohm parallel ergeben einen Gesamtwiderstand von zwanzig Ohm.}
\lies{$E = 2\eh{kW} \cdot 0{,}5\eh{h} = 1\eh{kWh}$}{A two-kilowatt kettle switched on for half an hour transfers one kilowatt-hour of energy.}{Ein Wasserkocher mit zwei Kilowatt, der eine halbe Stunde eingeschaltet ist, wandelt eine Kilowattstunde Energie um.}
\lies{Schalter}{Close the switch: the circuit is complete and the lamp lights up. Open the switch: the lamp goes out.}{Schließe den Schalter: Der Stromkreis ist geschlossen und die Lampe leuchtet. Öffne den Schalter: Die Lampe erlischt.}
\end{lesen}

%--------------------------------------------------------------------------------------------------
\newpage
\kapitel{Formeln lesen: Wärme, Wellen, Optik und Atome}{Reading formulae: heat, waves, optics and atoms}{"ri:dIN "fO:mjuli:, "hi:t, "weIvz, "6ptIks @n "\ae t@mz}

\begin{lesen}
\lies{$Q = c \cdot m \cdot \Delta T$}{$Q$ equals $c$ $m$ delta $T$: heat equals specific heat capacity times mass times temperature change}{$Q$ gleich $c$ mal $m$ mal Delta $T$: Wärme gleich spezifische Wärmekapazität mal Masse mal Temperaturänderung}
\lies{$\dfrac{p_1 \cdot V_1}{T_1} = \dfrac{p_2 \cdot V_2}{T_2}$}{$p$ one $V$ one over $T$ one equals $p$ two $V$ two over $T$ two}{$p$ eins mal $V$ eins durch $T$ eins gleich $p$ zwei mal $V$ zwei durch $T$ zwei}
\lies{$f = \dfrac{1}{T}$}{frequency equals one over the period}{Frequenz gleich eins durch Periodendauer}
\lies{$c = \lambda \cdot f$ \newline \footnotesize engl. $v = f\,\lambda$}{wave speed equals frequency times wavelength \;/\; $v$ equals $f$ lambda}{Ausbreitungsgeschwindigkeit gleich Wellenlänge mal Frequenz \;/\; $c$ gleich Lambda mal $f$}
\lies{$T = 2\pi \cdot \sqrt{\dfrac{l}{g}}$}{$T$ equals two pi times the square root of $l$ over $g$}{$T$ gleich zwei pi mal Wurzel aus $l$ durch $g$}
\lies{$\dfrac{1}{f} = \dfrac{1}{g} + \dfrac{1}{b}$ \newline \footnotesize engl. $\frac{1}{f} = \frac{1}{u} + \frac{1}{v}$}{one over $f$ equals one over $u$ plus one over $v$ (the lens equation)}{eins durch $f$ gleich eins durch $g$ plus eins durch $b$ (die Linsengleichung)}
\lies{$n = \dfrac{\sin\alpha}{\sin\beta}$}{the refractive index equals sine alpha over sine beta \;/\; sine $i$ over sine $r$}{die Brechzahl ist gleich Sinus Alpha durch Sinus Beta}
\lies{$E = h \cdot f$}{$E$ equals $h$ $f$: the energy of a photon equals the Planck constant times the frequency}{$E$ gleich $h$ mal $f$: Die Energie eines Photons ist das plancksche Wirkungsquantum mal die Frequenz}
\lies{$E = m \cdot c^2$}{$E$ equals $m$ $c$ squared}{$E$ gleich $m$ $c$ Quadrat}
\lies{$N = N_0 \cdot \left(\dfrac{1}{2}\right)^{t/T_{1/2}}$}{$N$ equals $N$ nought times one half to the power of $t$ over the half-life}{$N$ gleich $N$ null mal ein Halb hoch $t$ durch die Halbwertszeit}
\lies{$^{238}_{\ 92}\mathrm{U}$}{uranium-two three eight: mass number two hundred and thirty-eight, atomic number ninety-two}{Uran-238: Massenzahl zweihundertachtunddreißig, Kernladungszahl zweiundneunzig}
\end{lesen}

\platz\zwischen{Beispiele}{Examples}{Ig"zA:mp\textsyllabic{l}z}

\begin{lesen}
\lies{$\lambda = \dfrac{c}{f} = \dfrac{3 \cdot 10^8\eh{m/s}}{100\eh{MHz}} = 3\eh{m}$}{A radio station broadcasts at one hundred megahertz. The wavelength is three metres.}{Ein Radiosender sendet mit hundert Megahertz. Die Wellenlänge beträgt drei Meter.}
\lies{$Q = 0{,}5\eh{kg} \cdot 4{,}19\eh{\tfrac{kJ}{kg \cdot K}} \cdot 80\eh{K}$ \newline $Q \approx 168\eh{kJ}$}{About one hundred and sixty-eight kilojoules are needed to heat half a kilogram of water from twenty to one hundred degrees Celsius.}{Um ein halbes Kilogramm Wasser von zwanzig auf hundert Grad Celsius zu erwärmen, sind etwa hundertachtundsechzig Kilojoule nötig.}
\lies{$N = N_0 \cdot \left(\frac12\right)^3 = \dfrac{N_0}{8}$}{After three half-lives one eighth of the original nuclei are left.}{Nach drei Halbwertszeiten ist noch ein Achtel der ursprünglichen Kerne vorhanden.}
\lies{$n = \dfrac{c_0}{c} = \dfrac{3{,}00 \cdot 10^8}{2{,}25 \cdot 10^8} \approx 1{,}33$}{The refractive index of water is about one point three three.}{Die Brechzahl von Wasser beträgt etwa eins Komma drei drei.}
\end{lesen}

%--------------------------------------------------------------------------------------------------
\newpage
\kapitel{Experimente beschreiben}{Describing experiments}{dI"skraIbIN Ik"sperIm@nts}

\begin{center}
\begin{tikzpicture}[font=\small]
	\draw[black, line width=2.5pt] (-1.3,0) -- (1.2,0);
	\draw[black, line width=2pt] (-0.8,0) -- (-0.8,5.2); \draw[black, line width=2pt] (-0.8,4.8) -- (1.9,4.8);
	\draw[blue, thick, decorate, decoration={zigzag, segment length=5pt, amplitude=5pt}] (1.5,4.8) -- (1.5,2.5);
	\draw[blue, thick] (1.5,2.5) -- (1.5,2.2); \filldraw[fill=blue!20, draw=blue, very thick] (1.15,1.5) rectangle (1.85,2.2);
	\draw[thick, fill=yellow!20] (2.4,0.3) rectangle (2.8,5.0); \foreach \y in {0.5,0.8,...,4.9} {\draw (2.4,\y) -- (2.55,\y);}
	\draw[red, thick, -{Stealth}] (1.5,1.5) -- (1.5,0.5) node[left] {$\vec{F}_G$};
	\node[anchor=west, align=left] (a) at (4.2,4.7) {clamp stand -- \color{black!60}Stativ};
	\node[anchor=west, align=left] (b) at (4.2,3.8) {spring -- \color{black!60}Feder};
	\node[anchor=west, align=left] (c) at (4.2,2.9) {ruler -- \color{black!60}Lineal};
	\node[anchor=west, align=left] (d) at (4.2,2.0) {mass (slotted masses) -- \color{black!60}Massestück};
	\node[anchor=west, align=left] (e) at (4.2,1.1) {weight (force) -- \color{black!60}Gewichtskraft};
	\draw[-{Stealth}, thin] (a.west) -- (1.0,4.9); \draw[-{Stealth}, thin] (b.west) -- (1.75,3.7); \draw[-{Stealth}, thin] (c.west) -- (2.85,2.9);
	\draw[-{Stealth}, thin] (d.west) -- (1.95,1.9); \draw[-{Stealth}, thin] (e.west) -- (1.65,1.0);
\end{tikzpicture}
\end{center}

Die Teile eines Versuchsprotokolls (\textit{lab report}):

\begin{lesen}
\lies{\textbf{Aufgabe}}{\textbf{Aim:} To investigate how the extension of a spring depends on the force applied.}{Untersuche, wie die Verlängerung einer Feder von der wirkenden Kraft abhängt.}
\lies{\textbf{Vermutung}}{\textbf{Hypothesis / Prediction:} The greater the force, the greater the extension.}{Je größer die Kraft, desto größer die Verlängerung.}
\lies{\textbf{Geräte}}{\textbf{Apparatus:} clamp stand, spring, slotted masses, ruler}{Stativ, Feder, Massestücke, Lineal}
\lies{\textbf{Durchführung}}{\textbf{Method:} Set up the apparatus as shown in the diagram. Measure the length of the unloaded spring. Add the masses one at a time and record the length of the spring each time.}{Baue den Versuch wie in der Skizze auf. Miss die Länge der unbelasteten Feder. Hänge nacheinander die Massestücke an und notiere jeweils die Länge der Feder.}
\lies{\textbf{Messwerte}}{\textbf{Results:} Record your readings in a table.}{Trage die Messwerte in eine Tabelle ein.}
\lies{\textbf{Auswertung}}{\textbf{Analysis / Conclusion:} Plot a graph of force against extension. The graph is a straight line through the origin, so the force is directly proportional to the extension.}{Zeichne ein Kraft-Verlängerungs-Diagramm. Der Graph ist eine Ursprungsgerade, also ist die Kraft proportional zur Verlängerung.}
\lies{\textbf{Fehlerbetrachtung}}{\textbf{Evaluation:} Sources of error are reading the ruler (parallax error) and the oscillation of the spring.}{Fehlerquellen sind das Ablesen am Lineal (Parallaxenfehler) und das Schwingen der Feder.}
\end{lesen}

\platz\zwischen{Beispiele}{Examples}{Ig"zA:mp\textsyllabic{l}z}

\begin{lesen}
\lies{Variablen}{Change only one variable at a time and keep all the other variables constant (a fair test).}{Ändere immer nur eine Größe und halte alle anderen Größen konstant.}
\lies{Messen}{Measure the time for ten oscillations and divide by ten.}{Miss die Zeit für zehn Schwingungen und dividiere durch zehn.}
\lies{Wiederholen}{Repeat each measurement three times and calculate the mean.}{Wiederhole jede Messung dreimal und berechne den Mittelwert.}
\lies{Beobachten}{What do you observe? -- The lamp gets brighter as the voltage increases.}{Was beobachtest du? -- Die Lampe leuchtet heller, wenn die Spannung steigt.}
\lies{Sicherheit}{Safety: wear safety goggles. Do not look directly into the laser beam.}{Sicherheit: Trage eine Schutzbrille. Blicke nicht direkt in den Laserstrahl.}
\end{lesen}

%--------------------------------------------------------------------------------------------------
\newpage
\kapitel{Diagramme beschreiben}{Describing graphs}{dI"skraIbIN "grA:fs}

\begin{center}
\begin{tikzpicture}[font=\small, scale=1.05]
	\draw[black!15, very thin, step=0.5] (0,0) grid (6,4);
	\draw[achse] (-0.3,0) -- (6.5,0) node[right] {$s$ in cm}; \draw[achse] (0,-0.3) -- (0,4.5) node[above] {$F$ in N};
	\foreach \x in {1,...,6} {\node[below, font=\footnotesize] at (\x,0) {\x};}
	\foreach \y/\l in {1/0{,}5, 2/1{,}0, 3/1{,}5, 4/2{,}0} {\node[left, font=\footnotesize] at (0,\y) {$\l$};}
	\draw[red, thick] (0,0) -- (6.2,4.03);
	\foreach \x/\y in {1/0.75, 2/1.2, 3/2.05, 4/2.5, 5/3.35, 6/3.85} {\draw[blue, very thick] (\x-0.1,\y-0.1) -- (\x+0.1,\y+0.1) (\x-0.1,\y+0.1) -- (\x+0.1,\y-0.1);}
	\draw[black!70] (2,1.3) -- (5,1.3) node[midway, below, font=\footnotesize] {$\Delta s$} -- (5,3.25) node[midway, right, font=\footnotesize] {$\Delta F$};
	\node[anchor=west, align=left] (a) at (7.6,3.9) {line of best fit -- \color{black!60}Ausgleichsgerade};
	\node[anchor=west, align=left] (b) at (7.6,3.1) {data point (reading) -- \color{black!60}Messpunkt};
	\node[anchor=west, align=left] (c) at (7.6,2.3) {gradient triangle -- \color{black!60}Anstiegsdreieck};
	\node[anchor=west, align=left] (d) at (7.6,1.5) {$x$-axis: independent variable\\\color{black!60}Rechtsachse: unabhängige Größe};
	\node[anchor=west, align=left] (e) at (7.6,0.5) {origin -- \color{black!60}Koordinatenursprung};
	\draw[-{Stealth}, thin] (a.west) -- (5.9,3.9); \draw[-{Stealth}, thin] (b.west) -- (5.2,3.3); \draw[-{Stealth}, thin] (c.west) -- (5.15,2.3);
	\draw[-{Stealth}, thin] (e.west) -- (0.15,0.1);
\end{tikzpicture}
\end{center}

\begin{lesen}
\lies{Achsen}{The extension is plotted on the $x$-axis (horizontal axis), the force on the $y$-axis (vertical axis).}{Auf der Rechtsachse ist die Verlängerung abgetragen, auf der Hochachse die Kraft.}
\lies{Auftragen}{Plot a graph of force against extension \;/\; plot $F$ against $s$.}{Zeichne ein $F$-$s$-Diagramm \;/\; trage $F$ über $s$ auf.}
\lies{Gerade}{The graph is a straight line through the origin: $F$ is directly proportional to $s$.}{Der Graph ist eine Ursprungsgerade: $F$ ist proportional zu $s$.}
\lies{$D = \dfrac{\Delta F}{\Delta s}$}{The gradient (slope) of the line gives the spring constant.}{Der Anstieg der Geraden ergibt die Federkonstante.}
\lies{Verlauf}{The curve rises steeply at first and then levels off.}{Die Kurve steigt zunächst steil an und flacht dann ab.}
\lies{Verlauf}{The current increases \;/\; decreases \;/\; remains constant as the voltage increases.}{Die Stromstärke steigt \;/\; sinkt \;/\; bleibt konstant, wenn die Spannung zunimmt.}
\lies{Fläche}{The area under a velocity–time graph gives the distance travelled.}{Die Fläche unter dem Graphen im $v$-$t$-Diagramm entspricht dem zurückgelegten Weg.}
\lies{Ausreißer}{This reading is an anomalous result (an outlier); it is ignored when drawing the line of best fit.}{Dieser Messwert ist ein Ausreißer; er wird beim Zeichnen der Ausgleichsgeraden nicht berücksichtigt.}
\end{lesen}

\platz\zwischen{Beispiele}{Examples}{Ig"zA:mp\textsyllabic{l}z}

\begin{lesen}
\lies{$D = \dfrac{1{,}0\eh{N}}{3{,}1\eh{cm}} \approx 32\eh{N/m}$}{From the gradient triangle the spring constant is about thirty-two newtons per metre.}{Aus dem Anstiegsdreieck ergibt sich eine Federkonstante von etwa zweiunddreißig Newton pro Meter.}
\lies{\skST}{The distance–time graph is a straight line: the object moves at constant speed.}{Das Weg-Zeit-Diagramm ist eine Gerade: Der Körper bewegt sich mit konstanter Geschwindigkeit.}
\lies{\skVT}{The velocity increases uniformly with time: the acceleration is constant.}{Die Geschwindigkeit nimmt gleichmäßig mit der Zeit zu: Die Beschleunigung ist konstant.}
\lies{\skAntiprop}{The pressure is inversely proportional to the volume: if the volume is halved, the pressure doubles.}{Der Druck ist umgekehrt proportional zum Volumen: Wird das Volumen halbiert, verdoppelt sich der Druck.}
\lies{\skZerfall}{The number of undecayed nuclei decreases exponentially with time.}{Die Anzahl der noch nicht zerfallenen Kerne nimmt exponentiell mit der Zeit ab.}
\end{lesen}

%--------------------------------------------------------------------------------------------------
\newpage
\kapitel{Griechische Buchstaben in der Physik}{Greek letters in physics}{"gri:k "let@z In "fIzIks}

\begin{center}
{\small\renewcommand{\arraystretch}{1.25}
\begin{tabular}{@{}c l l l@{}}
	\toprule
	\textbf{Zeichen} & \textbf{Englisch} & \textbf{steht für (englisch)} & \textbf{steht für (deutsch)} \\ \midrule
	$\alpha$ & alpha \ipa{"\ae lf@} & angle; alpha particle & Winkel; Alphateilchen \\
	$\beta$ & beta \ipa{"bi:t@} & angle; beta particle & Winkel; Betateilchen \\
	$\gamma$ & gamma \ipa{"g\ae m@} & gamma radiation & Gammastrahlung \\
	$\Delta$ & delta \ipa{"delt@} & change in, difference & Änderung, Differenz \\
	$\varepsilon$ & epsilon \ipa{"epsIl6n} & permittivity & Permittivität (Dielektrizitätszahl) \\
	$\eta$ & eta \ipa{"i:t@} & efficiency & Wirkungsgrad \\
	$\vartheta$ & theta \ipa{"Ti:t@} & temperature; angle & Temperatur; Winkel \\
	$\lambda$ & lambda \ipa{"l\ae md@} & wavelength; decay constant & Wellenlänge; Zerfallskonstante \\
	$\mu$ & mu \ipa{mju:} & coefficient of friction; micro- & Reibungszahl; Mikro- \\
	$\varrho$ & rho \ipa{r@U} & density; resistivity & Dichte; spezifischer Widerstand \\
	$\sigma$ & sigma \ipa{"sIgm@} & stress; conductivity & mechanische Spannung; Leitfähigkeit \\
	$\tau$ & tau \ipa{taU} & time constant & Zeitkonstante \\
	$\varphi$ & phi \ipa{faI} & angle; phase; potential & Winkel; Phase; Potential \\
	$\Phi$ & capital phi & magnetic flux & magnetischer Fluss \\
	$\omega$ & omega \ipa{"@UmIg@} & angular velocity & Winkelgeschwindigkeit \\
	$\Omega$ & capital omega & ohm & Ohm \\
	\bottomrule
\end{tabular}}
\end{center}

\platz\zwischen{Beispiele}{Examples}{Ig"zA:mp\textsyllabic{l}z}

\begin{lesen}
\lies{$\Delta T = T_2 - T_1$}{delta $T$ equals $T$ two minus $T$ one: the change in temperature}{Delta $T$ gleich $T$ zwei minus $T$ eins: die Temperaturänderung}
\lies{$\eta = 0{,}85$}{eta equals nought point eight five: the efficiency is eighty-five per cent}{Eta gleich null Komma acht fünf: Der Wirkungsgrad beträgt fünfundachtzig Prozent}
\lies{$\lambda = 589\eh{nm}$}{lambda equals five hundred and eighty-nine nanometres (yellow sodium light)}{Lambda gleich fünfhundertneunundachtzig Nanometer (gelbes Natriumlicht)}
\lies{$\varrho = 7{,}87\eh{g/cm^3}$}{rho equals seven point eight seven grams per cubic centimetre (the density of iron)}{Rho gleich sieben Komma acht sieben Gramm pro Kubikzentimeter (die Dichte von Eisen)}
\lies{$R = 47\eh{\Omega}$}{$R$ equals forty-seven ohms}{$R$ gleich siebenundvierzig Ohm}
\end{lesen}


\newpage
%%%%  REGISTER  %%%%
\teil{Register}{Index}{"Indeks}
\section{Register Deutsch -- Englisch}
\begin{multicols}{2}\small\raggedright

\medskip{\large\textbf{\color{cornflower}A}}\par\nopagebreak
\textbf{Abbildungsmaßstab} (der) -- magnification\dotfill\hyperref[e:252]{\pageref*{e:252}}\par
\textbf{abgeleitete Einheit} (die) -- derived unit\dotfill\hyperref[e:5]{\pageref*{e:5}}\par
\textbf{absolute Nullpunkt} (der) -- absolute zero\dotfill\hyperref[e:115]{\pageref*{e:115}}\par
\textbf{absolute Temperatur} (die) -- absolute temperature, thermodynamic temperature\dotfill\hyperref[e:114]{\pageref*{e:114}}\par
\textbf{Abwärme} (die) -- waste heat\dotfill\hyperref[e:108]{\pageref*{e:108}}\par
\textbf{Aggregatzustand} (der) -- state of matter\dotfill\hyperref[e:125]{\pageref*{e:125}}\par
\textbf{Aktivität} (die) -- activity\dotfill\hyperref[e:283]{\pageref*{e:283}}\par
\textbf{Alphastrahlung} (die) -- alpha radiation\dotfill\hyperref[e:275]{\pageref*{e:275}}\par
\textbf{Amplitude} (die) -- amplitude\dotfill\hyperref[e:195]{\pageref*{e:195}}\par
\textbf{Angriffspunkt} (der) -- point of application\dotfill\hyperref[e:60]{\pageref*{e:60}}\par
\textbf{Äquivalentdosis} (die) -- equivalent dose\dotfill\hyperref[e:284]{\pageref*{e:284}}\par
\textbf{Arbeit} (die) -- work\dotfill\hyperref[e:96]{\pageref*{e:96}}\par
\textbf{Atom} (das) -- atom\dotfill\hyperref[e:259]{\pageref*{e:259}}\par
\textbf{Atomhülle} (die) -- electron shell\dotfill\hyperref[e:261]{\pageref*{e:261}}\par
\textbf{Atomkern} (der) -- nucleus\dotfill\hyperref[e:260]{\pageref*{e:260}}\par
\textbf{Auftrieb} (der) -- buoyancy, upthrust\dotfill\hyperref[e:92]{\pageref*{e:92}}\par
\textbf{Ausbreitungsgeschwindigkeit} (die) -- wave speed\dotfill\hyperref[e:203]{\pageref*{e:203}}\par
\textbf{Ausgleichsgerade} (die) -- line of best fit\dotfill\hyperref[e:24]{\pageref*{e:24}}\par
\textbf{Auslenkung} (die) -- displacement\dotfill\hyperref[e:194]{\pageref*{e:194}}\par
\textbf{Auswertung} (die) -- analysis, evaluation\dotfill\hyperref[e:19]{\pageref*{e:19}}\par

\medskip{\large\textbf{\color{cornflower}B}}\par\nopagebreak
\textbf{Bahn} (die) -- path, trajectory\dotfill\hyperref[e:33]{\pageref*{e:33}}\par
\textbf{Bahngeschwindigkeit} (die) -- tangential speed\dotfill\hyperref[e:55]{\pageref*{e:55}}\par
\textbf{Basiseinheit} (die) -- base unit\dotfill\hyperref[e:4]{\pageref*{e:4}}\par
\textbf{Batterie} (die) -- battery, cell\dotfill\hyperref[e:156]{\pageref*{e:156}}\par
\textbf{Beleuchtungsstärke} (die) -- illuminance\dotfill\hyperref[e:258]{\pageref*{e:258}}\par
\textbf{Beobachtung} (die) -- observation\dotfill\hyperref[e:18]{\pageref*{e:18}}\par
\textbf{Beschleunigung} (die) -- acceleration\dotfill\hyperref[e:40]{\pageref*{e:40}}\par
\textbf{Betastrahlung} (die) -- beta radiation\dotfill\hyperref[e:276]{\pageref*{e:276}}\par
\textbf{Betrag (einer Kraft)} (der) -- magnitude\dotfill\hyperref[e:61]{\pageref*{e:61}}\par
\textbf{Beugung} (die) -- diffraction\dotfill\hyperref[e:208]{\pageref*{e:208}}\par
\textbf{Bewegung} (die) -- motion\dotfill\hyperref[e:30]{\pageref*{e:30}}\par
\textbf{Bezugssystem} (das) -- frame of reference\dotfill\hyperref[e:32]{\pageref*{e:32}}\par
\textbf{Bild} (das) -- image\dotfill\hyperref[e:249]{\pageref*{e:249}}\par
\textbf{Bildweite} (die) -- image distance\dotfill\hyperref[e:247]{\pageref*{e:247}}\par
\textbf{Bindungsenergie} (die) -- binding energy\dotfill\hyperref[e:290]{\pageref*{e:290}}\par
\textbf{Brechung} (die) -- refraction\dotfill\hyperref[e:234]{\pageref*{e:234}}\par
\textbf{Brechungsgesetz} (das) -- Snell's law, law of refraction\dotfill\hyperref[e:237]{\pageref*{e:237}}\par
\textbf{Brechungswinkel} (der) -- angle of refraction\dotfill\hyperref[e:235]{\pageref*{e:235}}\par
\textbf{Brechzahl} (die) -- refractive index\dotfill\hyperref[e:236]{\pageref*{e:236}}\par
\textbf{Brennpunkt} (der) -- focal point, principal focus\dotfill\hyperref[e:244]{\pageref*{e:244}}\par
\textbf{Brennweite} (die) -- focal length\dotfill\hyperref[e:245]{\pageref*{e:245}}\par

\medskip{\large\textbf{\color{cornflower}C}}\par\nopagebreak
\textbf{chemische Energie} (die) -- chemical energy\dotfill\hyperref[e:102]{\pageref*{e:102}}\par

\medskip{\large\textbf{\color{cornflower}D}}\par\nopagebreak
\textbf{Dauermagnet} (der) -- permanent magnet\dotfill\hyperref[e:175]{\pageref*{e:175}}\par
\textbf{Dehnung, die Verlängerung} (die) -- extension\dotfill\hyperref[e:71]{\pageref*{e:71}}\par
\textbf{Diagramm} (das) -- graph\dotfill\hyperref[e:23]{\pageref*{e:23}}\par
\textbf{Dichte} (die) -- density\dotfill\hyperref[e:58]{\pageref*{e:58}}\par
\textbf{Doppler-Effekt} (der) -- Doppler effect\dotfill\hyperref[e:219]{\pageref*{e:219}}\par
\textbf{Drehmoment} (das) -- torque, moment (of a force)\dotfill\hyperref[e:81]{\pageref*{e:81}}\par
\textbf{Drehpunkt} (der) -- pivot, fulcrum\dotfill\hyperref[e:84]{\pageref*{e:84}}\par
\textbf{Drehzahl} (die) -- rotational speed\dotfill\hyperref[e:53]{\pageref*{e:53}}\par
\textbf{Druck} (der) -- pressure\dotfill\hyperref[e:89]{\pageref*{e:89}}\par
\textbf{Durchschnittsgeschwindigkeit} (die) -- average speed\dotfill\hyperref[e:38]{\pageref*{e:38}}\par

\medskip{\large\textbf{\color{cornflower}E}}\par\nopagebreak
\textbf{Echo} (das) -- echo\dotfill\hyperref[e:217]{\pageref*{e:217}}\par
\textbf{Eigenfrequenz} (die) -- natural frequency\dotfill\hyperref[e:199]{\pageref*{e:199}}\par
\textbf{Einfallslot} (das) -- normal\dotfill\hyperref[e:229]{\pageref*{e:229}}\par
\textbf{Einfallswinkel} (der) -- angle of incidence\dotfill\hyperref[e:230]{\pageref*{e:230}}\par
\textbf{Einheit} (die) -- unit\dotfill\hyperref[e:3]{\pageref*{e:3}}\par
\textbf{elektrische Energie} (die) -- electrical energy\dotfill\hyperref[e:101]{\pageref*{e:101}}\par
\textbf{elektrische Feld} (das) -- electric field\dotfill\hyperref[e:171]{\pageref*{e:171}}\par
\textbf{elektrische Feldstärke} (die) -- electric field strength\dotfill\hyperref[e:172]{\pageref*{e:172}}\par
\textbf{elektrische Ladung} (die) -- electric charge\dotfill\hyperref[e:143]{\pageref*{e:143}}\par
\textbf{elektrische Leistung} (die) -- electrical power\dotfill\hyperref[e:168]{\pageref*{e:168}}\par
\textbf{elektrische Spannung} (die) -- voltage, potential difference\dotfill\hyperref[e:145]{\pageref*{e:145}}\par
\textbf{elektrische Stromstärke} (die) -- electric current\dotfill\hyperref[e:144]{\pageref*{e:144}}\par
\textbf{elektrische Widerstand} (der) -- resistance\dotfill\hyperref[e:146]{\pageref*{e:146}}\par
\textbf{Elektromagnet} (der) -- electromagnet\dotfill\hyperref[e:176]{\pageref*{e:176}}\par
\textbf{elektromagnetische Induktion} (die) -- electromagnetic induction\dotfill\hyperref[e:184]{\pageref*{e:184}}\par
\textbf{Elektromotor} (der) -- electric motor\dotfill\hyperref[e:186]{\pageref*{e:186}}\par
\textbf{Elektron} (das) -- electron\dotfill\hyperref[e:262]{\pageref*{e:262}}\par
\textbf{Elementarladung} (die) -- elementary charge\dotfill\hyperref[e:266]{\pageref*{e:266}}\par
\textbf{Energie} (die) -- energy\dotfill\hyperref[e:95]{\pageref*{e:95}}\par
\textbf{Energieerhaltungssatz} (der) -- law of conservation of energy\dotfill\hyperref[e:104]{\pageref*{e:104}}\par
\textbf{Energieniveau} (das) -- energy level\dotfill\hyperref[e:294]{\pageref*{e:294}}\par
\textbf{Energieträger} (der) -- energy source\dotfill\hyperref[e:109]{\pageref*{e:109}}\par
\textbf{Energieumwandlung} (die) -- energy conversion, energy transfer\dotfill\hyperref[e:103]{\pageref*{e:103}}\par
\textbf{Entropie} (die) -- entropy\dotfill\hyperref[e:142]{\pageref*{e:142}}\par
\textbf{Erdung} (die) -- earth (BE), ground (AE)\dotfill\hyperref[e:189]{\pageref*{e:189}}\par
\textbf{erneuerbare Energie} (die) -- renewable energy\dotfill\hyperref[e:110]{\pageref*{e:110}}\par
\textbf{erstarren} -- freeze, solidify\dotfill\hyperref[e:130]{\pageref*{e:130}}\par

\medskip{\large\textbf{\color{cornflower}F}}\par\nopagebreak
\textbf{Fadenpendel} (das) -- simple pendulum\dotfill\hyperref[e:191]{\pageref*{e:191}}\par
\textbf{Fallbeschleunigung} (die) -- acceleration due to gravity\dotfill\hyperref[e:47]{\pageref*{e:47}}\par
\textbf{Farbzerlegung, die Dispersion} (die) -- dispersion\dotfill\hyperref[e:256]{\pageref*{e:256}}\par
\textbf{Feder} (die) -- spring\dotfill\hyperref[e:70]{\pageref*{e:70}}\par
\textbf{Federkonstante} (die) -- spring constant\dotfill\hyperref[e:72]{\pageref*{e:72}}\par
\textbf{Federschwinger} (der) -- mass on a spring\dotfill\hyperref[e:192]{\pageref*{e:192}}\par
\textbf{Feldlinie} (die) -- field line\dotfill\hyperref[e:173]{\pageref*{e:173}}\par
\textbf{fest} -- solid\dotfill\hyperref[e:126]{\pageref*{e:126}}\par
\textbf{Flaschenzug} (der) -- block and tackle\dotfill\hyperref[e:87]{\pageref*{e:87}}\par
\textbf{flüssig} -- liquid\dotfill\hyperref[e:127]{\pageref*{e:127}}\par
\textbf{Formelzeichen} (das) -- symbol\dotfill\hyperref[e:6]{\pageref*{e:6}}\par
\textbf{Fotoeffekt} (der) -- photoelectric effect\dotfill\hyperref[e:293]{\pageref*{e:293}}\par
\textbf{freie Fall} (der) -- free fall\dotfill\hyperref[e:46]{\pageref*{e:46}}\par
\textbf{Frequenz} (die) -- frequency\dotfill\hyperref[e:52]{\pageref*{e:52}}\par
\textbf{Frequenz} (die) -- frequency\dotfill\hyperref[e:197]{\pageref*{e:197}}\par

\medskip{\large\textbf{\color{cornflower}G}}\par\nopagebreak
\textbf{Gammastrahlung} (die) -- gamma radiation\dotfill\hyperref[e:277]{\pageref*{e:277}}\par
\textbf{gasförmig} -- gaseous\dotfill\hyperref[e:128]{\pageref*{e:128}}\par
\textbf{gedämpfte Schwingung} (die) -- damped oscillation\dotfill\hyperref[e:198]{\pageref*{e:198}}\par
\textbf{Gegenstandsweite} (die) -- object distance\dotfill\hyperref[e:246]{\pageref*{e:246}}\par
\textbf{geneigte Ebene} (die) -- inclined plane\dotfill\hyperref[e:85]{\pageref*{e:85}}\par
\textbf{Generator} (der) -- generator\dotfill\hyperref[e:187]{\pageref*{e:187}}\par
\textbf{Geschwindigkeit} (die) -- velocity, speed\dotfill\hyperref[e:37]{\pageref*{e:37}}\par
\textbf{Geschwindigkeit-Zeit-Diagramm} (das) -- velocity–time graph\dotfill\hyperref[e:45]{\pageref*{e:45}}\par
\textbf{Gesetz} (das) -- law\dotfill\hyperref[e:21]{\pageref*{e:21}}\par
\textbf{Gewichtskraft} (die) -- weight\dotfill\hyperref[e:63]{\pageref*{e:63}}\par
\textbf{gleichförmige Bewegung} (die) -- uniform motion\dotfill\hyperref[e:42]{\pageref*{e:42}}\par
\textbf{gleichmäßig beschleunigte Bewegung} (die) -- uniformly accelerated motion\dotfill\hyperref[e:43]{\pageref*{e:43}}\par
\textbf{Gleichstrom} (der) -- direct current (DC)\dotfill\hyperref[e:166]{\pageref*{e:166}}\par
\textbf{Gleitreibung} (die) -- sliding friction\dotfill\hyperref[e:68]{\pageref*{e:68}}\par
\textbf{Glühlampe} (die) -- lamp, bulb\dotfill\hyperref[e:158]{\pageref*{e:158}}\par
\textbf{Gravitation} (die) -- gravitation, gravity\dotfill\hyperref[e:93]{\pageref*{e:93}}\par
\textbf{Grenzwinkel} (der) -- critical angle\dotfill\hyperref[e:239]{\pageref*{e:239}}\par

\medskip{\large\textbf{\color{cornflower}H}}\par\nopagebreak
\textbf{Haftreibung} (die) -- static friction\dotfill\hyperref[e:67]{\pageref*{e:67}}\par
\textbf{Halbleiter} (der) -- semiconductor\dotfill\hyperref[e:152]{\pageref*{e:152}}\par
\textbf{Halbwertszeit} (die) -- half-life\dotfill\hyperref[e:280]{\pageref*{e:280}}\par
\textbf{Hauptquantenzahl} (die) -- principal quantum number\dotfill\hyperref[e:295]{\pageref*{e:295}}\par
\textbf{Hebel} (der) -- lever\dotfill\hyperref[e:82]{\pageref*{e:82}}\par
\textbf{Hebelarm} (der) -- lever arm\dotfill\hyperref[e:83]{\pageref*{e:83}}\par
\textbf{Hohlspiegel, der Wölbspiegel} (der) -- concave mirror, convex mirror\dotfill\hyperref[e:233]{\pageref*{e:233}}\par
\textbf{hookesche Gesetz} (das) -- Hooke's law\dotfill\hyperref[e:73]{\pageref*{e:73}}\par

\medskip{\large\textbf{\color{cornflower}I}}\par\nopagebreak
\textbf{ideale Gas} (das) -- ideal gas\dotfill\hyperref[e:140]{\pageref*{e:140}}\par
\textbf{Impuls} (der) -- momentum\dotfill\hyperref[e:79]{\pageref*{e:79}}\par
\textbf{Induktivität} (die) -- inductance\dotfill\hyperref[e:183]{\pageref*{e:183}}\par
\textbf{innere Energie} (die) -- internal energy\dotfill\hyperref[e:118]{\pageref*{e:118}}\par
\textbf{Interferenz} (die) -- interference\dotfill\hyperref[e:209]{\pageref*{e:209}}\par
\textbf{Ion} (das) -- ion\dotfill\hyperref[e:272]{\pageref*{e:272}}\par
\textbf{ionisierende Strahlung} (die) -- ionising radiation\dotfill\hyperref[e:278]{\pageref*{e:278}}\par
\textbf{Isolator, der Nichtleiter} (der) -- insulator\dotfill\hyperref[e:151]{\pageref*{e:151}}\par
\textbf{Isotop} (das) -- isotope\dotfill\hyperref[e:271]{\pageref*{e:271}}\par

\medskip{\large\textbf{\color{cornflower}K}}\par\nopagebreak
\textbf{Kapazität} (die) -- capacitance\dotfill\hyperref[e:170]{\pageref*{e:170}}\par
\textbf{Kernfusion} (die) -- nuclear fusion\dotfill\hyperref[e:286]{\pageref*{e:286}}\par
\textbf{Kernladungszahl, die Ordnungszahl} (die) -- atomic number, proton number\dotfill\hyperref[e:267]{\pageref*{e:267}}\par
\textbf{Kernreaktor} (der) -- nuclear reactor\dotfill\hyperref[e:288]{\pageref*{e:288}}\par
\textbf{Kernschatten, der Halbschatten} (der) -- umbra, penumbra\dotfill\hyperref[e:226]{\pageref*{e:226}}\par
\textbf{Kernspaltung} (die) -- nuclear fission\dotfill\hyperref[e:285]{\pageref*{e:285}}\par
\textbf{Kettenreaktion} (die) -- chain reaction\dotfill\hyperref[e:287]{\pageref*{e:287}}\par
\textbf{Kilowattstunde} (die) -- kilowatt-hour\dotfill\hyperref[e:112]{\pageref*{e:112}}\par
\textbf{kinetische Energie} (die) -- kinetic energy\dotfill\hyperref[e:98]{\pageref*{e:98}}\par
\textbf{Kondensator} (der) -- capacitor\dotfill\hyperref[e:169]{\pageref*{e:169}}\par
\textbf{kondensieren} -- condense\dotfill\hyperref[e:133]{\pageref*{e:133}}\par
\textbf{Konstante} (die) -- constant\dotfill\hyperref[e:27]{\pageref*{e:27}}\par
\textbf{Körper} (der) -- body, object\dotfill\hyperref[e:28]{\pageref*{e:28}}\par
\textbf{Kraft} (die) -- force\dotfill\hyperref[e:59]{\pageref*{e:59}}\par
\textbf{Kräftegleichgewicht} (das) -- equilibrium\dotfill\hyperref[e:75]{\pageref*{e:75}}\par
\textbf{Kraftstoß} (der) -- impulse\dotfill\hyperref[e:80]{\pageref*{e:80}}\par
\textbf{Kraftwerk} (das) -- power station\dotfill\hyperref[e:111]{\pageref*{e:111}}\par
\textbf{Kreisbewegung} (die) -- circular motion\dotfill\hyperref[e:50]{\pageref*{e:50}}\par
\textbf{Kurzschluss} (der) -- short circuit\dotfill\hyperref[e:164]{\pageref*{e:164}}\par

\medskip{\large\textbf{\color{cornflower}L}}\par\nopagebreak
\textbf{Längenausdehnung} (die) -- linear expansion\dotfill\hyperref[e:139]{\pageref*{e:139}}\par
\textbf{Längswelle (Longitudinalwelle)} (die) -- longitudinal wave\dotfill\hyperref[e:206]{\pageref*{e:206}}\par
\textbf{Lautstärke} (die) -- loudness, volume\dotfill\hyperref[e:215]{\pageref*{e:215}}\par
\textbf{Leistung} (die) -- power\dotfill\hyperref[e:105]{\pageref*{e:105}}\par
\textbf{Leiter} (der) -- conductor\dotfill\hyperref[e:150]{\pageref*{e:150}}\par
\textbf{Leitwert} (der) -- conductance\dotfill\hyperref[e:149]{\pageref*{e:149}}\par
\textbf{Licht} (das) -- light\dotfill\hyperref[e:220]{\pageref*{e:220}}\par
\textbf{Lichtbündel} (das) -- beam of light\dotfill\hyperref[e:223]{\pageref*{e:223}}\par
\textbf{Lichtgeschwindigkeit} (die) -- speed of light\dotfill\hyperref[e:224]{\pageref*{e:224}}\par
\textbf{Lichtquelle} (die) -- light source\dotfill\hyperref[e:221]{\pageref*{e:221}}\par
\textbf{Lichtstärke} (die) -- luminous intensity\dotfill\hyperref[e:257]{\pageref*{e:257}}\par
\textbf{Lichtstrahl} (der) -- ray of light, light ray\dotfill\hyperref[e:222]{\pageref*{e:222}}\par
\textbf{Linse} (die) -- lens\dotfill\hyperref[e:240]{\pageref*{e:240}}\par
\textbf{Linsengleichung} (die) -- lens equation\dotfill\hyperref[e:248]{\pageref*{e:248}}\par
\textbf{Lorentzkraft} (die) -- Lorentz force\dotfill\hyperref[e:185]{\pageref*{e:185}}\par
\textbf{Luftdruck} (der) -- atmospheric pressure\dotfill\hyperref[e:91]{\pageref*{e:91}}\par
\textbf{Luftwiderstand} (der) -- air resistance, drag\dotfill\hyperref[e:69]{\pageref*{e:69}}\par
\textbf{Lupe} (die) -- magnifying glass\dotfill\hyperref[e:253]{\pageref*{e:253}}\par

\medskip{\large\textbf{\color{cornflower}M}}\par\nopagebreak
\textbf{Magnet} (der) -- magnet\dotfill\hyperref[e:174]{\pageref*{e:174}}\par
\textbf{Magnetfeld} (das) -- magnetic field\dotfill\hyperref[e:178]{\pageref*{e:178}}\par
\textbf{magnetische Fluss} (der) -- magnetic flux\dotfill\hyperref[e:180]{\pageref*{e:180}}\par
\textbf{magnetische Flussdichte} (die) -- magnetic flux density\dotfill\hyperref[e:179]{\pageref*{e:179}}\par
\textbf{Masse} (die) -- mass\dotfill\hyperref[e:56]{\pageref*{e:56}}\par
\textbf{Massendefekt} (der) -- mass defect\dotfill\hyperref[e:289]{\pageref*{e:289}}\par
\textbf{Massenzahl} (die) -- mass number, nucleon number\dotfill\hyperref[e:268]{\pageref*{e:268}}\par
\textbf{messen} -- measure\dotfill\hyperref[e:8]{\pageref*{e:8}}\par
\textbf{Messgenauigkeit} (die) -- accuracy\dotfill\hyperref[e:13]{\pageref*{e:13}}\par
\textbf{Messgerät} (das) -- measuring instrument\dotfill\hyperref[e:11]{\pageref*{e:11}}\par
\textbf{Messung} (die) -- measurement\dotfill\hyperref[e:9]{\pageref*{e:9}}\par
\textbf{Messunsicherheit} (die) -- uncertainty (of measurement)\dotfill\hyperref[e:14]{\pageref*{e:14}}\par
\textbf{Messwert} (der) -- reading, measured value\dotfill\hyperref[e:10]{\pageref*{e:10}}\par
\textbf{Mittelwert} (der) -- mean (value)\dotfill\hyperref[e:15]{\pageref*{e:15}}\par
\textbf{Modell} (das) -- model\dotfill\hyperref[e:22]{\pageref*{e:22}}\par
\textbf{Momentangeschwindigkeit} (die) -- instantaneous velocity\dotfill\hyperref[e:39]{\pageref*{e:39}}\par

\medskip{\large\textbf{\color{cornflower}N}}\par\nopagebreak
\textbf{Neutron} (das) -- neutron\dotfill\hyperref[e:264]{\pageref*{e:264}}\par
\textbf{Neutronenzahl} (die) -- neutron number\dotfill\hyperref[e:269]{\pageref*{e:269}}\par
\textbf{newtonsche Grundgesetz} (das) -- Newton's second law\dotfill\hyperref[e:77]{\pageref*{e:77}}\par
\textbf{Nordpol, der Südpol} (der) -- north pole, south pole\dotfill\hyperref[e:177]{\pageref*{e:177}}\par
\textbf{Normalkraft} (die) -- normal force\dotfill\hyperref[e:64]{\pageref*{e:64}}\par
\textbf{Nukleon} (das) -- nucleon\dotfill\hyperref[e:265]{\pageref*{e:265}}\par
\textbf{Nuklid} (das) -- nuclide\dotfill\hyperref[e:270]{\pageref*{e:270}}\par
\textbf{Nutzenergie} (die) -- useful energy\dotfill\hyperref[e:107]{\pageref*{e:107}}\par

\medskip{\large\textbf{\color{cornflower}O}}\par\nopagebreak
\textbf{ohmsche Gesetz} (das) -- Ohm's law\dotfill\hyperref[e:147]{\pageref*{e:147}}\par
\textbf{optische Achse} (die) -- principal axis\dotfill\hyperref[e:243]{\pageref*{e:243}}\par
\textbf{Ort} (der) -- position\dotfill\hyperref[e:34]{\pageref*{e:34}}\par

\medskip{\large\textbf{\color{cornflower}P}}\par\nopagebreak
\textbf{Parallelschaltung} (die) -- parallel connection\dotfill\hyperref[e:163]{\pageref*{e:163}}\par
\textbf{Periodendauer} (die) -- period\dotfill\hyperref[e:196]{\pageref*{e:196}}\par
\textbf{Photon} (das) -- photon\dotfill\hyperref[e:291]{\pageref*{e:291}}\par
\textbf{physikalische Größe} (die) -- physical quantity\dotfill\hyperref[e:1]{\pageref*{e:1}}\par
\textbf{plancksche Wirkungsquantum} (das) -- Planck constant\dotfill\hyperref[e:292]{\pageref*{e:292}}\par
\textbf{potentielle Energie} (die) -- (gravitational) potential energy\dotfill\hyperref[e:97]{\pageref*{e:97}}\par
\textbf{Prisma} (das) -- prism\dotfill\hyperref[e:254]{\pageref*{e:254}}\par
\textbf{proportional} -- directly proportional\dotfill\hyperref[e:25]{\pageref*{e:25}}\par
\textbf{Proton} (das) -- proton\dotfill\hyperref[e:263]{\pageref*{e:263}}\par

\medskip{\large\textbf{\color{cornflower}Q}}\par\nopagebreak
\textbf{Querwelle (Transversalwelle)} (die) -- transverse wave\dotfill\hyperref[e:205]{\pageref*{e:205}}\par

\medskip{\large\textbf{\color{cornflower}R}}\par\nopagebreak
\textbf{Radialkraft} (die) -- centripetal force\dotfill\hyperref[e:94]{\pageref*{e:94}}\par
\textbf{radioaktive Zerfall} (der) -- radioactive decay\dotfill\hyperref[e:274]{\pageref*{e:274}}\par
\textbf{Radioaktivität} (die) -- radioactivity\dotfill\hyperref[e:273]{\pageref*{e:273}}\par
\textbf{reelle Bild} (das) -- real image\dotfill\hyperref[e:250]{\pageref*{e:250}}\par
\textbf{Reflexion} (die) -- reflection\dotfill\hyperref[e:207]{\pageref*{e:207}}\par
\textbf{Reflexion} (die) -- reflection\dotfill\hyperref[e:227]{\pageref*{e:227}}\par
\textbf{Reflexionsgesetz} (das) -- law of reflection\dotfill\hyperref[e:228]{\pageref*{e:228}}\par
\textbf{Reflexionswinkel} (der) -- angle of reflection\dotfill\hyperref[e:231]{\pageref*{e:231}}\par
\textbf{Reibungskraft} (die) -- friction, frictional force\dotfill\hyperref[e:65]{\pageref*{e:65}}\par
\textbf{Reibungszahl} (die) -- coefficient of friction\dotfill\hyperref[e:66]{\pageref*{e:66}}\par
\textbf{Reihenschaltung} (die) -- series connection\dotfill\hyperref[e:162]{\pageref*{e:162}}\par
\textbf{Resonanz} (die) -- resonance\dotfill\hyperref[e:200]{\pageref*{e:200}}\par
\textbf{resultierende Kraft} (die) -- resultant force, net force\dotfill\hyperref[e:74]{\pageref*{e:74}}\par
\textbf{Richtung} (die) -- direction\dotfill\hyperref[e:62]{\pageref*{e:62}}\par
\textbf{Rolle} (die) -- pulley\dotfill\hyperref[e:86]{\pageref*{e:86}}\par
\textbf{Ruhe} (die) -- rest\dotfill\hyperref[e:31]{\pageref*{e:31}}\par
\textbf{Ruhelage} (die) -- equilibrium position\dotfill\hyperref[e:193]{\pageref*{e:193}}\par

\medskip{\large\textbf{\color{cornflower}S}}\par\nopagebreak
\textbf{Sammellinse} (die) -- converging lens, convex lens\dotfill\hyperref[e:241]{\pageref*{e:241}}\par
\textbf{Schall} (der) -- sound\dotfill\hyperref[e:212]{\pageref*{e:212}}\par
\textbf{Schallgeschwindigkeit} (die) -- speed of sound\dotfill\hyperref[e:213]{\pageref*{e:213}}\par
\textbf{Schallpegel} (der) -- sound level\dotfill\hyperref[e:216]{\pageref*{e:216}}\par
\textbf{Schalter} (der) -- switch\dotfill\hyperref[e:157]{\pageref*{e:157}}\par
\textbf{Schaltplan} (der) -- circuit diagram\dotfill\hyperref[e:154]{\pageref*{e:154}}\par
\textbf{Schatten} (der) -- shadow\dotfill\hyperref[e:225]{\pageref*{e:225}}\par
\textbf{schmelzen} -- melt\dotfill\hyperref[e:129]{\pageref*{e:129}}\par
\textbf{Schmelztemperatur} (die) -- melting point\dotfill\hyperref[e:135]{\pageref*{e:135}}\par
\textbf{Schmelzwärme} (die) -- latent heat of fusion\dotfill\hyperref[e:137]{\pageref*{e:137}}\par
\textbf{Schweredruck} (der) -- hydrostatic pressure\dotfill\hyperref[e:90]{\pageref*{e:90}}\par
\textbf{Schwerpunkt} (der) -- centre of gravity\dotfill\hyperref[e:88]{\pageref*{e:88}}\par
\textbf{Schwingung} (die) -- oscillation, vibration\dotfill\hyperref[e:190]{\pageref*{e:190}}\par
\textbf{Sicherung} (die) -- fuse\dotfill\hyperref[e:165]{\pageref*{e:165}}\par
\textbf{sieden} -- boil\dotfill\hyperref[e:132]{\pageref*{e:132}}\par
\textbf{Siedetemperatur} (die) -- boiling point\dotfill\hyperref[e:136]{\pageref*{e:136}}\par
\textbf{Skala} (die) -- scale\dotfill\hyperref[e:12]{\pageref*{e:12}}\par
\textbf{Spannenergie} (die) -- elastic potential energy\dotfill\hyperref[e:99]{\pageref*{e:99}}\par
\textbf{Spannungsmesser} (der) -- voltmeter\dotfill\hyperref[e:161]{\pageref*{e:161}}\par
\textbf{Spannungsquelle} (die) -- voltage source, power supply\dotfill\hyperref[e:155]{\pageref*{e:155}}\par
\textbf{Spektrum} (das) -- spectrum\dotfill\hyperref[e:255]{\pageref*{e:255}}\par
\textbf{spezifische Wärmekapazität} (die) -- specific heat capacity\dotfill\hyperref[e:119]{\pageref*{e:119}}\par
\textbf{spezifische Widerstand} (der) -- resistivity\dotfill\hyperref[e:148]{\pageref*{e:148}}\par
\textbf{Spiegel} (der) -- mirror\dotfill\hyperref[e:232]{\pageref*{e:232}}\par
\textbf{Spule} (die) -- coil\dotfill\hyperref[e:181]{\pageref*{e:181}}\par
\textbf{stehende Welle} (die) -- standing wave, stationary wave\dotfill\hyperref[e:211]{\pageref*{e:211}}\par
\textbf{Stoff} (der) -- substance, material\dotfill\hyperref[e:29]{\pageref*{e:29}}\par
\textbf{Stromkreis} (der) -- electric circuit\dotfill\hyperref[e:153]{\pageref*{e:153}}\par
\textbf{Strommesser} (der) -- ammeter\dotfill\hyperref[e:160]{\pageref*{e:160}}\par
\textbf{sublimieren} -- sublime\dotfill\hyperref[e:134]{\pageref*{e:134}}\par

\medskip{\large\textbf{\color{cornflower}T}}\par\nopagebreak
\textbf{Teilchenmodell} (das) -- particle model\dotfill\hyperref[e:124]{\pageref*{e:124}}\par
\textbf{Temperatur} (die) -- temperature\dotfill\hyperref[e:113]{\pageref*{e:113}}\par
\textbf{thermische Energie} (die) -- thermal energy\dotfill\hyperref[e:100]{\pageref*{e:100}}\par
\textbf{Thermometer} (das) -- thermometer\dotfill\hyperref[e:116]{\pageref*{e:116}}\par
\textbf{Tonhöhe} (die) -- pitch\dotfill\hyperref[e:214]{\pageref*{e:214}}\par
\textbf{Totalreflexion} (die) -- total internal reflection\dotfill\hyperref[e:238]{\pageref*{e:238}}\par
\textbf{Trägheit} (die) -- inertia\dotfill\hyperref[e:76]{\pageref*{e:76}}\par
\textbf{Transformator} (der) -- transformer\dotfill\hyperref[e:188]{\pageref*{e:188}}\par

\medskip{\large\textbf{\color{cornflower}U}}\par\nopagebreak
\textbf{Überlagerung} (die) -- superposition\dotfill\hyperref[e:210]{\pageref*{e:210}}\par
\textbf{Ultraschall} (der) -- ultrasound\dotfill\hyperref[e:218]{\pageref*{e:218}}\par
\textbf{umgekehrt proportional} -- inversely proportional\dotfill\hyperref[e:26]{\pageref*{e:26}}\par
\textbf{Umlaufzeit} (die) -- period\dotfill\hyperref[e:51]{\pageref*{e:51}}\par

\medskip{\large\textbf{\color{cornflower}V}}\par\nopagebreak
\textbf{verdampfen} -- evaporate, vaporise\dotfill\hyperref[e:131]{\pageref*{e:131}}\par
\textbf{Verdampfungswärme} (die) -- latent heat of vaporisation\dotfill\hyperref[e:138]{\pageref*{e:138}}\par
\textbf{Vermutung, die Hypothese} (die) -- hypothesis\dotfill\hyperref[e:20]{\pageref*{e:20}}\par
\textbf{Versuch, das Experiment} (der) -- experiment\dotfill\hyperref[e:16]{\pageref*{e:16}}\par
\textbf{Versuchsaufbau} (der) -- experimental set-up\dotfill\hyperref[e:17]{\pageref*{e:17}}\par
\textbf{Verzögerung} (die) -- deceleration\dotfill\hyperref[e:41]{\pageref*{e:41}}\par
\textbf{virtuelle Bild} (das) -- virtual image\dotfill\hyperref[e:251]{\pageref*{e:251}}\par
\textbf{Volumen} (das) -- volume\dotfill\hyperref[e:57]{\pageref*{e:57}}\par
\textbf{Vorsatz} (der) -- prefix\dotfill\hyperref[e:7]{\pageref*{e:7}}\par

\medskip{\large\textbf{\color{cornflower}W}}\par\nopagebreak
\textbf{waagerechte Wurf} (der) -- horizontal projection\dotfill\hyperref[e:48]{\pageref*{e:48}}\par
\textbf{Wärme} (die) -- heat\dotfill\hyperref[e:117]{\pageref*{e:117}}\par
\textbf{Wärmedämmung} (die) -- thermal insulation\dotfill\hyperref[e:123]{\pageref*{e:123}}\par
\textbf{Wärmekraftmaschine} (die) -- heat engine\dotfill\hyperref[e:141]{\pageref*{e:141}}\par
\textbf{Wärmeleitung} (die) -- conduction\dotfill\hyperref[e:120]{\pageref*{e:120}}\par
\textbf{Wärmestrahlung} (die) -- thermal radiation\dotfill\hyperref[e:122]{\pageref*{e:122}}\par
\textbf{Wärmeströmung, die Konvektion} (die) -- convection\dotfill\hyperref[e:121]{\pageref*{e:121}}\par
\textbf{Wechselstrom} (der) -- alternating current (AC)\dotfill\hyperref[e:167]{\pageref*{e:167}}\par
\textbf{Wechselwirkungsgesetz} (das) -- Newton's third law (action and reaction)\dotfill\hyperref[e:78]{\pageref*{e:78}}\par
\textbf{Weg} (der) -- distance\dotfill\hyperref[e:35]{\pageref*{e:35}}\par
\textbf{Weg-Zeit-Diagramm} (das) -- distance–time graph\dotfill\hyperref[e:44]{\pageref*{e:44}}\par
\textbf{Welle} (die) -- wave\dotfill\hyperref[e:201]{\pageref*{e:201}}\par
\textbf{Wellenberg, das Wellental} (der) -- crest, trough\dotfill\hyperref[e:204]{\pageref*{e:204}}\par
\textbf{Wellenlänge} (die) -- wavelength\dotfill\hyperref[e:202]{\pageref*{e:202}}\par
\textbf{Widerstand (Bauteil)} (der) -- resistor\dotfill\hyperref[e:159]{\pageref*{e:159}}\par
\textbf{Windungszahl} (die) -- number of turns\dotfill\hyperref[e:182]{\pageref*{e:182}}\par
\textbf{Winkelgeschwindigkeit} (die) -- angular velocity\dotfill\hyperref[e:54]{\pageref*{e:54}}\par
\textbf{Wirkungsgrad} (der) -- efficiency\dotfill\hyperref[e:106]{\pageref*{e:106}}\par
\textbf{Wurfweite} (die) -- range\dotfill\hyperref[e:49]{\pageref*{e:49}}\par

\medskip{\large\textbf{\color{cornflower}Z}}\par\nopagebreak
\textbf{Zahlenwert} (der) -- numerical value\dotfill\hyperref[e:2]{\pageref*{e:2}}\par
\textbf{Zählrohr} (das) -- Geiger counter\dotfill\hyperref[e:279]{\pageref*{e:279}}\par
\textbf{Zeit} (die) -- time\dotfill\hyperref[e:36]{\pageref*{e:36}}\par
\textbf{Zerfallsgesetz} (das) -- decay law\dotfill\hyperref[e:281]{\pageref*{e:281}}\par
\textbf{Zerfallskonstante} (die) -- decay constant\dotfill\hyperref[e:282]{\pageref*{e:282}}\par
\textbf{Zerstreuungslinse} (die) -- diverging lens, concave lens\dotfill\hyperref[e:242]{\pageref*{e:242}}\par
\end{multicols}

\newpage
\section{Index English -- German}
\begin{multicols}{2}\small\raggedright

\medskip{\large\textbf{\color{cornflower}(}}\par\nopagebreak
\textbf{(gravitational) potential energy} -- die potentielle Energie\dotfill\hyperref[e:97]{\pageref*{e:97}}\par

\medskip{\large\textbf{\color{cornflower}A}}\par\nopagebreak
\textbf{absolute temperature, thermodynamic temperature} -- die absolute Temperatur\dotfill\hyperref[e:114]{\pageref*{e:114}}\par
\textbf{absolute zero} -- der absolute Nullpunkt\dotfill\hyperref[e:115]{\pageref*{e:115}}\par
\textbf{acceleration} -- die Beschleunigung\dotfill\hyperref[e:40]{\pageref*{e:40}}\par
\textbf{acceleration due to gravity} -- die Fallbeschleunigung\dotfill\hyperref[e:47]{\pageref*{e:47}}\par
\textbf{accuracy} -- die Messgenauigkeit\dotfill\hyperref[e:13]{\pageref*{e:13}}\par
\textbf{activity} -- die Aktivität\dotfill\hyperref[e:283]{\pageref*{e:283}}\par
\textbf{air resistance, drag} -- der Luftwiderstand\dotfill\hyperref[e:69]{\pageref*{e:69}}\par
\textbf{alpha radiation} -- die Alphastrahlung\dotfill\hyperref[e:275]{\pageref*{e:275}}\par
\textbf{alternating current (AC)} -- der Wechselstrom\dotfill\hyperref[e:167]{\pageref*{e:167}}\par
\textbf{ammeter} -- der Strommesser\dotfill\hyperref[e:160]{\pageref*{e:160}}\par
\textbf{amplitude} -- die Amplitude\dotfill\hyperref[e:195]{\pageref*{e:195}}\par
\textbf{analysis, evaluation} -- die Auswertung\dotfill\hyperref[e:19]{\pageref*{e:19}}\par
\textbf{angle of incidence} -- der Einfallswinkel\dotfill\hyperref[e:230]{\pageref*{e:230}}\par
\textbf{angle of reflection} -- der Reflexionswinkel\dotfill\hyperref[e:231]{\pageref*{e:231}}\par
\textbf{angle of refraction} -- der Brechungswinkel\dotfill\hyperref[e:235]{\pageref*{e:235}}\par
\textbf{angular velocity} -- die Winkelgeschwindigkeit\dotfill\hyperref[e:54]{\pageref*{e:54}}\par
\textbf{atmospheric pressure} -- der Luftdruck\dotfill\hyperref[e:91]{\pageref*{e:91}}\par
\textbf{atom} -- das Atom\dotfill\hyperref[e:259]{\pageref*{e:259}}\par
\textbf{atomic number, proton number} -- die Kernladungszahl, die Ordnungszahl\dotfill\hyperref[e:267]{\pageref*{e:267}}\par
\textbf{average speed} -- die Durchschnittsgeschwindigkeit\dotfill\hyperref[e:38]{\pageref*{e:38}}\par

\medskip{\large\textbf{\color{cornflower}B}}\par\nopagebreak
\textbf{base unit} -- die Basiseinheit\dotfill\hyperref[e:4]{\pageref*{e:4}}\par
\textbf{battery, cell} -- die Batterie\dotfill\hyperref[e:156]{\pageref*{e:156}}\par
\textbf{beam of light} -- das Lichtbündel\dotfill\hyperref[e:223]{\pageref*{e:223}}\par
\textbf{beta radiation} -- die Betastrahlung\dotfill\hyperref[e:276]{\pageref*{e:276}}\par
\textbf{binding energy} -- die Bindungsenergie\dotfill\hyperref[e:290]{\pageref*{e:290}}\par
\textbf{block and tackle} -- der Flaschenzug\dotfill\hyperref[e:87]{\pageref*{e:87}}\par
\textbf{body, object} -- der Körper\dotfill\hyperref[e:28]{\pageref*{e:28}}\par
\textbf{boil} -- sieden\dotfill\hyperref[e:132]{\pageref*{e:132}}\par
\textbf{boiling point} -- die Siedetemperatur\dotfill\hyperref[e:136]{\pageref*{e:136}}\par
\textbf{buoyancy, upthrust} -- der Auftrieb\dotfill\hyperref[e:92]{\pageref*{e:92}}\par

\medskip{\large\textbf{\color{cornflower}C}}\par\nopagebreak
\textbf{capacitance} -- die Kapazität\dotfill\hyperref[e:170]{\pageref*{e:170}}\par
\textbf{capacitor} -- der Kondensator\dotfill\hyperref[e:169]{\pageref*{e:169}}\par
\textbf{centre of gravity} -- der Schwerpunkt\dotfill\hyperref[e:88]{\pageref*{e:88}}\par
\textbf{centripetal force} -- die Radialkraft\dotfill\hyperref[e:94]{\pageref*{e:94}}\par
\textbf{chain reaction} -- die Kettenreaktion\dotfill\hyperref[e:287]{\pageref*{e:287}}\par
\textbf{chemical energy} -- die chemische Energie\dotfill\hyperref[e:102]{\pageref*{e:102}}\par
\textbf{circuit diagram} -- der Schaltplan\dotfill\hyperref[e:154]{\pageref*{e:154}}\par
\textbf{circular motion} -- die Kreisbewegung\dotfill\hyperref[e:50]{\pageref*{e:50}}\par
\textbf{coefficient of friction} -- die Reibungszahl\dotfill\hyperref[e:66]{\pageref*{e:66}}\par
\textbf{coil} -- die Spule\dotfill\hyperref[e:181]{\pageref*{e:181}}\par
\textbf{concave mirror, convex mirror} -- der Hohlspiegel, der Wölbspiegel\dotfill\hyperref[e:233]{\pageref*{e:233}}\par
\textbf{condense} -- kondensieren\dotfill\hyperref[e:133]{\pageref*{e:133}}\par
\textbf{conductance} -- der Leitwert\dotfill\hyperref[e:149]{\pageref*{e:149}}\par
\textbf{conduction} -- die Wärmeleitung\dotfill\hyperref[e:120]{\pageref*{e:120}}\par
\textbf{conductor} -- der Leiter\dotfill\hyperref[e:150]{\pageref*{e:150}}\par
\textbf{constant} -- die Konstante\dotfill\hyperref[e:27]{\pageref*{e:27}}\par
\textbf{convection} -- die Wärmeströmung, die Konvektion\dotfill\hyperref[e:121]{\pageref*{e:121}}\par
\textbf{converging lens, convex lens} -- die Sammellinse\dotfill\hyperref[e:241]{\pageref*{e:241}}\par
\textbf{crest, trough} -- der Wellenberg, das Wellental\dotfill\hyperref[e:204]{\pageref*{e:204}}\par
\textbf{critical angle} -- der Grenzwinkel\dotfill\hyperref[e:239]{\pageref*{e:239}}\par

\medskip{\large\textbf{\color{cornflower}D}}\par\nopagebreak
\textbf{damped oscillation} -- die gedämpfte Schwingung\dotfill\hyperref[e:198]{\pageref*{e:198}}\par
\textbf{decay constant} -- die Zerfallskonstante\dotfill\hyperref[e:282]{\pageref*{e:282}}\par
\textbf{decay law} -- das Zerfallsgesetz\dotfill\hyperref[e:281]{\pageref*{e:281}}\par
\textbf{deceleration} -- die Verzögerung\dotfill\hyperref[e:41]{\pageref*{e:41}}\par
\textbf{density} -- die Dichte\dotfill\hyperref[e:58]{\pageref*{e:58}}\par
\textbf{derived unit} -- die abgeleitete Einheit\dotfill\hyperref[e:5]{\pageref*{e:5}}\par
\textbf{diffraction} -- die Beugung\dotfill\hyperref[e:208]{\pageref*{e:208}}\par
\textbf{direct current (DC)} -- der Gleichstrom\dotfill\hyperref[e:166]{\pageref*{e:166}}\par
\textbf{direction} -- die Richtung\dotfill\hyperref[e:62]{\pageref*{e:62}}\par
\textbf{directly proportional} -- proportional\dotfill\hyperref[e:25]{\pageref*{e:25}}\par
\textbf{dispersion} -- die Farbzerlegung, die Dispersion\dotfill\hyperref[e:256]{\pageref*{e:256}}\par
\textbf{displacement} -- die Auslenkung\dotfill\hyperref[e:194]{\pageref*{e:194}}\par
\textbf{distance} -- der Weg\dotfill\hyperref[e:35]{\pageref*{e:35}}\par
\textbf{distance–time graph} -- das Weg-Zeit-Diagramm\dotfill\hyperref[e:44]{\pageref*{e:44}}\par
\textbf{diverging lens, concave lens} -- die Zerstreuungslinse\dotfill\hyperref[e:242]{\pageref*{e:242}}\par
\textbf{Doppler effect} -- der Doppler-Effekt\dotfill\hyperref[e:219]{\pageref*{e:219}}\par

\medskip{\large\textbf{\color{cornflower}E}}\par\nopagebreak
\textbf{earth (BE), ground (AE)} -- die Erdung\dotfill\hyperref[e:189]{\pageref*{e:189}}\par
\textbf{echo} -- das Echo\dotfill\hyperref[e:217]{\pageref*{e:217}}\par
\textbf{efficiency} -- der Wirkungsgrad\dotfill\hyperref[e:106]{\pageref*{e:106}}\par
\textbf{elastic potential energy} -- die Spannenergie\dotfill\hyperref[e:99]{\pageref*{e:99}}\par
\textbf{electric charge} -- die elektrische Ladung\dotfill\hyperref[e:143]{\pageref*{e:143}}\par
\textbf{electric circuit} -- der Stromkreis\dotfill\hyperref[e:153]{\pageref*{e:153}}\par
\textbf{electric current} -- die elektrische Stromstärke\dotfill\hyperref[e:144]{\pageref*{e:144}}\par
\textbf{electric field} -- das elektrische Feld\dotfill\hyperref[e:171]{\pageref*{e:171}}\par
\textbf{electric field strength} -- die elektrische Feldstärke\dotfill\hyperref[e:172]{\pageref*{e:172}}\par
\textbf{electric motor} -- der Elektromotor\dotfill\hyperref[e:186]{\pageref*{e:186}}\par
\textbf{electrical energy} -- die elektrische Energie\dotfill\hyperref[e:101]{\pageref*{e:101}}\par
\textbf{electrical power} -- die elektrische Leistung\dotfill\hyperref[e:168]{\pageref*{e:168}}\par
\textbf{electromagnet} -- der Elektromagnet\dotfill\hyperref[e:176]{\pageref*{e:176}}\par
\textbf{electromagnetic induction} -- die elektromagnetische Induktion\dotfill\hyperref[e:184]{\pageref*{e:184}}\par
\textbf{electron} -- das Elektron\dotfill\hyperref[e:262]{\pageref*{e:262}}\par
\textbf{electron shell} -- die Atomhülle\dotfill\hyperref[e:261]{\pageref*{e:261}}\par
\textbf{elementary charge} -- die Elementarladung\dotfill\hyperref[e:266]{\pageref*{e:266}}\par
\textbf{energy} -- die Energie\dotfill\hyperref[e:95]{\pageref*{e:95}}\par
\textbf{energy conversion, energy transfer} -- die Energieumwandlung\dotfill\hyperref[e:103]{\pageref*{e:103}}\par
\textbf{energy level} -- das Energieniveau\dotfill\hyperref[e:294]{\pageref*{e:294}}\par
\textbf{energy source} -- der Energieträger\dotfill\hyperref[e:109]{\pageref*{e:109}}\par
\textbf{entropy} -- die Entropie\dotfill\hyperref[e:142]{\pageref*{e:142}}\par
\textbf{equilibrium} -- das Kräftegleichgewicht\dotfill\hyperref[e:75]{\pageref*{e:75}}\par
\textbf{equilibrium position} -- die Ruhelage\dotfill\hyperref[e:193]{\pageref*{e:193}}\par
\textbf{equivalent dose} -- die Äquivalentdosis\dotfill\hyperref[e:284]{\pageref*{e:284}}\par
\textbf{evaporate, vaporise} -- verdampfen\dotfill\hyperref[e:131]{\pageref*{e:131}}\par
\textbf{experiment} -- der Versuch, das Experiment\dotfill\hyperref[e:16]{\pageref*{e:16}}\par
\textbf{experimental set-up} -- der Versuchsaufbau\dotfill\hyperref[e:17]{\pageref*{e:17}}\par
\textbf{extension} -- die Dehnung, die Verlängerung\dotfill\hyperref[e:71]{\pageref*{e:71}}\par

\medskip{\large\textbf{\color{cornflower}F}}\par\nopagebreak
\textbf{field line} -- die Feldlinie\dotfill\hyperref[e:173]{\pageref*{e:173}}\par
\textbf{focal length} -- die Brennweite\dotfill\hyperref[e:245]{\pageref*{e:245}}\par
\textbf{focal point, principal focus} -- der Brennpunkt\dotfill\hyperref[e:244]{\pageref*{e:244}}\par
\textbf{force} -- die Kraft\dotfill\hyperref[e:59]{\pageref*{e:59}}\par
\textbf{frame of reference} -- das Bezugssystem\dotfill\hyperref[e:32]{\pageref*{e:32}}\par
\textbf{free fall} -- der freie Fall\dotfill\hyperref[e:46]{\pageref*{e:46}}\par
\textbf{freeze, solidify} -- erstarren\dotfill\hyperref[e:130]{\pageref*{e:130}}\par
\textbf{frequency} -- die Frequenz\dotfill\hyperref[e:52]{\pageref*{e:52}}\par
\textbf{frequency} -- die Frequenz\dotfill\hyperref[e:197]{\pageref*{e:197}}\par
\textbf{friction, frictional force} -- die Reibungskraft\dotfill\hyperref[e:65]{\pageref*{e:65}}\par
\textbf{fuse} -- die Sicherung\dotfill\hyperref[e:165]{\pageref*{e:165}}\par

\medskip{\large\textbf{\color{cornflower}G}}\par\nopagebreak
\textbf{gamma radiation} -- die Gammastrahlung\dotfill\hyperref[e:277]{\pageref*{e:277}}\par
\textbf{gaseous} -- gasförmig\dotfill\hyperref[e:128]{\pageref*{e:128}}\par
\textbf{Geiger counter} -- das Zählrohr\dotfill\hyperref[e:279]{\pageref*{e:279}}\par
\textbf{generator} -- der Generator\dotfill\hyperref[e:187]{\pageref*{e:187}}\par
\textbf{graph} -- das Diagramm\dotfill\hyperref[e:23]{\pageref*{e:23}}\par
\textbf{gravitation, gravity} -- die Gravitation\dotfill\hyperref[e:93]{\pageref*{e:93}}\par

\medskip{\large\textbf{\color{cornflower}H}}\par\nopagebreak
\textbf{half-life} -- die Halbwertszeit\dotfill\hyperref[e:280]{\pageref*{e:280}}\par
\textbf{heat} -- die Wärme\dotfill\hyperref[e:117]{\pageref*{e:117}}\par
\textbf{heat engine} -- die Wärmekraftmaschine\dotfill\hyperref[e:141]{\pageref*{e:141}}\par
\textbf{Hooke's law} -- das hookesche Gesetz\dotfill\hyperref[e:73]{\pageref*{e:73}}\par
\textbf{horizontal projection} -- der waagerechte Wurf\dotfill\hyperref[e:48]{\pageref*{e:48}}\par
\textbf{hydrostatic pressure} -- der Schweredruck\dotfill\hyperref[e:90]{\pageref*{e:90}}\par
\textbf{hypothesis} -- die Vermutung, die Hypothese\dotfill\hyperref[e:20]{\pageref*{e:20}}\par

\medskip{\large\textbf{\color{cornflower}I}}\par\nopagebreak
\textbf{ideal gas} -- das ideale Gas\dotfill\hyperref[e:140]{\pageref*{e:140}}\par
\textbf{illuminance} -- die Beleuchtungsstärke\dotfill\hyperref[e:258]{\pageref*{e:258}}\par
\textbf{image} -- das Bild\dotfill\hyperref[e:249]{\pageref*{e:249}}\par
\textbf{image distance} -- die Bildweite\dotfill\hyperref[e:247]{\pageref*{e:247}}\par
\textbf{impulse} -- der Kraftstoß\dotfill\hyperref[e:80]{\pageref*{e:80}}\par
\textbf{inclined plane} -- die geneigte Ebene\dotfill\hyperref[e:85]{\pageref*{e:85}}\par
\textbf{inductance} -- die Induktivität\dotfill\hyperref[e:183]{\pageref*{e:183}}\par
\textbf{inertia} -- die Trägheit\dotfill\hyperref[e:76]{\pageref*{e:76}}\par
\textbf{instantaneous velocity} -- die Momentangeschwindigkeit\dotfill\hyperref[e:39]{\pageref*{e:39}}\par
\textbf{insulator} -- der Isolator, der Nichtleiter\dotfill\hyperref[e:151]{\pageref*{e:151}}\par
\textbf{interference} -- die Interferenz\dotfill\hyperref[e:209]{\pageref*{e:209}}\par
\textbf{internal energy} -- die innere Energie\dotfill\hyperref[e:118]{\pageref*{e:118}}\par
\textbf{inversely proportional} -- umgekehrt proportional\dotfill\hyperref[e:26]{\pageref*{e:26}}\par
\textbf{ion} -- das Ion\dotfill\hyperref[e:272]{\pageref*{e:272}}\par
\textbf{ionising radiation} -- die ionisierende Strahlung\dotfill\hyperref[e:278]{\pageref*{e:278}}\par
\textbf{isotope} -- das Isotop\dotfill\hyperref[e:271]{\pageref*{e:271}}\par

\medskip{\large\textbf{\color{cornflower}K}}\par\nopagebreak
\textbf{kilowatt-hour} -- die Kilowattstunde\dotfill\hyperref[e:112]{\pageref*{e:112}}\par
\textbf{kinetic energy} -- die kinetische Energie\dotfill\hyperref[e:98]{\pageref*{e:98}}\par

\medskip{\large\textbf{\color{cornflower}L}}\par\nopagebreak
\textbf{lamp, bulb} -- die Glühlampe\dotfill\hyperref[e:158]{\pageref*{e:158}}\par
\textbf{latent heat of fusion} -- die Schmelzwärme\dotfill\hyperref[e:137]{\pageref*{e:137}}\par
\textbf{latent heat of vaporisation} -- die Verdampfungswärme\dotfill\hyperref[e:138]{\pageref*{e:138}}\par
\textbf{law} -- das Gesetz\dotfill\hyperref[e:21]{\pageref*{e:21}}\par
\textbf{law of conservation of energy} -- der Energieerhaltungssatz\dotfill\hyperref[e:104]{\pageref*{e:104}}\par
\textbf{law of reflection} -- das Reflexionsgesetz\dotfill\hyperref[e:228]{\pageref*{e:228}}\par
\textbf{lens} -- die Linse\dotfill\hyperref[e:240]{\pageref*{e:240}}\par
\textbf{lens equation} -- die Linsengleichung\dotfill\hyperref[e:248]{\pageref*{e:248}}\par
\textbf{lever} -- der Hebel\dotfill\hyperref[e:82]{\pageref*{e:82}}\par
\textbf{lever arm} -- der Hebelarm\dotfill\hyperref[e:83]{\pageref*{e:83}}\par
\textbf{light} -- das Licht\dotfill\hyperref[e:220]{\pageref*{e:220}}\par
\textbf{light source} -- die Lichtquelle\dotfill\hyperref[e:221]{\pageref*{e:221}}\par
\textbf{line of best fit} -- die Ausgleichsgerade\dotfill\hyperref[e:24]{\pageref*{e:24}}\par
\textbf{linear expansion} -- die Längenausdehnung\dotfill\hyperref[e:139]{\pageref*{e:139}}\par
\textbf{liquid} -- flüssig\dotfill\hyperref[e:127]{\pageref*{e:127}}\par
\textbf{longitudinal wave} -- die Längswelle (Longitudinalwelle)\dotfill\hyperref[e:206]{\pageref*{e:206}}\par
\textbf{Lorentz force} -- die Lorentzkraft\dotfill\hyperref[e:185]{\pageref*{e:185}}\par
\textbf{loudness, volume} -- die Lautstärke\dotfill\hyperref[e:215]{\pageref*{e:215}}\par
\textbf{luminous intensity} -- die Lichtstärke\dotfill\hyperref[e:257]{\pageref*{e:257}}\par

\medskip{\large\textbf{\color{cornflower}M}}\par\nopagebreak
\textbf{magnet} -- der Magnet\dotfill\hyperref[e:174]{\pageref*{e:174}}\par
\textbf{magnetic field} -- das Magnetfeld\dotfill\hyperref[e:178]{\pageref*{e:178}}\par
\textbf{magnetic flux} -- der magnetische Fluss\dotfill\hyperref[e:180]{\pageref*{e:180}}\par
\textbf{magnetic flux density} -- die magnetische Flussdichte\dotfill\hyperref[e:179]{\pageref*{e:179}}\par
\textbf{magnification} -- der Abbildungsmaßstab\dotfill\hyperref[e:252]{\pageref*{e:252}}\par
\textbf{magnifying glass} -- die Lupe\dotfill\hyperref[e:253]{\pageref*{e:253}}\par
\textbf{magnitude} -- der Betrag (einer Kraft)\dotfill\hyperref[e:61]{\pageref*{e:61}}\par
\textbf{mass} -- die Masse\dotfill\hyperref[e:56]{\pageref*{e:56}}\par
\textbf{mass defect} -- der Massendefekt\dotfill\hyperref[e:289]{\pageref*{e:289}}\par
\textbf{mass number, nucleon number} -- die Massenzahl\dotfill\hyperref[e:268]{\pageref*{e:268}}\par
\textbf{mass on a spring} -- der Federschwinger\dotfill\hyperref[e:192]{\pageref*{e:192}}\par
\textbf{mean (value)} -- der Mittelwert\dotfill\hyperref[e:15]{\pageref*{e:15}}\par
\textbf{measure} -- messen\dotfill\hyperref[e:8]{\pageref*{e:8}}\par
\textbf{measurement} -- die Messung\dotfill\hyperref[e:9]{\pageref*{e:9}}\par
\textbf{measuring instrument} -- das Messgerät\dotfill\hyperref[e:11]{\pageref*{e:11}}\par
\textbf{melt} -- schmelzen\dotfill\hyperref[e:129]{\pageref*{e:129}}\par
\textbf{melting point} -- die Schmelztemperatur\dotfill\hyperref[e:135]{\pageref*{e:135}}\par
\textbf{mirror} -- der Spiegel\dotfill\hyperref[e:232]{\pageref*{e:232}}\par
\textbf{model} -- das Modell\dotfill\hyperref[e:22]{\pageref*{e:22}}\par
\textbf{momentum} -- der Impuls\dotfill\hyperref[e:79]{\pageref*{e:79}}\par
\textbf{motion} -- die Bewegung\dotfill\hyperref[e:30]{\pageref*{e:30}}\par

\medskip{\large\textbf{\color{cornflower}N}}\par\nopagebreak
\textbf{natural frequency} -- die Eigenfrequenz\dotfill\hyperref[e:199]{\pageref*{e:199}}\par
\textbf{neutron} -- das Neutron\dotfill\hyperref[e:264]{\pageref*{e:264}}\par
\textbf{neutron number} -- die Neutronenzahl\dotfill\hyperref[e:269]{\pageref*{e:269}}\par
\textbf{Newton's second law} -- das newtonsche Grundgesetz\dotfill\hyperref[e:77]{\pageref*{e:77}}\par
\textbf{Newton's third law (action and reaction)} -- das Wechselwirkungsgesetz\dotfill\hyperref[e:78]{\pageref*{e:78}}\par
\textbf{normal} -- das Einfallslot\dotfill\hyperref[e:229]{\pageref*{e:229}}\par
\textbf{normal force} -- die Normalkraft\dotfill\hyperref[e:64]{\pageref*{e:64}}\par
\textbf{north pole, south pole} -- der Nordpol, der Südpol\dotfill\hyperref[e:177]{\pageref*{e:177}}\par
\textbf{nuclear fission} -- die Kernspaltung\dotfill\hyperref[e:285]{\pageref*{e:285}}\par
\textbf{nuclear fusion} -- die Kernfusion\dotfill\hyperref[e:286]{\pageref*{e:286}}\par
\textbf{nuclear reactor} -- der Kernreaktor\dotfill\hyperref[e:288]{\pageref*{e:288}}\par
\textbf{nucleon} -- das Nukleon\dotfill\hyperref[e:265]{\pageref*{e:265}}\par
\textbf{nucleus} -- der Atomkern\dotfill\hyperref[e:260]{\pageref*{e:260}}\par
\textbf{nuclide} -- das Nuklid\dotfill\hyperref[e:270]{\pageref*{e:270}}\par
\textbf{number of turns} -- die Windungszahl\dotfill\hyperref[e:182]{\pageref*{e:182}}\par
\textbf{numerical value} -- der Zahlenwert\dotfill\hyperref[e:2]{\pageref*{e:2}}\par

\medskip{\large\textbf{\color{cornflower}O}}\par\nopagebreak
\textbf{object distance} -- die Gegenstandsweite\dotfill\hyperref[e:246]{\pageref*{e:246}}\par
\textbf{observation} -- die Beobachtung\dotfill\hyperref[e:18]{\pageref*{e:18}}\par
\textbf{Ohm's law} -- das ohmsche Gesetz\dotfill\hyperref[e:147]{\pageref*{e:147}}\par
\textbf{oscillation, vibration} -- die Schwingung\dotfill\hyperref[e:190]{\pageref*{e:190}}\par

\medskip{\large\textbf{\color{cornflower}P}}\par\nopagebreak
\textbf{parallel connection} -- die Parallelschaltung\dotfill\hyperref[e:163]{\pageref*{e:163}}\par
\textbf{particle model} -- das Teilchenmodell\dotfill\hyperref[e:124]{\pageref*{e:124}}\par
\textbf{path, trajectory} -- die Bahn\dotfill\hyperref[e:33]{\pageref*{e:33}}\par
\textbf{period} -- die Periodendauer\dotfill\hyperref[e:196]{\pageref*{e:196}}\par
\textbf{period} -- die Umlaufzeit\dotfill\hyperref[e:51]{\pageref*{e:51}}\par
\textbf{permanent magnet} -- der Dauermagnet\dotfill\hyperref[e:175]{\pageref*{e:175}}\par
\textbf{photoelectric effect} -- der Fotoeffekt\dotfill\hyperref[e:293]{\pageref*{e:293}}\par
\textbf{photon} -- das Photon\dotfill\hyperref[e:291]{\pageref*{e:291}}\par
\textbf{physical quantity} -- die physikalische Größe\dotfill\hyperref[e:1]{\pageref*{e:1}}\par
\textbf{pitch} -- die Tonhöhe\dotfill\hyperref[e:214]{\pageref*{e:214}}\par
\textbf{pivot, fulcrum} -- der Drehpunkt\dotfill\hyperref[e:84]{\pageref*{e:84}}\par
\textbf{Planck constant} -- das plancksche Wirkungsquantum\dotfill\hyperref[e:292]{\pageref*{e:292}}\par
\textbf{point of application} -- der Angriffspunkt\dotfill\hyperref[e:60]{\pageref*{e:60}}\par
\textbf{position} -- der Ort\dotfill\hyperref[e:34]{\pageref*{e:34}}\par
\textbf{power} -- die Leistung\dotfill\hyperref[e:105]{\pageref*{e:105}}\par
\textbf{power station} -- das Kraftwerk\dotfill\hyperref[e:111]{\pageref*{e:111}}\par
\textbf{prefix} -- der Vorsatz\dotfill\hyperref[e:7]{\pageref*{e:7}}\par
\textbf{pressure} -- der Druck\dotfill\hyperref[e:89]{\pageref*{e:89}}\par
\textbf{principal axis} -- die optische Achse\dotfill\hyperref[e:243]{\pageref*{e:243}}\par
\textbf{principal quantum number} -- die Hauptquantenzahl\dotfill\hyperref[e:295]{\pageref*{e:295}}\par
\textbf{prism} -- das Prisma\dotfill\hyperref[e:254]{\pageref*{e:254}}\par
\textbf{proton} -- das Proton\dotfill\hyperref[e:263]{\pageref*{e:263}}\par
\textbf{pulley} -- die Rolle\dotfill\hyperref[e:86]{\pageref*{e:86}}\par

\medskip{\large\textbf{\color{cornflower}R}}\par\nopagebreak
\textbf{radioactive decay} -- der radioaktive Zerfall\dotfill\hyperref[e:274]{\pageref*{e:274}}\par
\textbf{radioactivity} -- die Radioaktivität\dotfill\hyperref[e:273]{\pageref*{e:273}}\par
\textbf{range} -- die Wurfweite\dotfill\hyperref[e:49]{\pageref*{e:49}}\par
\textbf{ray of light, light ray} -- der Lichtstrahl\dotfill\hyperref[e:222]{\pageref*{e:222}}\par
\textbf{reading, measured value} -- der Messwert\dotfill\hyperref[e:10]{\pageref*{e:10}}\par
\textbf{real image} -- das reelle Bild\dotfill\hyperref[e:250]{\pageref*{e:250}}\par
\textbf{reflection} -- die Reflexion\dotfill\hyperref[e:207]{\pageref*{e:207}}\par
\textbf{reflection} -- die Reflexion\dotfill\hyperref[e:227]{\pageref*{e:227}}\par
\textbf{refraction} -- die Brechung\dotfill\hyperref[e:234]{\pageref*{e:234}}\par
\textbf{refractive index} -- die Brechzahl\dotfill\hyperref[e:236]{\pageref*{e:236}}\par
\textbf{renewable energy} -- die erneuerbare Energie\dotfill\hyperref[e:110]{\pageref*{e:110}}\par
\textbf{resistance} -- der elektrische Widerstand\dotfill\hyperref[e:146]{\pageref*{e:146}}\par
\textbf{resistivity} -- der spezifische Widerstand\dotfill\hyperref[e:148]{\pageref*{e:148}}\par
\textbf{resistor} -- der Widerstand (Bauteil)\dotfill\hyperref[e:159]{\pageref*{e:159}}\par
\textbf{resonance} -- die Resonanz\dotfill\hyperref[e:200]{\pageref*{e:200}}\par
\textbf{rest} -- die Ruhe\dotfill\hyperref[e:31]{\pageref*{e:31}}\par
\textbf{resultant force, net force} -- die resultierende Kraft\dotfill\hyperref[e:74]{\pageref*{e:74}}\par
\textbf{rotational speed} -- die Drehzahl\dotfill\hyperref[e:53]{\pageref*{e:53}}\par

\medskip{\large\textbf{\color{cornflower}S}}\par\nopagebreak
\textbf{scale} -- die Skala\dotfill\hyperref[e:12]{\pageref*{e:12}}\par
\textbf{semiconductor} -- der Halbleiter\dotfill\hyperref[e:152]{\pageref*{e:152}}\par
\textbf{series connection} -- die Reihenschaltung\dotfill\hyperref[e:162]{\pageref*{e:162}}\par
\textbf{shadow} -- der Schatten\dotfill\hyperref[e:225]{\pageref*{e:225}}\par
\textbf{short circuit} -- der Kurzschluss\dotfill\hyperref[e:164]{\pageref*{e:164}}\par
\textbf{simple pendulum} -- das Fadenpendel\dotfill\hyperref[e:191]{\pageref*{e:191}}\par
\textbf{sliding friction} -- die Gleitreibung\dotfill\hyperref[e:68]{\pageref*{e:68}}\par
\textbf{Snell's law, law of refraction} -- das Brechungsgesetz\dotfill\hyperref[e:237]{\pageref*{e:237}}\par
\textbf{solid} -- fest\dotfill\hyperref[e:126]{\pageref*{e:126}}\par
\textbf{sound} -- der Schall\dotfill\hyperref[e:212]{\pageref*{e:212}}\par
\textbf{sound level} -- der Schallpegel\dotfill\hyperref[e:216]{\pageref*{e:216}}\par
\textbf{specific heat capacity} -- die spezifische Wärmekapazität\dotfill\hyperref[e:119]{\pageref*{e:119}}\par
\textbf{spectrum} -- das Spektrum\dotfill\hyperref[e:255]{\pageref*{e:255}}\par
\textbf{speed of light} -- die Lichtgeschwindigkeit\dotfill\hyperref[e:224]{\pageref*{e:224}}\par
\textbf{speed of sound} -- die Schallgeschwindigkeit\dotfill\hyperref[e:213]{\pageref*{e:213}}\par
\textbf{spring} -- die Feder\dotfill\hyperref[e:70]{\pageref*{e:70}}\par
\textbf{spring constant} -- die Federkonstante\dotfill\hyperref[e:72]{\pageref*{e:72}}\par
\textbf{standing wave, stationary wave} -- die stehende Welle\dotfill\hyperref[e:211]{\pageref*{e:211}}\par
\textbf{state of matter} -- der Aggregatzustand\dotfill\hyperref[e:125]{\pageref*{e:125}}\par
\textbf{static friction} -- die Haftreibung\dotfill\hyperref[e:67]{\pageref*{e:67}}\par
\textbf{sublime} -- sublimieren\dotfill\hyperref[e:134]{\pageref*{e:134}}\par
\textbf{substance, material} -- der Stoff\dotfill\hyperref[e:29]{\pageref*{e:29}}\par
\textbf{superposition} -- die Überlagerung\dotfill\hyperref[e:210]{\pageref*{e:210}}\par
\textbf{switch} -- der Schalter\dotfill\hyperref[e:157]{\pageref*{e:157}}\par
\textbf{symbol} -- das Formelzeichen\dotfill\hyperref[e:6]{\pageref*{e:6}}\par

\medskip{\large\textbf{\color{cornflower}T}}\par\nopagebreak
\textbf{tangential speed} -- die Bahngeschwindigkeit\dotfill\hyperref[e:55]{\pageref*{e:55}}\par
\textbf{temperature} -- die Temperatur\dotfill\hyperref[e:113]{\pageref*{e:113}}\par
\textbf{thermal energy} -- die thermische Energie\dotfill\hyperref[e:100]{\pageref*{e:100}}\par
\textbf{thermal insulation} -- die Wärmedämmung\dotfill\hyperref[e:123]{\pageref*{e:123}}\par
\textbf{thermal radiation} -- die Wärmestrahlung\dotfill\hyperref[e:122]{\pageref*{e:122}}\par
\textbf{thermometer} -- das Thermometer\dotfill\hyperref[e:116]{\pageref*{e:116}}\par
\textbf{time} -- die Zeit\dotfill\hyperref[e:36]{\pageref*{e:36}}\par
\textbf{torque, moment (of a force)} -- das Drehmoment\dotfill\hyperref[e:81]{\pageref*{e:81}}\par
\textbf{total internal reflection} -- die Totalreflexion\dotfill\hyperref[e:238]{\pageref*{e:238}}\par
\textbf{transformer} -- der Transformator\dotfill\hyperref[e:188]{\pageref*{e:188}}\par
\textbf{transverse wave} -- die Querwelle (Transversalwelle)\dotfill\hyperref[e:205]{\pageref*{e:205}}\par

\medskip{\large\textbf{\color{cornflower}U}}\par\nopagebreak
\textbf{ultrasound} -- der Ultraschall\dotfill\hyperref[e:218]{\pageref*{e:218}}\par
\textbf{umbra, penumbra} -- der Kernschatten, der Halbschatten\dotfill\hyperref[e:226]{\pageref*{e:226}}\par
\textbf{uncertainty (of measurement)} -- die Messunsicherheit\dotfill\hyperref[e:14]{\pageref*{e:14}}\par
\textbf{uniform motion} -- die gleichförmige Bewegung\dotfill\hyperref[e:42]{\pageref*{e:42}}\par
\textbf{uniformly accelerated motion} -- die gleichmäßig beschleunigte Bewegung\dotfill\hyperref[e:43]{\pageref*{e:43}}\par
\textbf{unit} -- die Einheit\dotfill\hyperref[e:3]{\pageref*{e:3}}\par
\textbf{useful energy} -- die Nutzenergie\dotfill\hyperref[e:107]{\pageref*{e:107}}\par

\medskip{\large\textbf{\color{cornflower}V}}\par\nopagebreak
\textbf{velocity, speed} -- die Geschwindigkeit\dotfill\hyperref[e:37]{\pageref*{e:37}}\par
\textbf{velocity–time graph} -- das Geschwindigkeit-Zeit-Diagramm\dotfill\hyperref[e:45]{\pageref*{e:45}}\par
\textbf{virtual image} -- das virtuelle Bild\dotfill\hyperref[e:251]{\pageref*{e:251}}\par
\textbf{voltage source, power supply} -- die Spannungsquelle\dotfill\hyperref[e:155]{\pageref*{e:155}}\par
\textbf{voltage, potential difference} -- die elektrische Spannung\dotfill\hyperref[e:145]{\pageref*{e:145}}\par
\textbf{voltmeter} -- der Spannungsmesser\dotfill\hyperref[e:161]{\pageref*{e:161}}\par
\textbf{volume} -- das Volumen\dotfill\hyperref[e:57]{\pageref*{e:57}}\par

\medskip{\large\textbf{\color{cornflower}W}}\par\nopagebreak
\textbf{waste heat} -- die Abwärme\dotfill\hyperref[e:108]{\pageref*{e:108}}\par
\textbf{wave} -- die Welle\dotfill\hyperref[e:201]{\pageref*{e:201}}\par
\textbf{wave speed} -- die Ausbreitungsgeschwindigkeit\dotfill\hyperref[e:203]{\pageref*{e:203}}\par
\textbf{wavelength} -- die Wellenlänge\dotfill\hyperref[e:202]{\pageref*{e:202}}\par
\textbf{weight} -- die Gewichtskraft\dotfill\hyperref[e:63]{\pageref*{e:63}}\par
\textbf{work} -- die Arbeit\dotfill\hyperref[e:96]{\pageref*{e:96}}\par
\end{multicols}

\end{document}
